\documentclass[11pt]{article}
\usepackage[review]{acl}
\usepackage{times}
\usepackage{latexsym}
\usepackage[T1]{fontenc}
\usepackage[utf8]{inputenc}
\usepackage{microtype}
\usepackage{inconsolata}
\usepackage{graphicx}
\usepackage{booktabs}
\usepackage{amsmath,amssymb,amsthm}
\usepackage{array}
\usepackage{longtable}
\usepackage{rotating}
\usepackage{url}
\usepackage{xspace}
\usepackage{tikz}
\usetikzlibrary{calc,positioning,arrows.meta,shapes.misc}
\usepackage{enumitem}
\usepackage{alphalph}
\usepackage[capitalise,noabbrev,nameinlink]{cleveref}
\setlist[itemize]{leftmargin=*,nosep}
\setcounter{topnumber}{4}\setcounter{bottomnumber}{2}\setcounter{totalnumber}{6}\setcounter{dbltopnumber}{4}
\renewcommand{\topfraction}{0.95}\renewcommand{\bottomfraction}{0.5}\renewcommand{\textfraction}{0.05}
\renewcommand{\floatpagefraction}{0.85}\renewcommand{\dbltopfraction}{0.95}\renewcommand{\dblfloatpagefraction}{0.85}
\newtheorem{lemma}{Lemma}
\crefname{lemma}{Lemma}{Lemmas}
\crefname{appendix}{Appendix}{Appendices}
\crefname{section}{Section}{Sections}
\Crefname{section}{Section}{Sections}

\newcommand{\method}{NutriTURN\xspace}
\newcommand{\compact}{Compact-7\xspace}
\newcommand{\twovar}{the two-variable score\xspace}
\newcommand{\stable}{\textsc{stable}\xspace}
\newcommand{\forming}{\textsc{still-forming}\xspace}
\newcommand{\unstable}{\textsc{unstable}\xspace}
\newcommand{\unassess}{\textsc{unassessable}\xspace}
\newcommand{\ci}[2]{[#1,\,#2]}
\newcommand{\sgn}{\operatorname{sgn}}
\DeclareMathOperator{\logit}{logit}
\newcommand{\BB}{\mathrm{BetaBin}}

\title{High Scores, Fragile Targets: Auditing Change in\\LLM-Annotated Diet-Health Evidence}
\author{Anonymous submission}

\begin{document}
\maketitle

\begin{abstract}
When a benchmark derives both inputs and outcomes from automatic annotations, high discrimination
can reflect label composition and representation choices rather than annotation-stable predictive
value. We present \method, a benchmark and audit framework built from 112{,}453 PubMed
claim-record associations for 165 diet-health claims, annotated in full by two LLM pipelines, with
a calendar-respecting protocol that separates three phenomena. Composition-driven ranking: an
untrained margin rule reaches AUROC 0.881, and an exact pair decomposition places about 91\% of a
signed logistic score's estimated $+0.067$ gain in comparisons across pre-cutoff majority
classes. Representation-dependent failure: under a second annotation the score loses to the
margin rule, 82\% of the loss falls within classes, and replacing its two-branch agreement input
by a signed-only one raises its AUROC from 0.758 to 0.940 on common support without changing the
outcomes, while the same input establishes no development gain. Annotation-dependent targets:
the two annotations share 27 of 96 positive events, and the second agrees less with an
annotation-independent numeric anchor. No calendar-cohort gain is established, the score loses on
one external domain, and the excess of reversals over era-conditioned controls is confined to the
primary annotation and depends on the temporal partition.
\end{abstract}

\section{Introduction}
\label{sec:intro}
Systems that retrieve, extract and verify biomedical evidence are increasingly asked a temporal
question: given what has been published so far, is the evidence on a claim likely to change?
Review-update prediction \citep{Bashir2019,Bashir2021,AlharbiStevenson2026} and living reviews
\citep{Elliott2014,Nevill2025} motivate the question, and early findings that later attenuate or
reverse are documented \citep{Ioannidis2005JAMA,IoannidisTrikalinos2005,Prasad2013}. At the scale of a literature such
datasets are built from automatic annotation: a model labels every document, labels before a cutoff become the inputs, and a thresholded aggregate of labels after it becomes the outcome.
This construction admits high scores with no predictive value beyond what the labels already
carry: under an exchangeable label process the outcome's probability is a closed-form function of
the pre-cutoff counts (\cref{lem:stationary}), so composition alone ranks it well; a signed input can
learn which label class tends to change under one annotator, a class prior with no guarantee of
transfer; an input whose definition switches on a threshold can fail when a second annotator moves
units across it; a second annotator can define a different set of positive events; and retrospective
evaluation lets later outcomes shape the model applied at a cutoff. A benchmark that
does not separate them rewards them \citep{Gururangan2018,Poliak2018,McCoy2019}.

We present \method,\footnote{Data, code, and results:
\url{https://anonymous.4open.science/r/NutriTURN-0C33}.} a benchmark for forecasting
evidence-direction change built to test these explanations separately: a
longitudinal PubMed stream of 165 exposure-outcome claims, two complete LLM annotations of all
112{,}453 claim-record associations, an origin-defined calendar cohort, two frozen external domains and null controls, with a protocol that asks whether a gain over an analytic composition reference holds across annotations, in calendar time, within majority classes and under transfer (\cref{sec:protocol}).
Applying it to composition references, learned scores and prompted LLM forecasters yields four findings (\cref{sec:results}) that separate three phenomena: (1)~label composition
explains most of the ranking; (2)~the learned gain is concentrated across classes: an exact pair
decomposition places about 91\% of it in pairs with different pre-cutoff majority classes, a
class-rate difference (five to one) that persists within the reported margin strata; (3)~the shipped score's
failure under a second annotation is largely representation-dependent: 82\% of its loss falls
within classes, and replacing its two-branch agreement input by a signed-only one repairs the loss
($0.758$ to $0.940$ on common support) without changing the outcomes, while the same input
establishes no development gain; no gain is established in the calendar cohort and a loss is
established on one external domain; (4)~repairing the predictor does not repair the target:
the two annotations share 27 of 96 positive events, and the one that redefines the positive class
most agrees less with an annotation-independent numeric anchor. Our contributions are \textbf{(i)} the benchmark; \textbf{(ii)} a five-check
protocol with fixed primary cells and tiers; and \textbf{(iii)} diagnostics that locate where a
ranking comes from (composition references, pair decomposition, class-blind and branch-free
variants, sign exchanges, null controls, a numeric anchor for both annotations).

\section{Related Work}
\paragraph{Evidence change and review updating.} Conclusion-change prediction for systematic-review
updates \citep{Bashir2019,Bashir2021} was scaled by \citet{AlharbiStevenson2026} to 3{,}326 Cochrane
review-update pairs; it predicts change \emph{given} the new records, whereas our task uses
records published before the origin only; conclusiveness assessment \citep{Ban2026} and
living-evidence automation \citep{Song2026,Harasgama2026} address the same problem. Trialstreamer \citep{Marshall2020} and RobotReviewer
\citep{Marshall2015} build evidence sets, Evidence Inference \citep{Lehman2019,DeYoung2020} labels
effect direction per record, as our label does, and SciFact
\citep{Wadden2020SciFact}, PUBHEALTH \citep{Kotonya2020PubHealth} and HealthVer
\citep{Sarrouti2021HealthVer} decide current support. Our target is a vote-count sign change
\citep{HedgesOlkin1980}, not a cumulative meta-analytic estimate
\citep{Lau1992,Pogue1997,Wetterslev2008}.

\paragraph{Forecasting and temporal evaluation.} ForecastQA \citep{Jin2021ForecastQA}, Autocast
\citep{Zou2022Autocast} and retrieval-augmented LLM forecasters \citep{Halawi2024} forecast from
dated text; retrospective evaluation is vulnerable to temporal leakage \citep{Paleka2025}, and TemporalWiki,
StreamingQA and TempLAMA \citep{Jang2022TemporalWiki,Liska2022StreamingQA,Dhingra2022TempLAMA,Lazaridou2021}
evaluate models against knowledge dated after training. Our calendar cohort addresses the same risk
for a target in which the annotator, not the world, decides what counts as a change.

\paragraph{Shortcuts, label quality and LLM annotation.} Hypothesis-only baselines \citep{Gururangan2018,Poliak2018},
HANS \citep{McCoy2019}, contrast sets \citep{Gardner2020} and adversarial collection
\citep{Nie2020} expose decision rules that succeed on a benchmark and fail off it
\citep{Geirhos2020}, and models can learn the annotators who wrote the data \citep{Geva2019}. Those
shortcuts are input artefacts; here the target itself carries little beyond
composition, what it carries can be specific to one annotator, and what a score extracts from it
can be specific to one feature definition.
Annotator disagreement is often signal \citep{Pavlick2019,Plank2022}, distant labels trade
noise for scale \citep{Mintz2009}, label corrections can reorder rankings \citep{Northcutt2021}, LLM annotators require task-specific
validation \citep{Gilardi2023,Ziems2024,Pangakis2023}, and configuration choices change
downstream conclusions \citep{Baumann2025}. Measuring inputs and outcome with one instrument is a known
source of shared error \citep{Podsakoff2003}; we show its temporal consequence.

\section{The Benchmark}
\label{sec:task}

\begin{figure*}[!t]
\centering
\includegraphics[width=0.83\textwidth]{figures/nutrifig.pdf}
\caption{\textbf{\method: construction and evaluation.} A frozen PubMed stream of 165 claims is
annotated in full by two configurations that differ in model, prompt and read-out. Each
claim-cutoff unit $(c,t)$ takes its inputs from pre-cutoff labels, its future direction from
post-cutoff labels only, and its endpoint $R(c,t)$ from the comparison of the two signs (red
\textsc{harmful}, blue \textsc{protective}, grey \textsc{null}). Scores are compared with the
composition references by paired $\Delta$AUROC; Tier~A licenses a retrospective gain over
composition, Tier~B adds non-inferiority across annotations and domains (at a fixed margin, not
robustness), Tier~C adds the calendar and class-invariance checks (\cref{sec:protocol}). The
bottom band shows the four null controls and the A$\to$C agreement-input intervention.}
\label{fig:protocol}
\end{figure*}

\paragraph{Stream.} Each of 165 claims $c=(e,o)$ pairs an exposure with an outcome and has one fixed
PubMed query over MeSH and title/abstract terms, without date or publication-type filters. The
archived snapshot holds 112{,}453 claim-record associations over 96{,}091 distinct identifiers; the claim list is a curated artefact of 145 candidate claims and 20 documented additions
(\cref{app:bench}), which contribute 6 of the 71 events in two claims and so do not seed the result. The harvest is bounded to 9-10 September 2026, and the archived identifier sets,
not the queries, define the benchmark: re-running the stored queries returns 91.7\% of the archived associations; for 19 claims, flagged in the release, the stored query now returns nothing and is not confirmed to be the harvesting query; and 48 claims were truncated by a per-query record cap that removed their early history. Removing either set changes no
protocol conclusion (\cref{app:snapshot}). Coverage has a concrete cost: of eight claims with widely described
reversals, three observable ones have no eligible cutoff before the pivotal trial and no trial
report in their archive, so no forecast could have been posed for them (\cref{app:robust}). Source records are released by PMID, the external corpora as unit-level label counts only.

\paragraph{Annotations.} Every claim-record association carries a direction label
(\textsc{protective}, \textsc{null}, \textsc{harmful} or \textsc{unclear}) relative to the named
exposure and outcome, not to clinical desirability (benefit-valence outcomes invert the label
names; removing such claims leaves the result unchanged, \cref{app:orientation}). The \emph{primary} pipeline reads title and
abstract with Llama-3.1-8B-Instruct \citep{llama3} and compares the first-token logits of the four
labels. The \emph{second} pipeline re-annotates every association with
Qwen2.5-7B-Instruct \citep{qwen25} under the same task framing and text budget, with explicit
\textsc{null} and \textsc{unclear} definitions and full label-string probabilities, so the
configurations differ in model, prompt wording and read-out (\cref{app:reextract}). \emph{Intersection} labels keep a label only where the pipelines agree (disagreements become \textsc{unclear}), and
500 \emph{soft-label draws} sample labels from the second pipeline's probabilities. An
LLM-independent numeric anchor agrees with the primary labels on 0.807 \ci{0.782}{0.830} of 1{,}841
high-confidence records (1.6\% of associations) and with the second pipeline's on 0.758
\ci{0.733}{0.782}; the anchor is a selected consistency check, not corpus-wide accuracy
(\cref{app:anchor}).

\paragraph{Endpoint.} For a cutoff year $t$, let $E(c,t)$ hold the claim's records published before
$t$ and $E^{+}(c,t)$ those published at or after $t$; the latter never enter an input. With $h$ and
$l$ the numbers of \textsc{harmful} and \textsc{protective} records in $E(c,t)$, the pooled
direction is $p_t=(h-l)/(h+l)$ when $h+l\ge8$ and 0 otherwise. The future direction $p^{+}_t$ is the
same ratio over the signed records of $E^{+}(c,t)$ with no minimum count (0 when none follows; every eligible source unit has at least 11), and the endpoint is
\begin{equation}
R(c,t)=\mathbf{1}\big[\sgn(p^{+}_t)\neq\sgn(p_t)\big],
\label{eq:target}
\end{equation}
with $\sgn(0)=0$. \textsc{null} records are excluded from $p_t$, so $|p_t|$ near zero means a
balance of signed labels, not the absence of an association. A zero direction on exactly one side
is an event, so 11 of the 71 events are such units (seven pre-cutoff exact ties, one sparse
history, three post-cutoff ties); nothing bounds the later majority, and the horizon runs to the
end of the corpus. The endpoint stays primary because every cell was computed on it; tie-free,
minimum-signed-count and fixed-horizon variants reproduce the development ranking
(\cref{app:robust}). A unit is eligible when at least 20 records follow the
cutoff: 1{,}012 eligible units in 160 claims with 71 events in 23 claims (\cref{app:bench}), with
repeated cutoffs of a claim clustered in every interval.

\paragraph{Evaluation populations.} The \emph{development} cell is a retrospective claim-held-out
backtest on the 1{,}012 eligible units, with training outcomes extending to the end of the corpus. The \emph{calendar
cohort} is origin-defined: at origins $Y\in\{2004,2007,\ldots,2022\}$ and horizons
$H\in\{3,5,10\}$, a claim is a candidate if it has at least 40 records before $Y$, the window
$[Y,Y+H)$ must hold at least ten signed records, training is claim-disjoint and uses only outcomes
observable by $Y$, and only windows ending by 2025 are evaluated (\cref{app:cohort}); the design is calendar-respecting for documents, not for the annotator.
\emph{Frozen transfer} fits once on the source and applies the score, its calibrator and its count
model unchanged to External-1 (physical activity, sleep, environment) and External-2 (medication,
screening, occupational exposures, clinical risk states).

\section{Evaluation Protocol}
\label{sec:protocol}

\subsection{The composition references}
\label{sec:bb}
The \emph{margin rule} scores a unit by $1-|p_t|$ and needs no training. The \emph{stationary
reference} gives the endpoint's probability if a claim's signed labels were
exchangeable: suppose they are conditionally i.i.d.\ given $\theta=P(\textsc{harmful})$ with
$\theta\sim\mathrm{Beta}(\alpha,\alpha)$, and let $K$ count \textsc{harmful} labels among the next
$m$ signed records.
\begin{lemma}[stationary sign change]
\label{lem:stationary}
Given $(h,l,m)$, $K\sim\BB(m,h+\alpha,l+\alpha)$, and the probability that the future majority
differs from $\sgn(h-l)$ is
\begin{equation}
\pi_t(m)=\sum_{k\in\mathcal{K}_t}\BB(k;m,h{+}\alpha,l{+}\alpha),
\label{eq:bb}
\end{equation}
with $\mathcal{K}_t=\{k:\sgn(2k-m)\neq\sgn(h-l)\}$.
For fixed $n=h+l$, $m$ and $\alpha$, $\pi_t$ is non-increasing in $|h-l|/n$ among non-tied
histories, so when $n$, $m$ and $\alpha$ are common across units the margin rule ranks units as
$\pi_t$ does.
\end{lemma}
The proof and numerical checks are in \cref{app:bb}. \Cref{eq:bb} is stated for the raw majority
$\sgn(h-l)$; the scored reference replaces it by the operational sign $\sgn(p_t)$ of
\cref{eq:target}, which differs only on sparse histories, where $p_t$ is gated to 0
(\cref{app:bb}). The analytic reference is $\pi_t$; the margin rule is its untrained surrogate,
exact only under common $(n,m,\alpha)$, so both are reported. The \emph{realised-count} version
inserts the observed future $m$ (a diagnostic); the \emph{count-predicted} version averages
$\pi_t$ over draws of $m$ from a count model fitted on the training claims; both use $\alpha=1$.

\subsection{Scores under test}
The \emph{signed two-variable score} is a class-balanced $L_2$ logistic regression on standardised
$[p_t,a_t]$, where $a_t$ is the majority share among signed records, or the \textsc{null} share when
\textsc{null} records dominate. On the signed branch $a_t=(1+|p_t|)/2$, so its logit is
\begin{equation}
\logit s_2(c,t)=b_0+b_1p_t+b_2|p_t|,
\label{eq:logit}
\end{equation}
and $b_1$ is the only term through which the identity of the majority class enters. Its
\emph{class-blind variant} is the same model on $[|p_t|,a_t]$. The shipped agreement is
\emph{definition~A}, the archived reference: the \textsc{null} share when that share is at least
0.5, else the signed majority share. The \emph{branch-free variant} uses \emph{definition~C}, the
signed-only share $(1+|p_t|)/2$ on every unit (gated to 0 on sparse histories, as A is;
\cref{app:agreement}): a matched diagnostic intervention that changes the representation and
nothing else and inherits none of A's other results. Richer learned scores (\compact on seven
variables, eleven learners on all 55 pre-cutoff variables, order features added to the
stationary logit) stay at or below it (\cref{app:ablation,app:order}); prompted LLM forecasters read a digest
of the pre-cutoff labels (\cref{app:llm}).

Four \emph{null controls} rebuild the target from permuted or simulated labels (1{,}000 draws each, score refitted;
\cref{app:null}): Null~A permutes each claim's labels across its publication dates, Null~B draws
each label i.i.d.\ from the claim's label frequencies, Null~C permutes within claim and decade,
keeping each claim's era-level label mix, and Null~D draws labels from a per-claim logistic trend
of the \textsc{harmful} share on year, keeping the fitted drift without block conditioning.

\subsection{Checks, cells and tiers}
Each check has a primary cell and a pass condition on $\Delta$, the paired AUROC difference between
a score and a composition reference, with a claim-clustered percentile bootstrap of 2{,}000
resamples \citep{EfronTibshirani1993}; a difference is \emph{established} when its 95\%
interval excludes zero. With 23 positive clusters, percentile intervals can under-cover
\citep{Cameron2008}, so ``not established'' means insufficient evidence, never evidence of absence;
a wild cluster bootstrap-$t$ with Webb weights, a clustered-ROC sensitivity and a claim-level
bootstrap that refits the score in every resample leave every tested superiority and loss decision
unchanged, and one non-inferiority subdecision (External-1) procedure-sensitive
(\cref{app:metrics}), which shows stability to those choices and nothing more. A superiority check needs an established gain over both
the margin rule and the count-predicted stationary reference; a non-inferiority check is judged against the margin rule. We fixed the cells after observing the baseline results of
\cref{sec:results}, for which they are a post hoc description.
\begin{enumerate}[leftmargin=*,nosep,label=\arabic*.]
\item \textbf{Composition.} In the development cell, $\Delta$ is established as positive against both references.
\item \textbf{Annotators.} Under the second annotation and intersection labels, $\Delta$ against the margin rule is non-inferior: its lower 95\% bound exceeds $-\delta$, $\delta=0.02$, a stated convention below the smallest established difference in \cref{tab:protocol} (0.028).
\item \textbf{Calendar.} In the calendar cohort at $H=5$, $\Delta$ is established as positive against both references, with no established loss at $H=3$ or $H=10$.
\item \textbf{Class invariance.} Over pairs of units that share a pre-cutoff majority class ($p_t=0$ excluded), $\Delta$ is established as positive against both references (pooled), and the point difference is non-negative within each class; the class-blind variant is a diagnostic.
\item \textbf{Transfer.} In both external domains $\Delta$ against the margin rule is non-inferior at the same $\delta$.
\end{enumerate}
Passing all five establishes predictive value beyond composition, class information and one
annotator, not a legitimate forecaster. Checks~2 and~5 are non-inferiority conditions because
superiority is barely testable there (the margin rule already reaches 0.953 under the second
annotation; External-1 has 19 events); the margin $\delta$ makes a wide interval fail rather than
pass. A score
reaches \emph{Tier~A} (development gain) by passing check~1, \emph{Tier~B} (development gain,
non-inferior across annotations and domains) by passing checks 1, 2 and 5, and
\emph{Tier~C} (class-invariant calendar gain) by passing all five with a development $\Delta$ above
the 97.5th percentile of the paired differences the same score attains under Null~C. For a method fixed before its outcomes on the released cells are examined, checks 1, 3 and 4 form an
intersection-union test at joint level 5\% \citep{Berger1982}; the baselines here, developed
against the released outcomes, forfeit that guarantee. Label-feature and text-reading methods form two \emph{tracks} with different eligible checks (\cref{app:repro}).

\section{Results}
\label{sec:results}
\begin{table*}[!t]\centering\scriptsize
\setlength{\tabcolsep}{4pt}
\begin{tabular}{@{}llrrrrl@{}}
\toprule
Score & AUROC [95\% CI] & AUPRC & Brier & Log loss & Recall@10\% & $\Delta$ vs.\ margin [95\% CI] \\
\midrule
Margin rule $1-|p_t|$ (untrained) & .881 [.824, .936] & .468 & .0515 & .184 & .549 & reference \\
Stationary, count-predicted & .909 [.867, .949] & .497 & .0499 & .200 & .549 & $+$.028 [.008, .050] \\
\midrule
Signed two-variable ($p_t$, $a_t$) & .948 [.921, .971] & .535 & .0452 & .146 & .676 & $+$.067 [.017, .120] \\
Class-blind ($|p_t|$, $a_t$) & .897 [.839, .945] & .445 & -- & -- & -- & $+$.016 [$-$.017, .057] \\
\bottomrule
\end{tabular}
\caption{Development cell: 1{,}012 eligible units, 71 events in 23 claims (prevalence 0.070),
primary annotation, leave-one-claim-out. The margin rule is untrained as a ranking; its Brier and
log loss, like those of the learned rows, use a training-only Platt mapping \citep{Platt1999} of the
score to a probability, whereas the stationary row is uncalibrated (recalibrated: \cref{tab:matched}).
$\Delta$ is paired over 2{,}000 claim resamples; ``--'': not computed. The realised-count stationary
diagnostic (.882, $+$.000 \ci{-.033}{.034}) and the branch-free variant (.925, $+$.044
\ci{-.005}{.091}) are in \cref{tab:stationary,tab:protocol}; \compact (.946), the order
model (.915) and eleven further learners: \cref{app:metrics,app:ablation}.}
\label{tab:core}
\end{table*}

\begin{table*}[!t]\centering\scriptsize
\setlength{\tabcolsep}{3pt}
\begin{tabular}{@{}lllccll@{}}
\toprule
Check & Cell & Units (ev.) & Margin & Score & $\Delta$ vs margin [95\% CI] & Stationary vs margin [95\% CI] \\
\midrule
1 Composition & Development, A & 1{,}012 (71) & .881 & .948 & $\mathbf{+.067}$ [.017, .120] & $+$.028 [.008, .050] \\
 & Development, C & 1{,}012 (71) & .881 & .925 & $+$.044 [$-$.005, .091] & $+$.028 [.008, .050] \\
\midrule
2 Annotators & Second, A & 1{,}012 (52) & .953 & .857 & $\mathbf{-.096}$ [$-$.166, $-$.037] & $+$.010 \\
 & Second, common supp., A$^{\dagger}$ & 1{,}002 (42) & .942 & .758 & $\mathbf{-.184}$ [$-$.268, $-$.105] & $+$.012 \\
 & Second, common supp., C$^{\ddagger}$ & 1{,}002 (42) & .942 & .940 & $-$.002 [$-$.037, .023] & -- \\
 & Second, C & 1{,}012 (52) & .953 & .952 & $-$.001 [$-$.028, .018] & -- \\
 & Intersection, A & 1{,}012 (55) & .928 & .898 & $-$.030 [$-$.094, .023] & $+$.019 \\
 & Intersection, C & 1{,}012 (55) & .928 & .935 & $+$.007 [$-$.016, .030] & -- \\
\midrule
3 Calendar & $H{=}3$ & 756 (58) & .907 & .913 & $+$.006 [$-$.012, .023] & $+$.013 [$-$.002, .031] \\
 & $H{=}5$ (primary cell) & 625 (45) & .897 & .920 & $+$.023 [$-$.005, .053] & $+$.003 [$-$.021, .030] \\
 & $H{=}10$ & 492 (35) & .867 & .876 & $+$.008 [$-$.032, .047] & $-$.004 [$-$.040, .034] \\
\midrule
4 Class inv. & Within \textsc{prot}-majority & 730 (21) & .959 & .969 & $+$.011 [$-$.001, .024] & -- \\
 & Within \textsc{harm}-majority & 274 (42) & .833 & .863 & $+$.030 [$-$.041, .114] & -- \\
 & Pooled within-class pairs & 1{,}004 (63) & .909 & .927 & $+$.018 [$-$.012, .060] & -- \\
\midrule
5 Transfer & External-1, frozen & 322 (19) & .893 & .948 & $+$.055 [$-$.052, .106] & $+$.013 [$-$.010, .039] \\
 & External-2, frozen & 1{,}257 (95) & .896 & .823 & $\mathbf{-.074}$ [$-$.123, $-$.024] & $+$.002 [$-$.014, .017] \\
\bottomrule
\end{tabular}
\caption{The protocol applied to the signed two-variable score (``Score''; AUROC; bold:
established). A: the shipped two-branch agreement input; C: the signed-only (branch-free) input,
a matched intervention that changes nothing else (\cref{app:agreement}). Annotation rows rebuild
inputs and targets and refit; calendar rows refit per origin; external rows apply the source-fitted
model unchanged. Stationary AUROCs are in \cref{tab:stationary,app:reextract}; entries without an
interval are point differences of the reported AUROCs; against the stationary reference the
shipped score gains $+0.039$ \ci{0.004}{0.080} in development. Of checks 2 and 5, only the
intersection C row passes $\delta=0.02$; the second-annotation C row misses it by 0.008, the
common-support C row by 0.017, and every A row by at least 0.03. $^{\dagger}$At least eight signed
pre-cutoff records under both pipelines, refitted (unrefitted: .823, $-$.119). $^{\ddagger}$The
row above with A replaced by C on the same units and outcomes; C $-$ A is $+0.182$
\ci{0.120}{0.248} (\cref{tab:branchfree}).}
\label{tab:protocol}
\end{table*}

\subsection{Composition explains most of the ranking}
\label{sec:pred}
\Cref{tab:core} gives the development cell. The untrained margin rule ranks later change almost as
well as any score; the stationary reference with realised counts is indistinguishable from it, as
\cref{lem:stationary} anticipates, and its count-predicted form gains an established 0.028 from the
accrual model, not from \cref{eq:bb}. The lowest quintile of $|p_t|$ holds 54 of the 71 events
(\cref{app:robust}). Richer inputs do not move beyond two variables: \compact matches
the signed score to within 0.002, eleven learners on all 55 variables stay at or below it, and the
order features add nothing established (\cref{app:ablation,tab:matched}); shared
identifiers and single positive claims do not drive the development conclusion (\cref{app:dependence}).

\subsection{Most of the learned gain is a majority-class prior}
\label{sec:prior}
\begin{table}[!t]\centering\scriptsize
\setlength{\tabcolsep}{4pt}
\begin{tabular}{@{}lr@{}}
\toprule
Quantity & Value [95\% CI] \\
\midrule
\multicolumn{2}{@{}l}{\emph{Pair decomposition of $\Delta$AUROC (\cref{tab:decomp})}} \\
Across-class pairs (34{,}650) & $+$.0609 [.0159, .1081] \\
Within-class pairs (24{,}633) & $+$.0068 [$-$.0039, .0195] \\
Zero-direction pairs (7{,}528) & $-$.0010 [$-$.0032, $-$.0001] \\
Across-class share of the gain (.0667) & 91\% \\
Within-class loss, 2nd ann. ($-$.0957) & $-$.0781 [$-$.1247, $-$.0386] (82\%) \\
\midrule
Event rate, \textsc{harm}-majority & 42/274 (15.3\%) \\
Event rate, \textsc{prot}-majority & 21/730 (2.9\%) \\
Median $|p_t|$, \textsc{harm} / \textsc{prot} & 0.588 / 0.612 \\
Class odds ratio, adj.\ for $|p_t|$ & 7.80 [2.01, 33.82] \\
Fitted $\hat b_1$ / $\hat b_2$ (\cref{eq:logit}) & $+5.3$ ($+3.0$ std.) / $-12.2$ \\
Slope in $|p_t|$, \textsc{harm} / \textsc{prot} & $-6.9$ / $-17.5$ ($<0$ in 160/160 folds) \\
\midrule
Signed $-$ class-blind, all pairs & $+$.051 [.013, .092] \\
\bottomrule
\end{tabular}
\caption{Where the learned gain comes from (development cell, primary annotation; differences
against the margin rule are in \cref{tab:protocol}). The decomposition partitions the 66{,}811
positive-negative pairs by whether the two units share a pre-cutoff majority class; contributions
are additive and sum to the unrounded totals ($+0.0667$, $-0.0957$), intervals are claim-clustered,
and the shares are descriptive. \textsc{harm}, \textsc{prot}: pre-cutoff majority class ($p_t\neq0$;
1{,}004 units, 63 events). Odds ratio: class-balanced logistic regression on class and $|p_t|$,
claim-clustered (crude 6.11, Mantel-Haenszel 4.35; \cref{tab:classquint}); slopes:
$\hat b_2{\pm}\hat b_1$ over leave-one-claim-out folds. Sign-exchange and orientation-clean
results: \cref{app:signinv,app:orientation}.}
\label{tab:prior}
\end{table}
\Cref{tab:prior} locates the gain. An exact decomposition of the signed-minus-margin AUROC
difference by pair type places $+0.0609$ \ci{0.0159}{0.1081}, about 91\% of the estimated total,
in positive-negative pairs whose units have different pre-cutoff majority classes; the additive
within-class contribution, $+0.0068$ \ci{-0.0039}{0.0195}, is not established. The incremental gain is therefore
concentrated in across-majority-class comparisons, and the class-rate difference persists within
the reported margin strata: \textsc{harmful}-majority units change about five times as often as
\textsc{protective}-majority units, \textsc{harmful} majorities are not thinner (median $|p_t|$ 0.588 against 0.612), the class odds
ratio adjusted for $|p_t|$ is 7.80 \ci{2.01}{33.82}, and no \textsc{protective}-majority event is
observed above the lowest margin quintile (\cref{tab:classquint}). On the signed branch
\cref{eq:logit} is linear in $|p_t|$ within a class, with slope $b_2+b_1$ for \textsc{harmful} and
$b_2-b_1$ for \textsc{protective} majorities; the full-data fit gives $\hat b_2=-12.2$, so both
slopes are negative ($-6.9$ and $-17.5$, in every leave-one-claim-out fold), and $\hat b_1$ encodes
that \textsc{harmful} majorities change more often. The class-rate difference has an era-level
correlate: under the primary annotation the \textsc{harmful} share of signed labels falls within
claims by $-0.042$ \ci{-0.061}{-0.023} per decade (\cref{app:controlled}), so an early
\textsc{harmful} majority is the one most likely to be outvoted later, and a control preserving this drift reproduces the gain (\cref{sec:null}).
Whether the drift is the annotator's is not settled: on anchored records from 2000 onward the
within-claim slopes of the \textsc{harmful} share are negative for the anchor and both pipelines
($-0.053$, $-0.070$, $-0.043$ per decade), but only the primary's interval excludes zero and the
paired slope differences cover zero (\cref{tab:anchorera}).

\subsection{The learned gain is not consistent across conditions}
\label{sec:transport}
\Cref{tab:protocol} applies the five checks. The signed score passes check~1 against both
references and no other check; the calendar cells establish no gain, External-2 establishes a
loss, and the loss under the second annotation is largely representation-dependent.

\emph{Annotators.} Under the second annotation the ordering against the margin rule reverses:
rebuilt and refitted, the shipped score falls below the margin rule by an established margin, while
$\hat b_1$ keeps its sign and shrinks ($+5.3$ to $+2.1$); the loss survives restriction to the
common support and follows the annotator model (\cref{tab:controlled}); 82\% of it arises within
majority classes, so it is not the across-class advantage reversing. Agreement construction
contributes substantially to it. On the
1{,}002-unit cohort with adequate signed support under both pipelines, replacing the two-branch
agreement input (A) by the signed-only share (C) raises the second-annotation AUROC from 0.758 to
0.940, a paired improvement of $+0.182$ \ci{0.120}{0.248}, against a margin rule of 0.942 on the
same cohort and outcomes ($-0.002$ \ci{-0.037}{0.023} against it, which misses $\delta$ by 0.017),
so the point deficit nearly disappears with no outcome changed. The
improvement is largest on the \textsc{null}-dominant branch, on which $a_t$ is a \textsc{null}
share ($+0.246$ \ci{0.073}{0.405}, against $+0.041$ on the signed branch; \cref{tab:branchfree}).
Refitted on the full cohort, the branch-free
score's difference from the margin rule is $-0.001$ \ci{-0.028}{0.018} under the second annotation
and $+0.007$ \ci{-0.016}{0.030} under intersection labels: no loss is established in these two
annotation comparisons, and non-inferiority at $\delta$ is established for intersection labels but
not for the second annotation (a representation effect within that setting, not a proof that
annotation-induced signed membership plays no role). On the primary annotation the branch-free score establishes no
development gain, $+0.044$ \ci{-0.005}{0.091}, or $+0.051$ \ci{-0.006}{0.108} with the
\textsc{null} share as a third input, but its direct difference from the shipped score is not
established either ($-0.023$ \ci{-0.050}{0.002}; \cref{tab:astra2}), so the branch is not shown to
carry the primary signal; what is shown is that the only established gain over composition in
\cref{tab:protocol} comes from an input whose definition does not port across annotations.
Separately, the second annotation agrees less with the numeric anchor (\cref{sec:task}): $-0.049$
\ci{-0.068}{-0.030} against the primary, $-0.059$ \ci{-0.081}{-0.036} on the ratio-scale subset;
among signed anchors it assigns \textsc{null} more often, 9.4\% against 4.2\%, while opposite-sign
assignments have similar frequencies, 7.9\% against 7.6\% (\cref{tab:anchor2}): a plausible
route into the \textsc{null}-dominant branch, not a demonstrated mediation, and no verdict on
which pipeline defines the correct claim-level events.

\emph{Calendar.} In the origin-defined cohort no difference is established at any horizon
(\cref{app:cohort}); at five years the estimated advantage is
$+0.023$ \ci{-0.005}{0.053}, neither a gain nor a loss. No tested text-reading baseline (prompted 70B and 72B models, abstract embeddings) establishes a gain over composition (\cref{app:text}).

\emph{Class invariance and transfer.} Check~4 is not passed: over pairs that share a majority
class the gain is not established, in either class or pooled, and the class-blind diagnostic shows
no established gain (\cref{tab:prior}). In External-1 the signed score's estimate, $+0.055$
\ci{-0.052}{0.106}, establishes neither superiority nor non-inferiority; in External-2 it loses by an established margin in AUROC and in AUPRC (\cref{app:crossdomain}), so the prior does not transfer frozen; refitted inside External-2 the
same two inputs lose only $-0.017$ \ci{-0.047}{0.009}, so the gap is a transfer loss, not an
inability of the inputs to rank that domain. The signed score reaches Tier~A only; its branch-free form reaches no tier; no tested score reaches Tier~B or Tier~C.

\begin{table*}[!t]\centering\footnotesize
\setlength{\tabcolsep}{5pt}
\begin{tabular}{@{}lrllll@{}}
\toprule
Quantity & Obs. & Null~A & Null~B & Null~C & Null~D \\
\midrule
Events, primary & 71 & 35.6 [26,46]$^{\ast}$ & 36.5 [20,52]$^{\ast}$ & 61.2 [53,69]$^{\ast}$ & 70.1 [54,87] \\
Events, second & 52 & 34.3 [24,45]$^{\ast}$ & 36.0 [21,51]$^{\ast}$ & 51.3 [45,57] & -- \\
Events, intersect. & 55 & 40.2 [29,51]$^{\ast}$ & 40.9 [25,56] & 55.7 [49,62] & -- \\
\midrule
Signed AUROC$^{\ddagger}$ & .948 & .948 [.907,.974] & .947 [.898,.975] & -- & -- \\
Margin AUROC$^{\ddagger}$ & .881 & .957 [.929,.977]$^{\dagger}$ & .957 [.924,.979]$^{\dagger}$ & -- & -- \\
$\Delta$, primary$^{\ddagger}$ & $+$.067 & $-$.009$^{\ast}$ & $-$.010$^{\ast}$ & $+$.065 [.047,.083] & $+$.040 \\
\bottomrule
\end{tabular}
\caption{Observed streams against the null controls (null mean [95\% range], 1{,}000 draws;
$\Delta$: signed minus margin, paired within draws). $^{\ast}$observed above the 95\% null range
($p<.025$; intersection under Null~B $p=.054$); $^{\dagger}$below it ($p\le.002$);
$^{\ddagger}$archived run (Null~A/B event means 35.3, 36.1); other cells: reimplementation
(\cref{app:null}).}
\label{tab:null}
\end{table*}
\subsection{Annotation configuration can redefine the positive class}
\label{sec:annot}
The two annotations agree on most labels yet disagree on most positive events (\cref{tab:extractors}): 93.2\% of units receive the same
endpoint value, mostly shared non-events, while the positive class is largely redefined, on common
support, under intersection labels and across soft-label draws alike. A controlled re-annotation
that changes one factor at a time (\cref{tab:controlled}) attributes most of the redefinition to
the model family rather than the prompt or the read-out, and a re-run of the primary configuration
with a different batch size alone moves the event set (71 against 73 events, Jaccard 0.846).
Agreement on labels does not imply agreement on a thresholded target (\cref{app:cumulative}); the
configuration that redefines the positive class most is also the one that agrees less with the
numeric anchor (\cref{sec:transport}), evidence of a specific annotation discrepancy on the covered
records, not of which pipeline defines the correct events; and repairing the predictor's
representation changes none of this, since the 27 shared events are a property of the targets.
Annotation configuration changes hypothesis-test outcomes \citep{Baumann2025}; here it changes the
positive class.

\begin{table}[!t]\centering\scriptsize
\setlength{\tabcolsep}{2pt}
\resizebox{\columnwidth}{!}{%
\begin{tabular}{@{}lrrrrr@{}}
\toprule
Annotation & Events (cl.) & Shared & Jaccard & Sparse & Prim.-trained \\
\midrule
Primary (Llama) & 71 (23) & 71 & 1.000 & 1 & .948 \\
Second (Qwen) & 52 (16) & 27 & .281 & 10 & .878 \\
Intersection & 55 (19) & 30 & .313 & 14 & .869 \\
Common support, prim./sec. & 67 / 42 & 23 & .267 & -- & -- \\
Soft-label draws (500) & 43-57 & -- & -- & -- & -- \\
\bottomrule
\end{tabular}}
\caption{Positive events under each annotation of the same 1{,}012 units (label agreement 67.0\%, $\kappa=0.52$; pooled directions $r=0.943$; \cref{app:reextract}). ``Shared'' and Jaccard are relative to the primary
annotation; ``Sparse'': events on histories with fewer than eight signed records; ``Prim.-trained'': the primary-trained score evaluated unchanged against each endpoint.}
\label{tab:extractors}
\end{table}

\subsection{Excess reversals and the drift behind them}
\label{sec:null}
The composition-preserving controls (A and B; \cref{tab:null}) reproduce the fitted
discrimination but not the events: the streams reverse
about twice as often, an excess that also appears under the second annotation and, against Null~A,
under intersection labels. Null~C absorbs about 72\% of the primary excess and all of it under the
other two annotations, and reproduces the signed score's paired gain, which Nulls A and B do not;
the observed $+0.067$ lies below its 97.5th percentile (0.083), so the Tier~C null condition also
fails. Null~C conditions on claim-decade totals; only 9.6\% of units, holding 3 of the 71
primary events, have counts no permutation can change (\cref{app:null}), and what it leaves is
partition-dependent: with five-year blocks the
observed count falls inside the null range (64.9 \ci{59}{71}, $p=0.027$), with shifted boundaries
likewise (64.5 \ci{58}{71}, $p=0.038$), with twenty-year blocks the gap widens (56.3 \ci{48}{64},
$p=0.002$), and the second-annotation and intersection counts lie inside the range in all four
partitions (\cref{tab:nullblocks}). Null~D, a retrospective control that keeps each claim's fitted
\textsc{harmful}-share trend without block conditioning, is not rejected on the count (70.1
\ci{54}{87}, $p=0.458$) or on the gain ($+0.040$ against $+0.067$, $p=0.099$), without being shown to be
the generative explanation. The excess over era-conditioned controls is confined to the primary
annotation and is a partition-specific contrast, not an identified set of unexplained reversals;
across partitions, learned-score differences are evaluated under the primary annotation only.

\section{Conclusion}
\label{sec:conclusion}
On a twice-annotated PubMed stream, an untrained composition rule ranks later change at 0.881;
about 91\% of the best learned gain lies in comparisons across majority classes; most of that
score's loss under a second annotation is repaired by a signed-only agreement input that leaves
the outcomes untouched, while that input establishes no development gain; and the two
annotations share 27 of 96 positive events. A forecasting benchmark built from automatic
annotations must therefore separate why a score succeeds, why it fails and whether its events are
stable; \method is released as a challenge set for that, on which a method must satisfy each
applicable check's comparison rule on every annotation and report any threshold-switching input in
branch-free form. For any benchmark of this construction the audit reduces to five reports: a
composition reference beside every learned score, the gain decomposed by pair type, the target
rebuilt under a second annotation with event overlap, an annotation-independent check of both, and
event counts under composition- and trend-preserving controls.

\section*{Limitations}
The target is change in machine-generated direction labels, not expert-verified conclusion change,
because literature-scale forecasting data are built from automatic annotation; no human label
enters the benchmark, and the released adjudication frame (96 union events, 50 matched non-events)
is unexecuted, so the resource is an audit of a construct, not a validated one. The numeric anchor
covers 1.6\% of associations, is a record-level check that cannot validate the claim-level event
set, is weakest on \textsc{null}, and its era comparison starts in 2000 with an imprecise anchor
trend, so the source of the era drift is not identified. The two annotations differ in model,
prompt and read-out at once and span two 7--8B families, so the agreement intervention establishes
a representation-dependent contribution to the second-annotation loss, not its whole cause: the
branch-free score is non-inferior under neither the second annotation nor common support, and its
primary-annotation difference from the shipped score is not established, so the branch is not
shown to carry the primary signal. Both annotators were pretrained after most abstracts they label,
so the era diagnostics do not separate prior knowledge from change in the literature; Null~D is a
retrospective parametric control; the partition grid tests counts under three annotations but
learned-score differences under one; and the historical-case audit supports no timing claim. With
71 events in 23 claims, comparisons outside the primary cells are uncorrected, the percentile
bootstrap can under-cover, the additional procedures show decision stability rather than coverage,
the cells were fixed after the baseline results, and the text-reading baselines include no
fine-tuned model, so text is not shown to be uninformative. Retrieval is documented, not
regenerable: the claim-list assembly is unrecorded, 19 stored queries are unconfirmed, 48 claims
are capped, and the external corpora ship as primary-label counts without documents. The analytic
reference assumes exchangeable labels and a flat prior, and the calendar cohort conditions on
later accrual and is calendar-respecting for documents only. No clinical validity is claimed.

\section*{Ethical Considerations}
The benchmark uses public PubMed titles, abstracts and publication-level results, with no individual
patient data. Direction labels describe reported associations relative to a named exposure and
outcome; they do not establish clinical desirability, causal validity or personalised advice. A
high score denotes predicted change in an automatic annotation stream under the stated pipeline,
not evidence that an exposure is harmful, and this paper shows that such scores can encode
annotation-specific priors. Health communication that uses such signals should retain the cutoff
date, the uncertainty, the underlying publications and expert review. The resource is intended for
research on evidence streams, not for automated clinical decisions.

\bibliography{references_final}

\appendix
\renewcommand{\thesection}{\AlphAlph{\value{section}}}
\crefalias{section}{appendix}
\crefalias{subsection}{appendix}

\section{Benchmark Construction and Retrieval Audit}
\label{app:bench}
\paragraph{Fixed queries.} Each claim is an exposure-outcome pair. Its PubMed query is generated from
the two terms as \texttt{("<exposure>"[MeSH Terms] OR "<exposure>"[Title/Abstract]) AND
("<outcome>"[MeSH Terms] OR "<outcome>"[Title/Abstract])}. Retrieval uses ESearch with the history
server and EFetch in blocks of 500, without date or publication-type filters, so
the same search serves every cutoff of a claim. Retrieval stopped at a per-query cap of 800 records
for 48 queries, as recorded in the archive metadata; for those claims the archived set is what the
cap left (\cref{app:snapshot}). The archived snapshot holds the retrieval history as stored
for each fixed query: 112{,}453 claim-record associations over 96{,}091 distinct identifiers (median
737 records per claim). The dated re-retrieval of \cref{app:snapshot} returns a different identifier set for some claims,
so the archived sets, not the queries, define the benchmark. Nothing here implies that a mechanical query recalls every
scientifically relevant paper; a record belongs to every claim whose query returns it, and 12{,}411
identifiers (12.9\%) belong to two or more claims (\cref{app:dependence}). The exact retrieval timestamp is not recorded in the released metadata. Two post hoc
observations bound the original retrieval to 2026-09-09 through 2026-09-10 13:53 UTC. The highest archived
identifier (42{,}715{,}502) entered PubMed on 2026-09-09, and a harvest cannot predate the
arrival of a record it contains; the earliest available write of the 165 archived
record files is dated 2026-09-10 13:53~UTC, and a copy cannot precede what it copies. The lower
bound does not assume that identifiers are assigned in order of entry: it requires only that
these records are in the archive. The file times are bulk writes rather than a record-by-record
harvest, so they bound the harvest from above without dating it. These bounds constrain the
original harvest but do not make the archived identifier sets regenerable, since re-running the
stored queries returns a different identifier set (\cref{app:snapshot}). The corpus holds 5{,}818 records dated 2026 and two dated 2027
(ahead-of-print entries), so 2026 is a partially observed year; the fixed-horizon analyses of
\cref{app:robust,app:cumulative} treat 2026 as the last calendar year, whereas the origin-defined cohort of \cref{app:cohort} uses 2025 as the last complete year. The 145-claim candidate list is supplied
as a fixed artefact: the procedure that assembled it is not recorded in the released metadata, and
the 20 hand-added claims are listed below.

\paragraph{Cutoffs.} A claim enters with at least 45 records that carry a year and an
extractor-assigned direction. With $y_{(40)}$ the year of its 40th record in date order and
$y_{\max}$ the year of its last record, the grid is the set of distinct integers among seven
equally spaced points from $y_{(40)}+2$ to $y_{\max}-1$, provided $(y_{\max}-1)-y_{(40)}\ge4$. A
unit at cutoff $t$ sees every record with publication year strictly less than $t$. Every unit has at
least 40 pre-cutoff records; the median unit has 194 pre-cutoff and 287 post-cutoff records; the
median pre-cutoff $I^2$ is 0.935.

\paragraph{Post-cutoff fields and the target.} For every unit the post-cutoff records give
$n_{\text{post}}$, the post-cutoff pooled direction $p^{+}_t$ over their non-null labels, the
post-cutoff agreement, and $R$ under the sign convention of \cref{eq:target}. The fixed
horizons recompute the target inside $[t,t+H)$ for $H\in\{3,5,10\}$ with the same convention,
requiring at least ten non-null records in the window and a literature that reaches $t+H-1$.

\paragraph{Denominators.} The future-eligible sample is the 1{,}012 units with at least 20 post-cutoff
records (160 claims; five claims have no cutoff meeting the threshold), with 71 events in 23
claims. The 145-claim subset (948 units) has 885
future-eligible units with 65 events in 21 claims.

\paragraph{Hand-added claims.} Twenty claims were added by hand to the 145-claim candidate list:
added sugar $\rightarrow$ diabetes, coffee $\rightarrow$ cardiovascular disease, dairy $\rightarrow$
all-cause mortality, dietary cholesterol $\rightarrow$ cardiovascular disease, dietary pattern
$\rightarrow$ diabetes, eggs $\rightarrow$ cholesterol, folate $\rightarrow$ cardiovascular disease,
folate $\rightarrow$ neural tube defects, folic acid and B12 $\rightarrow$ lung cancer, folic acid
$\rightarrow$ colorectal cancer, hormone replacement $\rightarrow$ cardiovascular disease, legumes
$\rightarrow$ cardiovascular disease, milk $\rightarrow$ cardiovascular disease, niacin $\rightarrow$
cardiovascular disease, plant-based diet $\rightarrow$ cancer, selenium $\rightarrow$ prostate
cancer, sugar-sweetened beverages $\rightarrow$ obesity, trans fat $\rightarrow$ cardiovascular
disease, vitamin A $\rightarrow$ cancer and vitamin E $\rightarrow$ cardiovascular disease. Their
entry route is not recoverable, so the claim list is documented rather than regenerable. The analyses repeated on the 145-claim subset are those with an explicit 145-claim row or column in this appendix (the \compact leave-one-claim-out result, the tie-convention table and the rule-comparator table); the operational stationary baselines, the re-annotation and the common-support sensitivity were run on the full 165-claim corpus only.

\paragraph{External corpora.} External-1 (physical activity, sleep and environment: 46 claims, 322
units, all future-eligible, 19 events in 7 claims) and External-2 (medication, screening and
prevention, occupational exposures and clinical risk states: 180 claims, 1{,}260 units, 1{,}257
future-eligible, 95 events in 40 claims under \cref{eq:target}) were built with the same
retrieval, extractor, cutoff grid, feature builder and target. No claim is shared with the source
benchmark. \Cref{tab:bench} summarises every evaluation population.

\begin{table}[!t]\centering\small
\setlength{\tabcolsep}{3pt}
\resizebox{\columnwidth}{!}{%
\begin{tabular}{@{}lrrr@{}}
\toprule
 & Source & External-1 & External-2 \\
\midrule
Claims & 165 & 46 & 180 \\
Claim-cutoff units & 1{,}078 & 322 & 1{,}260 \\
Eligible units ($\ge20$ later records) & 1{,}012 & 322 & 1{,}257 \\
Eligible claims & 160 & 46 & 180 \\
Events, primary annotation (claims) & 71 (23) & 19 (7) & 95 (40) \\
Event prevalence & 0.070 & 0.059 & 0.076 \\
\midrule
Calendar cohort, 3 / 5 / 10 y: units & \multicolumn{3}{l}{756 / 625 / 492 (source)} \\
Calendar cohort, 3 / 5 / 10 y: events & \multicolumn{3}{l}{58 / 45 / 35 (source)} \\
\bottomrule
\end{tabular}}
\caption{The benchmark under the primary annotation (size of every evaluation population). The source stream holds 112{,}453
claim-record associations over 96{,}091 PubMed records; external corpora are released as
claim-cutoff units with primary-pipeline labels, for frozen transfer only.}
\label{tab:bench}
\end{table}

\section{Direction Extraction and Feature Representation}
\label{app:extract}
\subsection{Closed-output extraction prompt}
Direction is extracted with Llama-3.1-8B-Instruct in bfloat16, in evaluation mode with no gradient
and no generation. The title and abstract are concatenated, truncated to the first 3{,}000
characters, rendered through the chat template with the system and user turns of
\cref{fig:prompt}, tokenised with a 1{,}600-token limit and right truncation, and scored in a
single forward pass with batch size 16 and left padding. The logits at the final position are
restricted to the first token of each of the four label strings (\textsc{protective},
\textsc{harmful}, \textsc{null}, \textsc{unclear}); the argmax is the extractor-assigned claim–record direction and the softmax over the four is its confidence. Under the primary tokenizer only \textsc{null} is a single token; \textsc{protective}, \textsc{harmful} and \textsc{unclear} are split, but their first tokens are distinct, so the read-out is a well-defined argmax over four distinct tokens rather than a full-string likelihood. The second pipeline scores complete label strings and is compared with its own first-token read-out in \cref{app:reextract}. Every prompt is checked to contain the
claim's own outcome term before scoring. Because the output space is closed there is no sampling
temperature, no decoding and no output parser. \textsc{unclear} labels are retained for audit and
excluded from directional aggregation. Effect sizes, confidence limits and design indicators are
recovered by case-sensitive pattern matching over the abstract text, not by the language model.
The feature builder places each parsed interval on a working scale inferred from its limits:
when both limits are positive the interval is treated as a ratio measure and represented as
$y=(\ln L+\ln U)/2$ with $SE=(\ln U-\ln L)/(2\times1.96)$; otherwise it is treated as a
difference measure on the linear scale. The \compact inputs $\mu_{\text{eff}}$,
$\sigma_{\text{eff}}$ and $n_{\text{eff}}$ (\texttt{val\_mean}, \texttt{val\_sd}, \texttt{n\_rr} in
\cref{tab:feature-inventory}) are therefore computed over a mixture of log-ratio and additive
quantities, not over uniformly commensurate log-ratios; the typed anchor of
\cref{app:anchor} reads the measure from the text instead and is used for validation only.

\begin{figure}[!t]
\small
\begin{tabular}{@{}p{0.95\linewidth}@{}}
\toprule
\textsc{system}: You classify the direction of association reported between an exposure and an
outcome. Answer with exactly one word: PROTECTIVE if the exposure is reported to reduce the outcome,
HARMFUL if it increases it, NULL if no association is reported, UNCLEAR if the study does not report
this association. \\[4pt]
\textsc{user}: Exposure: \{exposure\} \textbackslash n Outcome: \{outcome\} \textbackslash n
\textbackslash n Study: \{title\}. \{abstract\} \textbackslash n\textbackslash n Direction: \\
\bottomrule
\end{tabular}
\caption{The direction-extraction prompt. The four label tokens are scored at the position
following \texttt{Direction:}; nothing is generated.}
\label{fig:prompt}
\end{figure}

\subsection{In-domain numeric-anchor validation}
\label{app:anchor}
We construct an LLM-independent reference directly from the quantitative reporting in the corpus;
no language model and no human annotation is used to create it. Every ``95\% CI'' or ``95\%
confidence interval'' followed by two decimal limits $(L,U)$, $U>L$, is a candidate. A candidate is
\emph{typed} by the effect-measure token in the 120 characters before it (40 after): ratio measures
(odds, risk, hazard, rate and prevalence ratios and their abbreviations) fix the null at 1; additive
measures (mean, weighted mean and standardised mean differences, risk differences, regression
coefficients and ``difference'') fix it at 0. A candidate with no token is
\textsc{effect-type-unidentified}, one with tokens of both families \textsc{effect-type-conflict},
and a ratio interval with a non-positive limit \textsc{invalid}. A typed candidate is
\emph{attributed} to the claim when the sentence holding it, or the preceding sentence, names both
an exposure term and an outcome term from a per-claim lexicon derived mechanically from the claim
query; otherwise it is \textsc{anchor-ambiguous-relation}. The anchor direction of an attributed
interval is \textsc{null} if it covers the null, \textsc{harmful} if $L$ is above it and
\textsc{protective} if $U$ is below it; because the primary prompt defines \textsc{protective} as
the exposure reducing the outcome and \textsc{harmful} as increasing it, no judgement about the
valence of the outcome is involved, and the anchor measures agreement on the sign relative to the
null. A record is \emph{high-confidence} when it has at least one attributed typed interval and all
such intervals agree. Of 112{,}453 claim-record associations, 15{,}144 carry an interval, 10{,}650 a typed one,
2{,}384 an attributed typed one, and 1{,}841 in 148 claims are high-confidence (1.6\% of the
corpus; 543 excluded for disagreeing attributed intervals). \cref{tab:astra1} gives the
results against the primary labels; every interval is a claim-clustered percentile bootstrap with 2{,}000 resamples.

\begin{table*}[!t]\centering\scriptsize
\setlength{\tabcolsep}{3pt}
\resizebox{\textwidth}{!}{%
\begin{tabular}{@{}lrrrrrrrrrr@{}}
\toprule
Evaluation & Associations & Claims & Agreement & [95\% CI] & Bal.\ acc. & Macro-F1 & [95\% CI] & R$_\textsc{prot}$ & R$_\textsc{harm}$ & R$_\textsc{null}$ \\
\midrule
high-confidence anchor & 1841 & 148 & .807 & [0.782,\,0.830] & .754 & .766 & [0.735,\,0.796] & .923 & .836 & .504 \\
exactly one attributed interval & 999 & 133 & .770 & [0.742,\,0.796] & .716 & .723 & [0.687,\,0.756] & .906 & .818 & .424 \\
excluding extractor-UNCLEAR & 1832 & 147 & .811 & [0.787,\,0.833] & .758 & .768 & [0.735,\,0.798] & .924 & .841 & .510 \\
ratio-scale intervals only & 1691 & 134 & .813 & [0.789,\,0.836] & .757 & .768 & [0.732,\,0.800] & .929 & .863 & .479 \\
additive-scale intervals only & 167 & 62 & .737 & [0.650,\,0.819] & .655 & .649 & [0.560,\,0.726] & .889 & .270 & .806 \\
design: synthesis & 397 & 122 & .793 & [0.740,\,0.840] & .736 & .763 & [0.707,\,0.813] & .942 & .667 & .600 \\
design: other & 1355 & 130 & .810 & [0.782,\,0.835] & .738 & .748 & [0.706,\,0.785] & .922 & .858 & .434 \\
design: trial & 89 & 47 & .809 & [0.709,\,0.900] & .816 & .805 & [0.698,\,0.898] & .829 & .895 & .724 \\
\bottomrule
\end{tabular}}
\caption{LLM-independent numeric anchor against the primary labels. The effect measure is read from the text next to each 95\% interval and the interval is attributed to the claim only when the exposure and outcome terms both occur in the interval's sentence or in the preceding sentence. Claim-clustered 95\% intervals, 2,000 resamples.}
\label{tab:astra1}
\end{table*}

\paragraph{Confusion.} On the high-confidence records (rows anchor, columns extractor
\textsc{protective}/\textsc{harmful}/\textsc{null}/\textsc{unclear}): \textsc{protective} 679 / 29 /
27 / 1; \textsc{harmful} 84 / 627 / 35 / 4; \textsc{null} 97 / 75 / 179 / 4. Signed anchors are
recovered at 0.92 and 0.84; the weak class is \textsc{null} (recall 0.50), where the extractor
assigns a direction to intervals that cross the null in both directions, consistent with abstracts
that state a direction in words while the interval covers the null. Additive measures are rare
(167 records) and their \textsc{harmful} recall is 0.27, because ``difference'' intervals often carry
a sign convention (a reduction reported as a positive difference) that the anchor cannot resolve;
this is a limitation of the anchor. Read by column, 84 of the 860 records the primary pipeline
labels \textsc{protective} carry a \textsc{harmful} anchor (9.8\%), and 29 of the 731 it labels
\textsc{harmful} carry a \textsc{protective} anchor (4.0\%). By design, agreement is 0.79 for synthesis records, 0.81 for
observational or other records and 0.81 \ci{0.709}{0.900} for the 89 trial records.

\paragraph{Both annotations against the anchor.} \Cref{tab:anchor2} scores the second pipeline and
the intersection labels on the same 1{,}841 records. Paired on identical records, primary minus
second is $+0.049$ \ci{0.030}{0.068} on agreement and $+0.030$ \ci{0.012}{0.048} on balanced
accuracy, both established; on the 1{,}486 signed anchors, which removes the \textsc{null}-rate
confound, the gap is $+0.074$ \ci{0.053}{0.094}. By anchor scale, the paired second-minus-primary
agreement difference is $-0.059$ \ci{-0.081}{-0.036} on the 1{,}691 ratio-scale records, whose
anchors carry no sign-convention ambiguity, and $-0.010$ (point estimate; interval not computed) on
the 167 additive-scale records; the two subgroup estimates are not tested against each other. The
excess disagreement is a signed-versus-\textsc{null} discrepancy rather than sign inversion: among
signed anchors the second pipeline assigns \textsc{null} to 9.4\% against the primary's 4.2\%,
whereas opposite-sign assignments have similar frequencies, 7.9\% against 7.6\% (a descriptive
contrast; equality of the opposite-sign rates is not established), and its recall on true
\textsc{null} anchors rises only from 0.504 to 0.561. \textsc{null} is the prompt's substantive
``no association'' label, distinct from \textsc{unclear}, so these are label assignments, not
abstentions. Intersection labels retain a signed or \textsc{null} label on 1{,}587 of the 1{,}841
anchored records (86.2\% coverage) and turn the other 254 into \textsc{unclear}, which the anchor
cannot express; the intersection row scores those 254 as disagreements, so it is full-population
agreement with \textsc{unclear} counted as an error, not an accuracy bound of any kind, and on the
retained 1{,}587 records agreement is 0.842 \ci{0.821}{0.861}, higher than either pipeline's
because those are the records on which the two pipelines already agree. The second pipeline's
confusion on the 1{,}841 records (rows anchor, columns \textsc{protective}/\textsc{harmful}/
\textsc{null}/\textsc{unclear}, row totals 736/750/355): \textsc{protective} 637 / 30 / 66 / 3;
\textsc{harmful} 87 / 559 / 74 / 30; \textsc{null} 85 / 61 / 199 / 10; read by column, 87 of the
809 records it labels \textsc{protective} carry a \textsc{harmful} anchor (10.8\%, against the
primary's 9.8\%) and 30 of the 650 it labels \textsc{harmful} a \textsc{protective} anchor (4.6\%,
against 4.0\%). The anchor is a record-level consistency check on a reporting-selected subset; it
does not validate the claim-level event set of either annotation.

\begin{table}[!t]\centering\scriptsize
\setlength{\tabcolsep}{3pt}
\resizebox{\columnwidth}{!}{%
\begin{tabular}{@{}llllc@{}}
\toprule
Pipeline & Agreement [95\% CI] & Bal.\ acc.\ [95\% CI] & Macro-F1 & R$_\textsc{p}$ / R$_\textsc{h}$ / R$_\textsc{n}$ \\
\midrule
Primary & .807 [.782, .830] & .754 [.725, .785] & .766 & .923 / .836 / .504 \\
Second & .758 [.733, .782] & .724 [.695, .749] & .732 & .865 / .745 / .561 \\
Intersection, all 1{,}841 & .726 [.698, .753] & .682 [.650, .712] & .742 & .845 / .728 / .473 \\
Intersection, retained 1{,}587 & .842 [.821, .861] & .792 [.763, .820] & .805 & .950 / .860 / .566 \\
\bottomrule
\end{tabular}}
\caption{The numeric anchor against each annotation on the same 1{,}841 high-confidence records in
148 claims (claim-clustered 95\% intervals, 2{,}000 resamples). Paired primary minus second:
$+0.049$ \ci{0.030}{0.068} on agreement, $+0.074$ \ci{0.053}{0.094} on the 1{,}486 signed anchors,
$+0.059$ \ci{0.036}{0.081} on the 1{,}691 ratio-scale records. Among signed anchors: \textsc{null}
assigned by primary / second 4.2\% / 9.4\%, opposite sign 7.6\% / 7.9\%.
Intersection: the first row scores the 254 \textsc{unclear} records as disagreements, the second
scores the 1{,}587 records that retain a label, a selected subset on which both pipelines agree.}
\label{tab:anchor2}
\end{table}

\paragraph{Era-stratified comparison.} \Cref{tab:anchorera} bins the anchored records by era; the
61 records dated before 2000 form no bin. Agreement with the anchor rises across eras for both
pipelines. On the same records the pooled signed \textsc{harmful} share of the anchor, the primary
and the second pipeline all decline across the three eras, a descriptive pattern in this sample.
The within-claim slopes per decade are $-0.053$ \ci{-0.125}{0.013}, $-0.070$
\ci{-0.140}{-0.002} and $-0.043$ \ci{-0.112}{0.024}: all three point estimates are negative, but
only the primary pipeline's trend is established, the anchor's and the second pipeline's intervals
include zero, and the paired slope differences on the same claim resamples, anchor minus primary
$+0.017$ \ci{-0.019}{0.056} and anchor minus second $-0.010$ \ci{-0.067}{0.038} per decade, both
cover zero: the anchor's rate of drift is not distinguishable from either pipeline's, and the
comparison establishes neither a common within-claim drift process nor equal rates. The primary pipeline's over-\textsc{harmful} rate (records it
labels \textsc{harmful} whose anchor is \textsc{null} or \textsc{protective}) falls from 0.098 in the
2000s to 0.035 in the 2020s, so the annotator's bias relative to the anchor is era-dependent and
shrinks. The anchored subset is reporting-selected, and on it the primary slope is steeper than
corpus-wide ($-0.042$; \cref{app:controlled}).
\begin{table}[!t]\centering\scriptsize
\setlength{\tabcolsep}{4pt}
\resizebox{\columnwidth}{!}{%
\begin{tabular}{@{}lrccc@{}}
\toprule
Era & $n$ & Agreement primary / second & \multicolumn{2}{c}{\textsc{harmful} share anchor / primary / second} \\
\midrule
2000s & 394 & .751 / .678 & \multicolumn{2}{c}{.619 / .565 / .548} \\
2010s & 646 & .800 / .762 & \multicolumn{2}{c}{.502 / .463 / .448} \\
2020s & 740 & .841 / .796 & \multicolumn{2}{c}{.449 / .392 / .382} \\
\bottomrule
\end{tabular}}
\caption{Anchored records by era (61 records before 2000 form no bin). Within-claim slopes of
the \textsc{harmful} share per decade on these records: anchor $-0.053$ \ci{-0.125}{0.013},
primary $-0.070$ \ci{-0.140}{-0.002}, second $-0.043$ \ci{-0.112}{0.024}; paired differences,
anchor minus primary $+0.017$ \ci{-0.019}{0.056}, anchor minus second $-0.010$ \ci{-0.067}{0.038}.}
\label{tab:anchorera}
\end{table}

\paragraph{Scope.} The anchor covers a small, reporting-selected subset, cannot express
\textsc{unclear}, and measures consistency with reported effect estimates on that subset, not the
accuracy of the full label set.

\subsection{The full pre-cutoff representation}
\label{app:features}
The frozen feature builder computes 55 pre-cutoff columns per unit from the records published
before the cutoff (\cref{tab:feature-inventory}), organised in six blocks: direction (11), of
which the four trajectory variables are a subset; heterogeneity (6); decline (4); design (8);
volume (7); and 19 derived variables on numeric reporting and effect estimates. Three groups of
columns are numerically identical by construction (\texttt{n\_rr}, \texttt{n\_rr\_studies},
\texttt{n\_numeric}; \texttt{n\_pre\_t}, \texttt{n\_records}; \texttt{spread}, \texttt{val\_spread}),
so the 55 columns span 51 distinct-valued variables, and \texttt{decline\_sufficient} is constant at
zero. The representation is what every unit of the released benchmark carries, what every rule
comparator and single signal reads, and the input of the learned baselines marked ``all pre-cutoff
variables''. \compact reads
seven of the columns, marked $^{\ast}$ in the inventory; the set was fixed before any external
evaluation, contains no trajectory, heterogeneity, decline or design variable, and is tested rather
than assumed (\cref{app:ablation}). The inventory is reproduced at the end of the appendix
as a data dictionary.

\subsection{Specification chronology and freezing}
\label{app:chronology}
\cref{tab:chronology} records, for each fixed component, how it was selected, what data were
visible when it was fixed and when it was frozen. The seven \compact variables were written as a
single set operationalising static direction-and-effect aggregation and were never varied; they are
not the output of any variable search, and the logistic specification is the library default
(standardised inputs, $L_2$, $C=1$, class balancing), whose regularisation strength was deliberately
not tuned because a sweep scored on the same claims would be optimistic. Both were fixed while the
source benchmark, including its subsequent-stream target column, existed, and before any
\compact-versus-later-change result had been computed, so they are treated as a frozen source
specification rather than a prospectively preregistered feature set. The source method was then
frozen in a hash-signed configuration record before the External-1 and External-2 corpora were
built, and no external outcome was used to alter it. The two-variable score was defined after
\compact as a diagnostic of its inputs and acquires no retrospective preregistration status.

\begin{table*}[!t]\centering\scriptsize
\setlength{\tabcolsep}{3pt}
\begin{tabular}{@{}p{0.16\textwidth}p{0.26\textwidth}p{0.215\textwidth}p{0.11\textwidth}p{0.155\textwidth}@{}}
\toprule
Component & Selection basis & Visible when fixed & Frozen & Later analyses (sensitivity only) \\
\midrule
\compact variables ($p$, $a$, $n_{\text{dir}}$, $n_{\text{pre}}$, $\mu_{\log RR}$, $\sigma_{\log RR}$, $n_{RR}$) & Hand-composed as one set operationalising static direction-and-effect aggregation; never varied; no subset search & Source pre-cutoff variables and the subsequent-stream target column; no \compact later-change result; no external data & Hash-signed source-method record, before External-1/2 were built & Forward and backward selection, nested $L_1$ selection, learned models (App.~\ref{app:ablation}) \\
Logistic specification ($L_2$, $C=1$, class balancing) & Library default; $C$ not tuned (a same-claim sweep declined as optimistic) & As above: source pre-cutoff variables and the subsequent-stream target column; no \compact later-change result & Same record & Learned-model comparisons (App.~\ref{app:learned}) \\
\bottomrule
\end{tabular}
\caption{Specification chronology and freezing. ``Visible'' lists the data that existed when the
component was fixed; ``Frozen'' names the stage at which it was locked.}
\label{tab:chronology}
\end{table*}

\section{Second Annotation, Snapshot Versioning and Notation}
\label{app:release}
\subsection{Second annotation pipeline and cross-pipeline sensitivity}
\label{app:reextract}
Every claim-record association with a primary annotation and a publication year (112{,}453) was
re-annotated with Qwen2.5-7B-Instruct in bfloat16 using the same task framing, the same
3{,}000-character title-plus-abstract budget and the same 1{,}600-token truncation as the primary
extractor, with the \textsc{null} and \textsc{unclear} definitions made explicit
(\cref{fig:prompt2}); the primary prompt (\cref{fig:prompt}) is preserved verbatim, including its
historical wording, and the two executed configurations therefore differ in prompt wording, model
family and read-out. Instead of the first-token read-out, the probability of each label is the
teacher-forced probability of the complete label string normalised over the four labels; the
first-token read-out and the share of next-token mass on the four first tokens (mean 0.999) are
stored alongside, and the two read-outs agree on 99.2\% of associations and define the same 52
events. Agreement is reported with its denominator: all 112{,}453 associations, the associations
both pipelines resolve (non-\textsc{unclear}), and the 56{,}423 both call signed. Agreement measures how stable the target is across annotators, not its accuracy. Unit-level statistics rebuild the pre-cutoff inputs and the
endpoint of the 1{,}012 future-eligible units under the second pipeline, under label intersection
(associations on which the pipelines disagree become \textsc{unclear}) and under 500 soft-label
draws (one label per claim-record association per draw from the second pipeline's label
probabilities, reused across every cutoff of the claim; a record shared by two claims is drawn
independently per claim), and refit the two-variable score, the margin rule and the count-predicted
baseline leave-one-claim-out each time (\cref{tab:reextract}). Two comparisons are distinguished:
a score \emph{rebuilt and refitted} on an annotation's inputs and endpoint, and the \emph{primary}
score (primary inputs, trained on the primary endpoint) evaluated against another annotation's
endpoint. Soft-label ranges are variation across annotation draws with no claim-bootstrap layer.
The \textsc{null}-dominant agreement branch holds 42 units under the primary pipeline and 332
under the second; event retention by branch is given in the table. \Cref{tab:branchfree} shows
that the branch change in the agreement input contributes substantially to the reversed ordering
within the second-annotation setting; it does not show that changed signed membership or the
imputed ties play no role in the primary-to-second difference.

\begin{figure}[!t]
\small
\begin{tabular}{@{}p{0.95\linewidth}@{}}
\toprule
\textsc{system}: You classify the direction of association reported between an exposure and an
outcome. Answer with exactly one word: PROTECTIVE if the exposure is reported to reduce the outcome,
HARMFUL if it increases it, NULL if the abstract directly reports no association for the named
exposure and outcome, UNCLEAR if the pair is absent from the abstract or its direction cannot be
determined. \\[4pt]
\textsc{user}: Exposure: \{exposure\} \textbackslash n Outcome: \{outcome\} \textbackslash n
\textbackslash n Study: \{title\}. \{abstract\} \textbackslash n\textbackslash n Direction: \\
\bottomrule
\end{tabular}
\caption{The executed second-pipeline prompt (Qwen2.5-7B-Instruct, chat template, bfloat16). The
four complete label strings are scored by teacher forcing after \texttt{Direction:}; nothing is
generated.}
\label{fig:prompt2}
\end{figure}

\begin{table*}[!t]\centering\scriptsize
\setlength{\tabcolsep}{4pt}
\begin{tabular}{@{}lrrrr@{}}
\toprule
Primary (rows) $\backslash$ second (columns) & \textsc{protective} & \textsc{harmful} & \textsc{null} & \textsc{unclear} \\
\midrule
\textsc{protective} & 36{,}782 & 1{,}421 & 12{,}374 & 473 \\
\textsc{harmful} & 2{,}726 & 15{,}494 & 8{,}775 & 2{,}290 \\
\textsc{null} & 2{,}329 & 402 & 22{,}177 & 672 \\
\textsc{unclear} & 581 & 174 & 4{,}900 & 883 \\
\bottomrule
\end{tabular}
\hspace{4mm}
\begin{tabular}{@{}lrrrr@{}}
\toprule
Agreement branch of the unit & Units & Primary events & Second events & Both \\
\midrule
Primary: signed branch & 969 & 59 & 42 & 22 \\
Primary: \textsc{null}-dominant & 42 & 11 & 9 & 4 \\
Primary: fewer than 8 signed & 1 & 1 & 1 & 1 \\
Second: signed branch & 670 & 28 & 24 & 10 \\
Second: \textsc{null}-dominant & 332 & 39 & 18 & 13 \\
Second: fewer than 8 signed & 10 & 4 & 10 & 4 \\
\bottomrule
\end{tabular}
\vspace{4pt}

\begin{tabular}{@{}lrrr@{}}
\toprule
Endpoint cross-tabulation & Second: event & Second: non-event & Total \\
\midrule
Primary: event & 27 & 44 & 71 \\
Primary: non-event & 25 & 916 & 941 \\
Total & 52 & 960 & 1012 \\
\bottomrule
\end{tabular}
\hspace{4mm}
\begin{tabular}{@{}lrr@{}}
\toprule
Pipeline & Events on sparse histories ($<8$ signed) & Primary score on this endpoint, AUROC \\
\midrule
Primary & 1 & .948 \\
Second & 10 & .878 \\
Intersection & 14 & .869 \\
\bottomrule
\end{tabular}
\caption{Cross-pipeline sensitivity, detail. Claim-record association-level agreement with its denominators: 67.0\% of all 112{,}453 associations ($\kappa=0.52$), 72.7\% of associations both pipelines resolve ($\kappa=0.58$), 92.7\% of the 56{,}423 both call signed ($\kappa=0.83$); agreement is not accuracy. 21{,}149 primary signed annotations become \textsc{null} and 4{,}147 switch sign (\textsc{null} share 22.7\% against 42.9\%). Exact ties with at least eight signed records: 7 units under the primary pipeline (all events), 0 under the second. The last block reports the primary-trained score (primary inputs, trained on the primary endpoint) evaluated unchanged against each endpoint, as in \cref{tab:extractors}; the rebuilt-and-refitted scores (.857, .898) are those of \cref{tab:protocol}. Soft-label draws (500; variation across annotation draws, no claim-bootstrap layer): 43-57 events, signed score .779-.880, margin rule .938-.964. The Qwen first-token read-out defines the same 52 events as the full-string read-out.}
\label{tab:reextract}
\end{table*}


\paragraph{Common-support sensitivity.} \Cref{tab:support} restricts the full-cohort out-of-fold predictions to the 1002 units (160 claims) with at least eight signed pre-cutoff records under both pipelines; exact ties with adequate support are retained. Low event overlap persists (Jaccard 0.267) and the method ordering under the second pipeline is unchanged (signed minus margin $-0.119$ \ci{-0.215}{-0.045}; signed minus stationary $-0.131$ \ci{-0.236}{-0.052}). The reversal persists when sparse histories are excluded from the evaluation population, and it strengthens when the signed score is refitted leave-one-claim-out on the 1002 units: the signed score falls to 0.758 and signed minus margin to $-0.184$ \ci{-0.268}{-0.105}. Sparse histories therefore produce the reversal neither through evaluation nor through training.

\begin{table*}[!t]\centering\scriptsize
\setlength{\tabcolsep}{4pt}
\resizebox{\textwidth}{!}{%
\begin{tabular}{@{}llllllllll@{}}
\toprule
Support / training scope & Pipeline & Units / claims & Events / positive claims & Shared events & Jaccard & Margin AUROC / AUPRC & Stationary AUROC / AUPRC & Signed AUROC / AUPRC & Signed $-$ margin [95\% CI] \\
\midrule
\multicolumn{10}{@{}l}{\emph{Full archived cohort, out-of-fold scores}} \\
Full cohort & Primary & 1{,}012 / 160 & 71 / 23 & 71 & 1.000 & .881 / .468 & .909 / .497 & .948 / .535 & $+$.067 [.017, .120] \\
Full cohort & Second & 1{,}012 / 160 & 52 / 16 & 27 & .281 & .953 / .611 & .963 / .626 & .857 / .369 & $-$.096 [$-$.166, $-$.037] \\
\midrule
\multicolumn{10}{@{}l}{\emph{Common support ($\ge8$ signed pre-cutoff records under both pipelines), full-cohort-trained out-of-fold scores restricted}} \\
Common support, restricted & Primary & 1002 / 160 & 67 / 23 & 23 & .267 & .889 / .460 & .915 / .487 & .948 / .556 & $+$.059 [.016, .103] \\
Common support, restricted & Second & 1002 / 160 & 42 / 14 & 23 & .267 & .942 / .422 & .954 / .425 & .823 / .137 & $-$.119 [$-$.215, $-$.045] \\
\midrule
\multicolumn{10}{@{}l}{\emph{Common support, signed score refitted leave-one-claim-out on the 1002 units}} \\
Common support, refitted & Primary & 1002 / 160 & 67 / 23 & 23 & .267 & .889 / .460 & -- & -- & $+$.059 [.016, .104] \\
Common support, refitted & Second & 1002 / 160 & 42 / 14 & 23 & .267 & .942 / .422 & -- & .758 / -- & $-$.184 [$-$.268, $-$.105] \\
\bottomrule
\end{tabular}}
\caption{Common-support sensitivity. The support mask (hash in the released manifest) keeps units with at least eight signed pre-cutoff records under both pipelines (1002 units; 10 excluded, of which 10 sparse under the second pipeline and 1 under the primary pipeline; the primary sparse unit is among the second pipeline's sparse units). Each pipeline keeps its own operational endpoint; in the restricted rows, scores are the full-cohort leave-one-claim-out predictions restricted to the mask; in the refitted row, the signed score is refitted on the mask. All 2{,}000 claim-clustered resamples were evaluable for both pipeline comparisons.}
\label{tab:support}
\end{table*}

\paragraph{Within-class comparison.} \Cref{tab:withinclass} pools pairs of units that share a
pre-cutoff majority class ($p_t\neq0$) under each annotation, using the leave-one-claim-out
predictions of the development cell. Under the primary annotation the signed score's within-class gain
over the margin rule is not established, and the class-blind variant matches it; under the second
annotation and intersection labels the signed score falls below the margin rule within classes by an
established margin. The class-blind variant's within-class values (0.926, 0.929, 0.889) lie within 0.017 of the margin rule's (0.909, 0.929, 0.897) under every annotation, and no within-class gain of the class-blind variant is established. Under the
second annotation the signed score falls to 0.478 within \textsc{harmful}-majority units, an
established loss; this comparison includes \textsc{null}-dominant units, whose agreement input changes definition, and \cref{tab:branchfree} attributes the loss to that change.

\begin{table}[!t]\centering\footnotesize
\setlength{\tabcolsep}{3pt}
\resizebox{\columnwidth}{!}{%
\begin{tabular}{@{}lrrrl@{}}
\toprule
Annotation & Margin & Signed & Class-blind & Signed $-$ margin [95\% CI] \\
\midrule
Primary, pooled & .909 & .927 & .926 & $+$.018 [$-$.012, .060] \\
\quad \textsc{harmful}-majority (274, 42) & .833 & .863 & -- & $+$.030 [$-$.041, .114] \\
\quad \textsc{protective}-majority (730, 21) & .959 & .969 & -- & $+$.011 [$-$.001, .024] \\
Second, pooled & .929 & .631 & .929 & $-$.298 [$-$.538, $-$.147] \\
\quad \textsc{harmful}-majority & .906 & .478 & -- & $-$.428 [$-$.666, $-$.216] \\
\quad \textsc{protective}-majority & .955 & .810 & -- & $-$.145 [$-$.373, .011] \\
Intersection, pooled & .897 & .801 & .889 & $-$.095 [$-$.191, $-$.016] \\
\bottomrule
\end{tabular}}
\caption{Within-class AUROC, development cell, each annotation rebuilt and refitted: pooled over
pairs sharing a majority class, and within each class (units, events in parentheses). Intervals: claim-clustered bootstrap, 2{,}000 resamples.}
\label{tab:withinclass}
\end{table}

\paragraph{Branch-free agreement under every annotation.} \Cref{tab:branchfree} substitutes the
agreement definitions of \cref{app:agreement} upstream of the two-variable score under each
annotation and refits leave-one-claim-out. Under the shipped two-branch definition A the score
gains $+0.067$ over the margin rule on the primary annotation and loses $-0.096$ on the second;
under the branch-free definition C the primary gain is $+0.044$ \ci{-0.005}{0.091}, not
established, and the second-annotation difference is $-0.001$ \ci{-0.028}{0.018}; adding the
\textsc{null} share as a third input (C+null) gives $+0.051$ \ci{-0.006}{0.108} and $+0.006$
\ci{-0.022}{0.032}. Stratifying the shipped score by branch under the second annotation locates the
loss: on the 670 signed-branch units the difference is $-0.041$ with the full-cohort fit
restricted and $-0.004$ refitted within the branch, whereas on the 332 \textsc{null}-dominant units it
is $-0.231$ \ci{-0.481}{-0.017} restricted and $+0.047$ refitted.

\paragraph{The matched intervention on common support.} The headline intervention holds the cohort
and outcomes fixed and changes only the agreement input: on the 1{,}002 common-support units
(\cref{tab:support}), refitting the two-variable score with definition C instead of A raises its
second-annotation AUROC from 0.758 to 0.940, a paired improvement (C minus A) of $+0.182$
\ci{0.120}{0.248}; the margin rule reaches 0.942 on the same units, and C minus margin is
$-0.002$ \ci{-0.037}{0.023} on the same 2{,}000 claim resamples, whose lower bound misses the
non-inferiority margin $\delta=0.02$ by 0.017. In
full-cohort slices the improvement is $+0.246$ \ci{0.073}{0.405} on the \textsc{null}-dominant
branch and $+0.041$ on the signed branch; the difference between the two branch effects is not
tested, so ``largest on \textsc{null}-dominant units'' is descriptive. Across annotations 296
units move from the primary signed branch into the second pipeline's \textsc{null}-dominant branch
and none move in the reverse direction. The full-cohort slices and the common-support refit are
complementary analyses, not an exact decomposition of one another. The intervention establishes a
substantial representation-dependent contribution to the second-annotation loss within that
setting; it does not show that the branch is its only cause, it does not show that
annotation-induced signed membership plays no role in the primary-to-second difference, and it
does not change the 27 shared events of \cref{tab:extractors}. On the primary annotation the
branch-free score's development gain is not established ($+0.044$), but neither is its
difference from the shipped score ($-0.023$ \ci{-0.050}{0.002}; \cref{tab:astra2}), so the
result is that the branch-free model does not establish a development gain, not that removing the
branch eliminates the primary signal. The branch-free score's full-cohort AUROCs are 0.952
under the second annotation and 0.935 under intersection labels.

\begin{table*}[!t]\centering\footnotesize
\setlength{\tabcolsep}{4pt}
\begin{tabular}{@{}llll@{}}
\toprule
\multicolumn{4}{@{}l}{\emph{(a) Signed two-variable score $-$ margin rule by agreement definition (full cohort, refitted)}} \\
Annotation & Shipped A & Branch-free C & C $+$ \textsc{null} share \\
\midrule
Primary & $+$.067 [.017, .120] & $+$.044 [$-$.005, .091] & $+$.051 [$-$.006, .108] \\
Second & $-$.096 [$-$.166, $-$.037] & $-$.001 [$-$.028, .018] & $+$.006 [$-$.022, .032] \\
Intersection & $-$.030 [$-$.094, .023] & $+$.007 [$-$.016, .030] & $+$.013 [$-$.032, .051] \\
\midrule
\multicolumn{4}{@{}l}{\emph{(b) Shipped score $-$ margin rule by agreement branch, second annotation, full cohort}} \\
Branch & Units & Restricted & Refitted within branch \\
\midrule
Signed & 670 & $-$.041 & $-$.004 \\
\textsc{null}-dominant & 332 & $-$.231 [$-$.481, $-$.017] & $+$.047 \\
\midrule
\multicolumn{4}{@{}l}{\emph{(c) Intervention A $\to$ C, second annotation, outcomes fixed}} \\
Cohort & Units (events) & AUROC A $\to$ C & C $-$ A [95\% CI] \\
\midrule
Common support (margin .942; C $-$ margin $-$.002 [$-$.037, .023]) & 1{,}002 (42) & .758 $\to$ .940 & $+$.182 [.120, .248] \\
Full cohort, \textsc{null}-dominant branch & 332 & -- & $+$.246 [.073, .405] \\
Full cohort, signed branch & 670 & -- & $+$.041 \\
\bottomrule
\end{tabular}
\caption{The agreement-input intervention. (a) Paired differences from the margin rule under each
agreement definition and annotation (leave-one-claim-out, claim-clustered 95\% intervals). (b)
The shipped score stratified by agreement branch under the second annotation (``restricted'':
full-cohort out-of-fold predictions restricted to the branch; ``refitted'': refit within the
branch). (c) The matched intervention on the same cohort and outcomes, agreement input replaced and
score refitted; C $-$ margin on common support is $-0.002$ \ci{-0.037}{0.023}. Across annotations
296 units move from the primary signed branch into the second pipeline's \textsc{null}-dominant
branch and none move in the reverse direction.}
\label{tab:branchfree}
\end{table*}


\subsection{Controlled annotation and era effects}
\label{app:controlled}
To separate the factors that distinguish the two archived configurations, every association was
re-annotated in cells that change one factor at a time: the primary configuration re-run with
different batching, Llama-3.1-8B with the primary prompt and full-string read-out, Llama with the
second prompt, and Qwen2.5-7B with each prompt. \Cref{tab:controlled} compares neighbouring cells.
Changing the model moves the event set most (Jaccard 0.312 under either prompt) and shifts the
\textsc{null} share most ($+0.154$ and $+0.315$); changing the prompt moves it less (0.520 and
0.538); changing the read-out (0.899) stays close to the run-to-run level (0.846). The method ordering
follows the model family: the signed score falls below the margin rule in both Qwen cells ($-0.092$
and $-0.176$) and in neither Llama cell, consistent with the representation effect that
\cref{tab:branchfree} demonstrates for the archived second configuration and with the larger
\textsc{null}-dominant branch that the Qwen configurations produce. Across the archived annotations the fitted class coefficient
keeps its sign but shrinks: raw $\hat b_1$ is $+5.331$ (primary), $+2.138$ (second) and $+1.912$
(intersection), standardised $+3.010$, $+1.298$ and $+1.223$; the primary full-data fit has
$\hat b_0=+3.467$, $\hat b_1=+5.331$ and $\hat b_2=-12.193$.

\begin{table}[!t]\centering\scriptsize
\setlength{\tabcolsep}{2pt}
\resizebox{\columnwidth}{!}{%
\begin{tabular}{@{}llcccc@{}}
\toprule
Change & Between & $\kappa$ & Jaccard & \textsc{null} shift & Events \\
\midrule
Batching & primary re-run & .980 & .846 & .000 & 71$\rightarrow$73 \\
Read-out & Llama P1: token$\rightarrow$string & .986 & .899 & $+$.006 & 73$\rightarrow$77 \\
Prompt & Llama: P1$\rightarrow$P2 & .633 & .520 & $-$.117 & 77$\rightarrow$75 \\
Prompt & Qwen: P1$\rightarrow$P2 & .823 & .538 & $+$.044 & 49$\rightarrow$51 \\
Model & P1: Llama$\rightarrow$Qwen & .561 & .312 & $+$.154 & 77$\rightarrow$49 \\
Model & P2: Llama$\rightarrow$Qwen & .402 & .312 & $+$.315 & 75$\rightarrow$51 \\
\bottomrule
\end{tabular}}
\caption{Controlled annotation, one factor per row. P1, P2: the primary and second
prompts (\cref{fig:prompt,fig:prompt2}); cells other than the re-run use full-string read-out. The
signed score falls below the margin rule in both Qwen cells ($-0.092$, $-0.176$) and in neither
Llama cell.}
\label{tab:controlled}
\end{table}

\paragraph{Run-to-run variation.} The re-run keeps the checkpoint, prompt, tokenisation, bfloat16
inference and first-token read-out of the archived primary configuration and changes only the batch
size (the archived run used 16, with left padding). The label changes are consistent with bfloat16
numerical variation: a different batch size changes the padded tensor shapes and, with them, the
order of floating-point reductions, which moves the logits slightly and flips labels whose two highest
first-token logits nearly tie. At the label level the effect is small ($\kappa=0.980$), yet it moves
the event set (Jaccard 0.846; 71 against 73 events), because the thresholded endpoint turns small
label changes near the majority boundary into changes of the positive class. Released labels are fixed
by hash, so the benchmark itself does not vary; the variation concerns re-annotation.

\paragraph{Era effects.} Cross-pipeline agreement rises from $\kappa=0.411$ for records from the
1960s to 0.549 for the 2000s. Pooled over claims, the primary pipeline's \textsc{harmful} share falls
from 0.463 to 0.225 across decades while its \textsc{protective} share rises from 0.289 to 0.503;
these pooled shares mix claims. With claim fixed effects the drift is smaller but remains under the
primary annotation: the within-claim slope of the \textsc{harmful} share is $-0.042$
\ci{-0.061}{-0.023} per decade against a pooled slope of $-0.118$ \ci{-0.170}{-0.060}, so claim
composition accounts for about two thirds of the pooled shift; under the second annotation the
within-claim slope is $-0.014$ \ci{-0.034}{0.004} against $-0.101$ \ci{-0.171}{-0.022}. Across the
114 claims that have records in both the 1990s and the 2020s, the primary \textsc{harmful} share falls
from 0.401 to 0.351 ($-0.050$ \ci{-0.092}{-0.008}). Re-annotating a random sample
of 2{,}000 associations with the exposure and outcome names masked changes 25.1\% of primary labels: 35.3\%
for records before 2000 against 24.0\% after (difference $+0.114$ \ci{0.038}{0.190}). The probe
removes textual cues as well as named-claim knowledge, so it shows that label behaviour differs by
era without identifying why; on the anchored subset from 2000 onward the anchor's own
\textsc{harmful}-share slope has a similarly signed point estimate, with the stated uncertainty
(\cref{tab:anchorera}).

\subsection{Snapshot versioning and dataset card}
\label{app:snapshot}
The stored corpus (version 1) was assembled without recording the retrieval timestamp, the
per-query PubMed count or the query translation, so its identifier sets cannot be regenerated; it
is released as a \emph{non-regenerable derived snapshot} whose files are fixed by SHA-256 hashes in
the manifest. Post hoc provenance evidence nevertheless constrains the original harvest to a narrow
interval: the highest archived PubMed identifier (42{,}715{,}502) entered PubMed on 2026-09-09,
establishing a lower bound, and the earliest available write of the 165 archived record
files occurred on 2026-09-10 13:53~UTC, establishing an upper bound. The exact original retrieval
timestamp was not stored, so 2026-09-09 through 2026-09-10 13:53 UTC is the interval the surviving evidence supports and
not a recovered value; the bounds say when the harvest ran and do not make its identifier sets
regenerable. On 2026-09-22 (UTC) every stored query string was re-run through PubMed ESearch at
one timestamp, storing the count, the query translation, the complete identifier list and its hash
(version 2 identifier set; abstracts were not re-fetched, the benchmark was not rebuilt, and no performance figure derives from it).
\Cref{tab:snapshot} compares the versions at the association level (claim-identifier pairs) and
at the level of distinct identifiers: 103{,}175 of 112{,}453 archived associations (91.7\%) and
88{,}649 of 96{,}091 distinct archived identifiers (92.3\%) are returned again.
Fifty-four claims have per-claim Jaccard overlap below 0.8, through two disjoint mechanisms. Nineteen
stored queries return no records at re-retrieval; they are exactly the nineteen claims that appear
to lose more than 10\% of their archived identifiers. Why they return nothing is not recoverable from
the surviving metadata: an edited query string, a changed automatic term translation or a query that
never ran as stored are all consistent with the evidence; the release flags these 19 claims
(manifest field \texttt{query\_unconfirmed}, true for the 19 and false for the other 146), whose
stored query is not confirmed to be the harvesting query, and states the discrepancy without
assigning a mechanism.
The other 35 return more records than were archived and lose none; their low overlap reflects the per-query cap described below (the archive
holds 56 to 4{,}888 records per claim). Neither mechanism is
change in the literature. Removing all 54 claims changes no conclusion in any of the twelve protocol
cells.

\paragraph{Retrieval cap.} The archive metadata records, for each query, whether retrieval stopped
at the per-query cap: 48 queries did, and their archives hold 780-800 records each. Forty of the 48
are among the 54 low-overlap claims above (the 35 whose stored queries now return more records than
were archived, and 5 of the 19 whose stored queries now return none), and 8 overlap with today's
retrieval at 0.8 or more.
The cap removed early history. For 43 capped claims with a re-retrieved set to compare against, the
archive begins later than the earliest record available today in 42 (median gap 45 years), covers
none of the oldest quartile of today's records for the median claim, and covers less than half of
today's records for 36 claims; the beta-carotene-cancer claim, for example, holds 797 archived records from
2012-2026 against a literature that begins in 1977. Capped claims accordingly have later cutoffs
(median 2021 against 2015). Removing all 48 capped claims leaves 715 units. The signed score's gain over the margin rule
under the primary annotation remains established, $+0.070$ \ci{0.009}{0.131} against $+0.067$
\ci{0.017}{0.120} on the full benchmark, and so does its loss under the second annotation, $-0.128$
\ci{-0.203}{-0.049} against $-0.096$ \ci{-0.166}{-0.037}. The streams most plausibly affected by
the cap therefore do not produce either central comparison. Rebuild levels: released
predictions and outcomes support metric and interval reproduction; association-level labels,
identifiers and years support the directional rebuilds listed in the manifest, including counts,
operational direction, agreement, endpoints and stationary probabilities, using the documented
count-model outputs or fitting procedure; recreating text-derived numeric and design features and
rerunning annotation requires the original abstract text and the documented pipeline
configurations, and these stages are not reproducible from the released package alone. The dated
query-identifier audit does not recreate the original text snapshot. The 145-claim curated list is a fixed artefact whose
assembly procedure is not recoverable; the 20 later additions are documented individually. The
external corpora were shipped as unit-level feature files without query strings and cannot be
re-retrieved; their selection (disjoint exposure domains chosen before any transfer result) and
freeze (feature builder and extractor revision identical to the source) are recorded in the
dataset card.

\begin{table}[!t]\centering\scriptsize
\setlength{\tabcolsep}{4pt}
\resizebox{\columnwidth}{!}{%
\begin{tabular}{@{}lr@{}}
\toprule
Quantity & Value \\
\midrule
Archived snapshot (v1): claim queries & 165 \\
Archived claim-record associations & 112{,}453 \\
Archived records dated 2026 (partial year) & 5{,}818 \\
Original snapshot retrieval window & 2026-09-09 to 2026-09-10 13:53 UTC \\
Exact original retrieval timestamp & not recorded \\
Re-retrieval timestamp (UTC) & 2026-09-22 06:01:37 \\
Associations returned at re-retrieval (v2) & 153{,}428 \\
Archived associations returned again & 103{,}175 (91.7\%) \\
Distinct archived identifiers / returned again & 96{,}091 / 88{,}649 (92.3\%) \\
Distinct identifiers at re-retrieval & 132{,}629 \\
Archived associations no longer returned & 9{,}278 \\
Associations new at re-retrieval & 50{,}253 \\
Per-claim Jaccard, median [min, max] & 0.992 [0.000, 1.000] \\
Claims losing $>10$\% of stored identifiers & 19 \\
External corpora re-retrievable & no (unit-level files only) \\
\bottomrule
\end{tabular}}
\caption{Snapshot versioning. Version 1 is the stored, non-regenerable derived snapshot on which every experiment runs; version 2 is the timestamped ESearch re-retrieval of the same query strings (identifier lists and hashes only). Per-query counts, translations and identifier hashes are in the released manifest.}
\label{tab:snapshot}
\end{table}


\subsection{Notation and terminology}
\label{app:notation}
\Cref{tab:notation} lists the notation and terminology used throughout the paper.
\begin{table}[!t]\centering\scriptsize
\setlength{\tabcolsep}{3pt}
\begin{tabular}{@{}p{0.27\columnwidth}p{0.7\columnwidth}@{}}
\toprule
Term or symbol & Meaning \\
\midrule
PubMed record & One title-abstract-year-identifier tuple returned by a fixed query; called a record in manuscript prose. The historical executed extraction prompt uses the word ``Study'' verbatim. \\
automatic annotation & The extractor-assigned direction of a claim–record association (\textsc{protective} $-1$, \textsc{null} $0$, \textsc{harmful} $+1$, \textsc{unclear}). \\
claim-cutoff unit $(c,t)$ & A claim and a calendar year; inputs from records with $\tau_i<t$. \\
future-eligible historical unit & A unit with at least 20 later records; the main evaluation population (1{,}012 source units). \\
$n^{+}_t,n^{-}_t,n^{0}_t,n_t$ & Pre-cutoff counts of \textsc{harmful}, \textsc{protective}, \textsc{null} and signed records. \\
$h$, $l$, $n=h+l$ & Actual \textsc{harmful}, \textsc{protective} and signed pre-cutoff counts; never replaced by zeros on sparse histories. \\
$r_t$ & Raw signed mean $(h-l)/n$, defined for $n>0$. \\
$p_t$, $|p_t|$ & Archived operational feature: $r_t$ when $n\ge8$, and 0 otherwise; $|p_t|$ its absolute margin. \\
$s_t$ & Operational current sign $\sgn(p_t)$; sparse evidence and a supported exact tie both give 0 but are identified separately. \\
$R_{\mathrm{raw}}$ / $\pi^{\mathrm{raw}}_t(m)$ & Theoretical raw-majority event (future raw majority differs from $\sgn(h-l)$) and its probability, \cref{eq:bb}; a labelled diagnostic. \\
$R$ / $\pi^{\mathrm{op}}_t(m)$ & Archived operational event (future sign differs from $s_t$; no future gate) and its probability with the same posterior; the scored stationary reference. \\
uncalibrated / recalibrated & $\pi^{\mathrm{op}}_t$ without Platt scaling (still operational) versus the fitted logistic recalibration of $\logit\pi^{\mathrm{op}}_t$ in \cref{tab:matched}. \\
$a_t$ & Agreement: majority share among signed records, or the \textsc{null} share when \textsc{null} dominates (definition A); the branch-free variant uses the signed-only share on every unit (definition C). \\
$p^{+}_t$, $m$ & Pooled direction and signed count of the subsequent stream. \\
subsequent-stream sign change $R$ & Primary endpoint, \cref{eq:target}. \\
cumulative updated-direction change $R_{\mathrm{cum}}$ & Sign change of the pooled direction over all records after updating. \\
margin rule & The untrained score $1-|p_t|$; the reference of every $\Delta$. \\
stationary baseline & Count-predicted operational Beta-Binomial probability $\pi^{\mathrm{op}}_t$ (uncalibrated). \\
order features $z_t$ & Seven sign-invariant features of the pre-cutoff label order (\cref{app:order}). \\
retrospective claim-held-out backtest & Leave-one-claim-out fitting with training outcomes to the end of the corpus. \\
origin-defined calendar cohort & Cohort defined at origin $Y$ from pre-origin eligibility; accrual reported as censoring; claim-disjoint training. \\
$\Delta$ & Each table states its reference and direction: ranking differences are model minus reference (margin rule unless stated); proper-score improvement columns use reference loss minus model loss, so positive values indicate improvement; all intervals use the stated paired resampling. \\
realised-count / count-predicted & $\pi^{\mathrm{op}}_t$ with the observed future signed count (diagnostic) or averaged over a training-fitted count model (deployable); $\pi^{\mathrm{raw}}_t$ is reported only as a diagnostic. \\
rebuilt vs.\ primary score & A score refitted on an annotation's own inputs and endpoint, or the primary score evaluated against another annotation's endpoint. \\
interval types & Claim-bootstrap CI; wild cluster bootstrap-$t$ CI; null range over 1{,}000 control draws; flip range over 200 recodings; annotation-draw range over 500 soft-label draws. \\
\bottomrule
\end{tabular}
\caption{Notation and terminology used throughout the paper.}
\label{tab:notation}
\end{table}

\section{Stationary Baseline, Order Features and Sign Invariance}
\label{app:bb}
\subsection{Derivation and checks}
\paragraph{Setting.} Fix a claim and a cutoff with $h=n^{+}_t$ \textsc{harmful} and $l=n^{-}_t$
\textsc{protective} signed pre-cutoff annotations, $n=h+l>0$, and $m\ge1$ signed records in the
future window. Under the model of \cref{sec:bb} the posterior is
$\theta\mid h,l\sim\mathrm{Beta}(h+\alpha,\,l+\alpha)$ and the number $K$ of future \textsc{harmful}
labels is Beta-Binomial, $\BB(k;m,a,b)=\binom{m}{k}\mathrm{B}(k+a,\,m-k+b)/\mathrm{B}(a,b)$ with
$a=h+\alpha$, $b=l+\alpha$. The future pooled direction is $p^{+}_t=(2K-m)/m$. For the theoretical raw-majority event,
$R_{\mathrm{raw}}=1$ iff $\sgn(2K-m)\neq\sgn(h-l)$, which gives \cref{eq:bb}. The archived
operational event instead compares the future sign with $s_t=\sgn(p_t)$; its probability is
$\pi^{\mathrm{op}}_t(m)=\sum_k\mathbf{1}[\sgn(2k-m)\ne s_t]\,\BB(k;m,h{+}\alpha,l{+}\alpha)$, with the same
posterior, and is the scored stationary reference throughout. The formula is symmetric under
$h\leftrightarrow l$ (checked numerically to $10^{-15}$), finite and in $[0,1]$ on the tested
domain.

\paragraph{Full statement.} Write $\pi^{\mathrm{raw}}_t(m)$ for $\pi_t(m)$ of \cref{lem:stationary}.
(a) For fixed $n$, $m$ and $\alpha$, $\pi^{\mathrm{raw}}_t(m)$ is non-increasing in the raw absolute
margin $|h-l|/n$ among histories with a non-zero raw majority; (b) as $m\to\infty$,
$\pi^{\mathrm{raw}}_t(m)\to I_{1/2}(\max(h,l)+\alpha,\,\min(h,l)+\alpha)$ for $h\ne l$ and $\to1$ for
$h=l$, where $I_x(a,b)$ is the regularised incomplete Beta function; for $h=l$ and finite $m$ the
probability is 1 for odd $m$ and $1-P(K=m/2\mid h,l)$ for even $m$; (c) when $n$, $m$ and $\alpha$
are common across units, the margin rule is a ranking surrogate for $\pi^{\mathrm{raw}}_t$ on the
non-tied subset. \Cref{lem:stationary} states (a) and (c).

\begin{proof}
(a) Let $n$, $m$ and $\alpha$ be fixed and $h>l$ (the case $h<l$ is symmetric). Moving one record
from $l$ to $h$ replaces the posterior $\mathrm{Beta}(a,b)$ by $\mathrm{Beta}(a+1,b-1)$, which is
larger in the likelihood-ratio order and therefore stochastically larger. A Beta-Binomial
variable is a mixture of $\mathrm{Binomial}(m,\theta)$ laws, each stochastically increasing in
$\theta$, so $K$ is stochastically larger under the larger posterior and the lower tail
$P(2K-m\le0)$, the change event for a \textsc{harmful} raw majority, is non-increasing in $h$,
that is in $|h-l|/n$. The argument uses a non-zero raw majority: the change event of a tied history
is $\{2K\neq m\}$, and the ordering does not extend across the tie
($\pi^{\mathrm{raw}}(30,30;2)=32/63\approx0.508$ but $\pi^{\mathrm{raw}}(31,29;2)=475/651\approx0.730$
with $\alpha=1$).
(b) For $h>l$ the change event is $\{2K-m\le0\}$. Conditionally on $\theta$, $K/m\to\theta$ almost
surely as $m\to\infty$; the posterior is continuous, so $P(\theta=1/2\mid h,l)=0$, and hence
$\pi^{\mathrm{raw}}_t(m)\to P(\theta<1/2\mid h,l)=I_{1/2}(h+\alpha,l+\alpha)$. For $h<l$ the change
event is $\{2K-m\ge0\}$, giving $\pi^{\mathrm{raw}}_t(m)\to P(\theta>1/2\mid h,l)=1-I_{1/2}(h+\alpha,l+\alpha)=I_{1/2}(l+\alpha,h+\alpha)$.
The two cases give $I_{1/2}(\max(h,l)+\alpha,\min(h,l)+\alpha)$. For $h=l$, change occurs whenever
the future sample is not tied; it therefore has probability 1 for odd $m$ and $1-P(K=m/2\mid h,l)$
for even $m$, and because the posterior has no atom at $1/2$ the probability of a future exact tie
vanishes and the change probability tends to 1.
(c) If $n$, $m$ and $\alpha$ are common across units, $\pi^{\mathrm{raw}}_t$ is a function of $h-l$
alone, non-increasing in $|h-l|/n$ by (a) on the non-tied subset and symmetric under
$h\leftrightarrow l$; ranking by $\pi^{\mathrm{raw}}_t$ on that subset is therefore the ranking by
$1-|h-l|/n$, with ties within equal margins. The ranking equivalence is restricted to non-tied
raw-majority histories; no general ordering of tied histories relative to non-tied histories is
asserted.
\end{proof}

\paragraph{Executed endpoint contract, ties and sparse evidence.} The endpoint uses $\sgn(0)=0$.
The frozen builder returns $p_t=r_t=(h-l)/(h+l)$ when $h+l\ge8$ and $p_t=0$ otherwise; the
operational current sign is $s_t=\sgn(p_t)$, so $s_t=0$ for an exact tie with at least eight
signed records and for a sparse history. The future side uses the raw signed mean of any
non-empty signed sample and 0 when no signed record follows; no future gate is applied, and no
future-eligible source unit has an empty signed future sample. The archived event equals the
reproduced operational event on all 1{,}012 units (0 mismatches). The scored
baselines use the operational indicator $\pi^{\mathrm{op}}_t$; the raw-majority indicator
$\pi^{\mathrm{raw}}_t$ differs on sparse histories (one unit under the primary pipeline, ten under
the second, fourteen under intersection labels), where $s_t=0$ makes any non-zero future sign a change while
the posterior still uses the actual counts (for $h=7$, $l=0$ the posterior is $\mathrm{Beta}(8,1)$
and the operational event has probability 1 for odd $m$). Seven source units under the primary
pipeline have an exact tie with at least eight signed records; all seven are events. The
raw-majority scores are retained as labelled diagnostic columns in the released per-unit files;
aligning the indicator changed the count-predicted score on 1 source units
(maximum absolute change 0.959). In the count-predicted implementation the drawn future
count is floored at one, so the $m=0$ branch never occurs; in the Dirichlet-multinomial draws a
future window without signed records is a change iff $s_t\ne0$, the endpoint's rule.

\paragraph{Count-predicted implementation.} The realised $m$ is unavailable at the cutoff. On the
training claims of each fold we fit $\log(1+m)=b_0+b_1\log(1+n)+b_2\log(1+\text{accrual}_t)+\varepsilon$
by least squares, where accrual is the number of signed records per year over the five years before
$t$ and $\varepsilon$ is Gaussian with the residual standard deviation; 400 draws of $m$ are
rounded to integers and floored at one; for each distinct drawn count the implementation evaluates
the operational probability $\pi^{\mathrm{op}}_t(m)$ with the current sign $s_t$ and averages these
values with the draw multiplicities (no Monte Carlo sampling of labels). The posterior retains the
actual pre-cutoff signed counts. The raw-majority probabilities of \cref{eq:bb} are retained
separately as diagnostic outputs.
The held-out claim is excluded from the count-model fit. An earlier reference draws the ratio $m/n$
from the training claims instead (\emph{ratio} row of \cref{tab:stationary}). The
Dirichlet-multinomial version places $\mathrm{Dir}(h+1,l+1,n^{0}_t+1)$ on the three resolved
labels, draws the future total count from the same count model fitted to signed-plus-\textsc{null}
counts, and compares the sign of the future signed majority with the current sign; a future window
with no signed record counts as a change when the current sign is non-zero.

\paragraph{Numerical checks.} The implementation passes the boundary tests of the corrected
proposition (exact $32/63$ and $475/651$ at $\alpha=1$; symmetry under $h\leftrightarrow l$;
probability 1 for tied histories and odd $m$; $\pi^{\mathrm{raw}}(8,0;m)\to1/512$ as $m$ grows;
non-increasing in the raw margin on non-tied histories; probability mass sums to one). Only finite
sums enter any prediction; the limiting expression is not used. A posterior-predictive Monte Carlo
check draws $\theta^{(b)}\sim\mathrm{Beta}(h+1,l+1)$ and $K^{(b)}\sim\mathrm{Binomial}(m,\theta^{(b)})$,
$b=1,\ldots,2{,}000$, for every source unit with its realised $m$ and the operational current sign:
the mean absolute difference between the empirical change frequency and $\pi^{\mathrm{op}}_t$ is
0.0012 (maximum 0.025, mean Monte Carlo standard error 0.0015; 98.8\%
of units within two standard errors). A separate plug-in simulation
fixes $\theta=h/n$ and draws 1{,}000 future samples of the realised size for the operational event; it
is a plug-in Binomial comparison, not a posterior draw, and \cref{fig:analytic}A reports it as
such: the plug-in Binomial probability matches its own simulation at $r=0.999$, and the
flat-prior operational probability correlates with the same simulation at $r=0.949$ with
a higher expected event total (39.9 against 26.6)
because the prior places mass toward balance. Expected event totals are run-specific: the
operational realised-count probability sums to 39.9 on the 1{,}012 units, the
raw-majority probability of \cref{eq:bb} to 38.9 (the sparse unit
contributes the difference), and the raw-majority plug-in to 25.6.
The within-claim permutation of this check (1{,}000 draws that re-draw the pre-cutoff labels as
well, operational indicator) gives 35.4 expected events and a per-unit
correlation of $r=0.44$ with the count-conditional operational probability; the coherent
Null A run of \cref{app:null} (1{,}000 draws, \cref{tab:astra7}) has mean 35.3 events and is a
separate run, not the same quantity. Non-tied monotonicity holds in all 624 tested $(n,m)$ configurations.

\begin{figure}[!t]
\centering
\includegraphics[width=\linewidth]{figures/fig4_analytic_vs_simulation.pdf}
\caption{Comparison with a plug-in Binomial simulation. A: per-unit analytic change probability
against the change frequency in 1{,}000 draws with $\theta$ fixed at the pre-cutoff proportion and
the realised future count; the plug-in Binomial matches this simulation and the flat-prior
Beta-Binomial adds prior mass toward balance. B: \cref{eq:bb} as a function of the pre-cutoff
margin for several pre-cutoff and future sample sizes on non-tied histories.}
\label{fig:analytic}
\end{figure}

\begin{table*}[!t]\centering\scriptsize
\setlength{\tabcolsep}{3pt}
\resizebox{\textwidth}{!}{%
\begin{tabular}{@{}llllllllr@{}}
\toprule
Score & AUROC [95\% CI] & AUPRC [95\% CI] & Log loss & Brier & Cal.\ slope & ECE & $\Delta$AUROC vs.\ margin & $\Delta$AUROC vs.\ two-variable \\
\midrule
Beta-Binomial, realised future count (diagnostic) & .882 [.816, .942] & .490 [.374, .619] & .272 & .0491 & .31 & .033 & $+$.000 [$-$.033, .034] & $-$.066 [$-$.125, $-$.013] \\
Stationary Beta-Binomial, count-predicted (operational) & .909 [.867, .949] & .497 [.379, .618] & .200 & .0499 & .53 & .033 & $+$.028 [.008, .050] & $-$.039 [$-$.080, $-$.004] \\
Beta-Binomial, training ratio of counts & .877 [.804, .938] & .467 [.337, .600] & .194 & .0498 & .64 & .025 & $-$.004 [$-$.055, .044] & $-$.071 [$-$.133, $-$.016] \\
Dirichlet-multinomial, realised count (diagnostic) & .849 [.768, .929] & .483 [.368, .605] & .297 & .0489 & .28 & .031 & $-$.032 [$-$.080, .009] & $-$.099 [$-$.176, $-$.025] \\
Dirichlet-multinomial, count-predicted & .893 [.840, .943] & .492 [.379, .615] & .223 & .0497 & .40 & .030 & $+$.011 [$-$.011, .033] & $-$.055 [$-$.100, $-$.011] \\
Margin rule $1-|p_t|$, untrained & .881 [.824, .936] & .468 [.359, .589] & .184 & .0515 & .92 & .012 & reference & -- \\
Two-variable logistic ($p$, $a$) & .948 [.921, .971] & .535 [.391, .676] & .146 & .0452 & .94 & .027 & $+$.067 [.017, .120] & reference \\
\compact logistic & .946 [.918, .969] & .527 [.378, .669] & .149 & .0460 & .94 & .014 & $+$.064 [.011, .122] & $-$.002 [$-$.011, .006] \\
\bottomrule
\end{tabular}}
\caption{Stationary posterior-predictive references for the archived operational endpoint on the source cohort (1{,}012 units, 71 events); all rows evaluate the archived operational endpoint $\pi^{\mathrm{op}}_t$. The stationary posterior-predictive rows use uncalibrated probabilities; probability metrics for the margin, two-variable and \compact scores use training-only Platt calibration (\cref{tab:core,app:metrics}). Realised-count rows use observed future counts and are diagnostic rather than origin-available forecasts; count-predicted rows integrate over the training-fitted count model of \cref{app:bb}; the ratio row draws the future/past count ratio from the training claims. The raw-majority probability of \cref{eq:bb} is a separately labelled diagnostic in the released per-unit files (columns \texttt{*\_rawsign}), not a row of this table. A prevalence-only forecast has log loss .254 and Brier .0652. The calibrated two-variable score improves on the deployable baseline by .0047 in Brier [$-$.004, .014] and .054 in log loss [$-$.001, .123] (paired, positive = better).}
\label{tab:stationary}
\end{table*}


\subsection{Order features and the offset diagnostic}
\label{app:order}
Let $z_1,\ldots,z_n$ be the signed pre-cutoff labels in publication order (records with equal
years keep the archived file order), multiplied by the sign of the current majority so that $+1$
agrees with the majority and $-1$ contradicts it, and let $h=\lfloor n/2\rfloor$. The seven
features are: (1) \emph{recent minus historical dissent}, the share of $-1$ among records dated
within six years before $t$ minus the share among earlier records; (2) \emph{recency-weighted
margin}, $\sum_i w_i z_i/\sum_i w_i$ with $w_i=\exp(-(t-1-\tau_i)/5)$; (3)
\emph{cumulative-margin slope}, the running mean of $z$ at $n$ minus its value at $h$; (4)
\emph{opposing run length}, the number of consecutive $-1$ labels at the end of the sequence; (5)
\emph{contraction}, $|\bar z_{1:h}|-|\bar z_{h+1:n}|$; (6) \emph{CUSUM range},
$(\max_j S_j-\min_j S_j)/n$ with $S_j=\sum_{i\le j}(z_i-\bar z)$; (7) \emph{acceleration of
contradiction}, the dissent share in the last quarter of the sequence minus the share in the
preceding quarter. Units with fewer than four signed records receive zeros. All seven are
unchanged when every signed label of the claim is exchanged. The offset model is fitted by
unweighted penalised likelihood with the offset fixed; the ridge strength $C\in\{0.1,1,10\}$ is
chosen inside each training fold by grouped five-fold log loss over the training claims. Two
reference fits use the same offset with composition variables ($|p_t|$, $a_t$, $\log(1+n)$) and
with composition plus order features (\cref{tab:order}).

\paragraph{Matched recalibration comparison.} Because the offset model's intercept recalibrates the
uncalibrated operational stationary probability, its proper-score improvement over that offset
does not isolate order information. \Cref{tab:matched} therefore compares, with identical estimator (unweighted
$L_2$ logistic, $C=1$), identical leave-one-claim-out folds and identical bootstrap resamples, a
logistic regression on $\logit\pi^{\mathrm{op}}_t$ alone (slope and intercept recalibration) and
on $\logit\pi^{\mathrm{op}}_t$ plus the seven order features. This is a different estimand from the
fixed-offset model of \cref{tab:order}, whose offset slope is fixed at one.

\begin{table}[!t]\centering\scriptsize
\setlength{\tabcolsep}{4pt}
\begin{tabular}{@{}lrrr@{}}
\toprule
Metric & Recalibrated OS & $+$ order features & Improvement [95\% CI] \\
\midrule
AUROC & .899 & .916 & $+$.017 [$-$.012, .045] \\
AUPRC & .454 & .520 & $+$.066 [$-$.065, .195] \\
Log loss & .167 & .158 & $+$.009 [$-$.010, .032] \\
Brier score & .0487 & .0458 & $+$.003 [$-$.004, .011] \\
\bottomrule
\end{tabular}
\caption{Matched recalibration comparison on the source cohort. The reference is a fitted logistic recalibration of the uncalibrated operational stationary logit $\logit\pi^{\mathrm{op}}_t$ (a different procedure from the uncalibrated baseline of \cref{tab:stationary} and from the fixed-offset model of \cref{tab:order}); the augmented model additionally uses the seven order features. Both use unweighted $L_2$ logistic regression with $C=1$, the same leave-one-claim-out folds and the same claim-bootstrap resamples. Positive values favour the augmented model: augmented minus reference for AUROC and AUPRC, reference minus augmented for log loss and Brier score, computed from unrounded paired predictions. None of the paired intervals excludes zero. \textbf{OS: operational stationary}}
\label{tab:matched}
\end{table}



\begin{table*}[!t]\centering\scriptsize
\setlength{\tabcolsep}{4pt}
\resizebox{\textwidth}{!}{%
\begin{tabular}{@{}lllllll@{}}
\toprule
Model & AUROC [95\% CI] & AUPRC & Log loss & $\Delta$AUROC vs.\ offset & $\Delta$AUPRC vs.\ offset & $\Delta$log loss vs.\ offset (positive = better) \\
\midrule
Stationary offset only ($\pi_t$) & .909 [.867, .949] & .497 & .200 & reference & reference & .000 [.000, .000] \\
Offset $+$ order features & .915 [.877, .952] & .514 & .167 & $+$.006 [$-$.016, .028] & $+$.017 [$-$.063, .097] & .034 [$-$.009, .083] \\
Offset $+$ composition & .911 [.871, .946] & .442 & .167 & $+$.002 [$-$.015, .021] & $-$.055 [$-$.108, $-$.007] & .033 [$-$.006, .078] \\
Offset $+$ composition $+$ order & .914 [.877, .949] & .508 & .166 & $+$.005 [$-$.019, .031] & $+$.011 [$-$.098, .117] & .035 [$-$.013, .089] \\
Order features only, no offset & .898 [.851, .942] & .468 & .169 & $-$.011 [$-$.036, .015] & $-$.028 [$-$.185, .132] & .031 [$-$.021, .092] \\
Two-variable logistic & .948 [.921, .971] & .535 & .146 & $+$.039 [.004, .080] & $+$.038 [$-$.085, .155] & -- \\
Margin rule & .881 [.824, .936] & .468 & .184 & $-$.028 [$-$.050, $-$.008] & $-$.029 [$-$.057, $-$.002] & -- \\
\bottomrule
\end{tabular}}
\caption{Fixed-offset diagnostic: incremental value of temporal order beyond the uncalibrated operational stationary probability (source cohort, leave-one-claim-out, nested ridge selection). The offset is the uncalibrated operational stationary probability. Standardised coefficients of the full-data fit: dissent recent vs hist -0.01, recency weighted margin +0.97, cum margin slope -0.29, opp run len +0.01, contraction +0.79, cusum max +0.11, contra accel +0.37.}
\label{tab:order}
\end{table*}


\subsection{Sign-flip experiments}
\label{app:signinv}
Direction-coding sensitivity on the source cohort; the invariant model and the augmented models have not been evaluated under transfer.
\begin{table*}[!t]\centering\scriptsize
\setlength{\tabcolsep}{4pt}
\resizebox{\textwidth}{!}{%
\begin{tabular}{@{}llllll@{}}
\toprule
Model & AUROC [95\% CI] & AUPRC & $\Delta$AUROC vs.\ margin & $\Delta$AUROC vs.\ signed two-variable & AUROC under claim-wise flips, mean [min, max] \\
\midrule
Margin rule & .881 [.824, .936] & .468 & reference & $-$.067 [$-$.120, $-$.017] & -- \\
Two-variable, signed & .948 [.921, .971] & .535 & $+$.067 [.017, .120] & reference & .896 [.886, .917] \\
\compact, signed & .946 [.918, .969] & .527 & $+$.064 [.011, .122] & $-$.002 [$-$.011, .006] & .897 [.885, .921] \\
Sign-invariant ($|p_t|$, $a_t$) & .897 [.839, .945] & .445 & $+$.016 [$-$.017, .057] & $-$.051 [$-$.092, $-$.013] & unchanged \\
Sign-invariant, full & .913 [.872, .951] & .466 & $+$.032 [$-$.022, .091] & $-$.035 [$-$.068, $-$.005] & unchanged \\
Two-variable, flip-augmented & .901 [.844, .947] & .436 & -- & $-$.047 [$-$.089, $-$.011] & .901 (full flip) \\
\compact, flip-augmented & .900 [.846, .947] & .441 & -- & $-$.048 [$-$.090, $-$.012] & .900 (full flip) \\
\bottomrule
\end{tabular}}
\caption{Sign-invariance experiments on the source cohort. Claim-wise flips exchange the signed labels of exactly half of the claims (all cutoffs of a claim together) in each of 200 draws and refit the signed models leave-one-claim-out; a full exchange of every claim leaves the signed scores unchanged because the fitted sign of $\beta$ reverses. Direction strata of the signed two-variable score: \textsc{protective}-majority units (730 units, 21 events) AUROC .969 (margin .959, sign-invariant .925); \textsc{harmful}-majority units (274, 42) .863 (margin .833, sign-invariant .905).}
\label{tab:signflip}
\end{table*}

\subsection{Exact pair decomposition of the AUROC difference}
\label{app:decomp}
The AUROC of a score is the share of positive-negative unit pairs it orders correctly (ties counted
one half), so a paired AUROC difference between two scores is a mean over the same pairs of the
difference in their pairwise ordering indicators, and it partitions exactly by pair type. We
partition the pairs of the development cell by whether the positive and the negative unit share a
pre-cutoff majority class (within-class), have different classes (across-class), or involve at
least one zero-direction unit ($p_t=0$). In the development cell the 71 positives (42
\textsc{harmful}-majority, 21 \textsc{protective}-majority, 8 zero-direction) and 941 negatives
(232 and 709; no zero-direction unit is a negative) form $71\times941=66{,}811$ pairs:
$42\times709+21\times232=34{,}650$ across-class, $42\times232+21\times709=24{,}633$ within-class and
$8\times941=7{,}528$ zero-direction, and each component's contribution is the sum of its pairs'
indicator differences divided by 66{,}811. \Cref{tab:decomp} reports the signed-minus-margin
difference by component under both annotations with claim-clustered intervals. Under the primary
annotation the across-class contribution is the only established positive component; the
within-class contribution is additive and has a different denominator from the within-class AUROC
difference of \cref{tab:withinclass}, which is normalised over within-class pairs alone; the
zero-direction contribution is small and established negative, $-0.0010$ \ci{-0.0032}{-0.0001},
which the margin rule's perfect ranking of the eight $p_t=0$ events (all scored $1-|p_t|=1$)
produces; the components sum to the unrounded total $+0.0667$. The percentage shares are
descriptive proportions of the estimated total, not interval-bearing quantities. Under the second
annotation the within-class component is $-0.0781$ \ci{-0.1247}{-0.0386}, 82\% of the total
$-0.0957$ and established, the across-class component $-0.0175$ \ci{-0.0543}{0.0132} is not
established, and the zero-direction component is exactly zero because the ten sparse-history
events are ranked above every negative by both scores; the loss is therefore not the primary
across-class advantage reversing, and the branch-free intervention of \cref{tab:branchfree} acts
on the within-class component. The margin-stratified rates of \cref{tab:classquint} show the
class-rate difference read by strata: no \textsc{protective}-majority event is observed above the
lowest quintile, among units that repeat claim cutoffs and are not independent.

\begin{table}[!t]\centering\scriptsize
\setlength{\tabcolsep}{4pt}
\begin{tabular}{@{}lrlr@{}}
\toprule
Component & Pairs & Contribution [95\% CI] & Share \\
\midrule
\multicolumn{4}{@{}l}{\emph{Primary, shipped A ($\Delta=+.0667$ \ci{.017}{.120})}} \\
Across-class & 34{,}650 & $+$.0609 [.0159, .1081] & 91\% \\
Within-class & 24{,}633 & $+$.0068 [$-$.0039, .0195] & 10\% \\
Zero-direction & 7{,}528 & $-$.0010 [$-$.0032, $-$.0001] & $-$2\% \\
\midrule
\multicolumn{4}{@{}l}{\emph{Second annotation, shipped A ($\Delta=-.0957$ \ci{-.166}{-.037})}} \\
Within-class & -- & $-$.0781 [$-$.1247, $-$.0386] & 82\% \\
Across-class & -- & $-$.0175 [$-$.0543, .0132] & 18\% \\
Zero-direction & -- & 0 [0, 0] & 0\% \\
\bottomrule
\end{tabular}
\caption{Pair decomposition of the paired signed-minus-margin AUROC difference (development cell,
leave-one-claim-out predictions; claim-clustered 95\% intervals over the same 2{,}000 resamples).
Pair counts follow from the class split of \cref{tab:prior} and were verified against the
decomposition script; contributions are additive and sum to the unrounded totals. Second-annotation
pair counts are in the release.}
\label{tab:decomp}
\end{table}

\section{Origin-Defined Cohort and Full Metric Suite}
\label{app:cohortmetrics}
\subsection{Full metric suite}
\label{app:metrics}
\Cref{tab:metrics_source,tab:metrics_forward,tab:metrics_external} report AUROC, AUPRC, log loss,
Brier score, calibration slope and intercept, expected calibration error, and precision and recall
at 5, 10 and 20\% review budgets for the model-cohort combinations evaluated. The offset and sign-invariant
diagnostics are source-only. Stationary rows use uncalibrated operational probabilities; the
score-calibration treatments are specified in the captions (training-only Platt scaling:
leave-one-claim-out on the source, per-origin training folds in the calendar cohort, the frozen
source calibrator on the external corpora). \Cref{tab:matched} is a separate matched-recalibration
experiment. Pooled review-budget metrics are retrospective ranking summaries and rank units by the
uncalibrated score.

\paragraph{Inference sensitivity.} \Cref{tab:wild} recomputes the paired $\Delta$AUROC intervals
of \cref{tab:protocol} under procedures suited to few positive clusters: per-unit DeLong placement
contributions to $\Delta$ are aggregated to claim level and a wild cluster bootstrap-$t$ with Webb
six-point weights (9{,}999 draws) and a leave-one-claim-out jackknife are applied to seven cells
(the development, second-annotation, intersection and three calendar contrasts of
\cref{tab:protocol}, plus signed minus class-blind); a clustered-ROC sensitivity treats claims as
the sampling unit of the ROC estimates themselves: the per-unit DeLong placement contributions
of both scores are aggregated to claim level, and the variance of the paired difference is the
claim-clustered variance of the per-claim contributions to $\hat A_1-\hat A_2$, so that the
covariance between the two AUROC estimates enters through the paired contributions
($\operatorname{Var}(\hat A_1-\hat A_2)=\operatorname{Var}\hat A_1+\operatorname{Var}\hat A_2-2\operatorname{Cov}(\hat A_1,\hat A_2)$)
rather than being dropped; it was run on those seven cells and on the pooled within-class and
two external cells of \cref{tab:protocol}; and a claim-level bootstrap that refits the score inside
every resample, rather than resampling fixed out-of-fold predictions, was run for the four
source-cohort cells (the calendar cells would require rebuilding the per-origin training grid
inside every replicate and were not run)
(2{,}000 resamples; every duplicate copy of a held-out claim drawn by the resampling is removed
from that claim's training folds; no resample is discarded, and a fold whose training side holds a
single class assigns that class's mean). \Cref{tab:wild} names the
procedure behind each column and marks the cells each covered. The tested superiority and loss
decisions against zero are unchanged: the development gain remains established
(\ci{0.018}{0.115} under the wild cluster bootstrap, \ci{0.016}{0.117} under the clustered ROC
variance, \ci{0.014}{0.117} under refitting), the second-annotation loss remains established
(\ci{-0.158}{-0.033} and \ci{-0.242}{-0.039}), the External-2 loss remains established under the
clustered variance (\ci{-0.122}{-0.025}), and no non-established cell becomes established
against zero. One non-inferiority subdecision is procedure-sensitive: External-1 is not
established as non-inferior under the primary percentile interval (lower bound $-0.052$) but is
under the clustered-ROC interval (lower bound $-0.012$, above $-\delta=-0.02$); the protocol
specifies the percentile interval as primary, so the \cref{tab:protocol} verdict stands, and the
complete transfer check and the tier assignments are unchanged because External-2 fails. The
refitting interval for the intersection contrast, \ci{-0.177}{0.020}, is about 1.7 times as wide
as the percentile interval (0.197 against 0.117), an observed widening on a cell with 19 positive
claims rather than an established relationship. The procedures do not give numerically identical
intervals. Their agreement on the decisions against zero supports the stability of those
decisions to the tested inference choices; it does not establish nominal coverage, remove the
retrospective specification of the cells, or bear on annotation validity. The two single-class
within-class cells of \cref{tab:protocol} received the percentile bootstrap only (under the second
annotation the same-class-pairs contrast of \cref{tab:withinclass}, $-0.298$, has clustered
interval \ci{-0.461}{-0.135}), and the A-to-C intervention of \cref{tab:branchfree} received the
percentile bootstrap for C minus A and C minus margin only.

\begin{table*}[!t]\centering\scriptsize
\setlength{\tabcolsep}{3pt}
\begin{tabular}{@{}lrlllllr@{}}
\toprule
Cell & $\Delta$ & Percentile & Wild cluster-$t$ & Jackknife & Clustered ROC & Refitting & Pos.\ cl. \\
\midrule
Prim., signed $-$ margin & $+$.067 & [.017, .120] & [.018, .115] & [.012, .121] & [.016, .117] & [.014, .117] & 23/160 \\
Prim., signed $-$ class-blind & $+$.051 & [.013, .092] & [.013, .088] & [.010, .091] & [.013, .088] & [.015, .120] & 23/160 \\
Second, signed $-$ margin & $-$.096 & [$-$.166, $-$.037] & [$-$.157, $-$.034] & [$-$.163, $-$.028] & [$-$.158, $-$.033] & [$-$.242, $-$.039] & 16/160 \\
Intersect., signed $-$ margin & $-$.030 & [$-$.094, .023] & [$-$.082, .022] & [$-$.094, .034] & [$-$.089, .029] & [$-$.177, .020] & 19/160 \\
Calendar $H{=}3$ & $+$.006 & [$-$.012, .023] & [$-$.011, .023] & [$-$.012, .024] & [$-$.011, .023] & -- & 25/158 \\
Calendar $H{=}5$ & $+$.023 & [$-$.005, .053] & [$-$.003, .050] & [$-$.007, .053] & [$-$.005, .051] & -- & 19/142 \\
Calendar $H{=}10$ & $+$.008 & [$-$.032, .047] & [$-$.028, .044] & [$-$.033, .049] & [$-$.029, .046] & -- & 17/122 \\
\midrule
Pooled within-class pairs (prim.) & $+$.018 & [$-$.012, .060] & -- & -- & [$-$.012, .049] & -- & 23/160 \\
External-1, frozen & $+$.055 & [$-$.052, .106] & -- & -- & [$-$.012, .121] & -- & 7/46 \\
External-2, frozen & $-$.074 & [$-$.123, $-$.024] & -- & -- & [$-$.122, $-$.025] & -- & 40/180 \\
\bottomrule
\end{tabular}
\caption{Paired $\Delta$AUROC intervals under four inference procedures beyond the claim-clustered
percentile bootstrap of the main text (2{,}000 resamples): a wild cluster bootstrap-$t$ over
claims with Webb six-point weights (9{,}999 draws), a leave-one-claim-out jackknife, a clustered-ROC
sensitivity (claim-level DeLong variance of each AUROC; run for all ten cells), and a claim-level
bootstrap with the score refitted in every resample (2{,}000 resamples, seed 3; run for the four
source-cohort cells). ``--'': not run. ``Pos.\ cl.'': claims with at least one event / claims in
the cell. The two single-class within-class cells of \cref{tab:protocol} received the percentile
bootstrap only.}
\label{tab:wild}
\end{table*}

\begin{table*}[!t]\centering\scriptsize
\setlength{\tabcolsep}{2.5pt}
\resizebox{\textwidth}{!}{%
\begin{tabular}{@{}lrrrrrrrrrrrrrrr@{}}
\toprule
Score & AUROC & AUPRC & Log loss & Brier & Cal.\ slope & Cal.\ int. & ECE & P@5 & R@5 & P@10 & R@10 & P@20 & R@20 & NNR@10 & Share@R.8 \\
\midrule
\multicolumn{16}{@{}l}{\emph{Source, claim-held-out (1{,}012 units, 71 events), 1012 units, 71 events}} \\
Margin rule $1-|p_t|$, untrained & .881 & .468 & .184 & .0515 & .924 & -.169 & .012 & .490 & .352 & .386 & .549 & .267 & .761 & 2.590 & .222 \\
Stationary Beta-Binomial, count-predicted (operational) & .909 & .497 & .200 & .0499 & .534 & .002 & .033 & .510 & .366 & .386 & .549 & .282 & .803 & 2.590 & .179 \\
Beta-Binomial, realised future count (diagnostic) & .882 & .490 & .272 & .0491 & .314 & -.515 & .033 & .549 & .394 & .406 & .577 & .272 & .775 & 2.463 & .215 \\
Dirichlet-multinomial, count-predicted & .893 & .492 & .223 & .0497 & .395 & -.346 & .030 & .510 & .366 & .386 & .549 & .277 & .789 & 2.590 & .205 \\
Two-variable logistic ($p$, $a$) & .948 & .535 & .146 & .0452 & .936 & -.090 & .027 & .529 & .380 & .475 & .676 & .327 & .930 & 2.104 & .139 \\
\compact logistic & .946 & .527 & .149 & .0460 & .936 & -.092 & .014 & .549 & .394 & .455 & .648 & .327 & .930 & 2.196 & .141 \\
Random forest (all 55) & .938 & .545 & .149 & .0452 & .949 & -.059 & .014 & .529 & .380 & .475 & .676 & .302 & .859 & 2.104 & .151 \\
XGBoost (all 55) & .932 & .543 & .152 & .0440 & .958 & -.057 & .013 & .608 & .437 & .485 & .690 & .302 & .859 & 2.061 & .151 \\
Stationary offset $+$ order features & .915 & .514 & .167 & .0477 & .668 & -.461 & .022 & .549 & .394 & .436 & .620 & .282 & .803 & 2.295 & .181 \\
Sign-invariant logistic & .913 & .466 & .168 & .0495 & .952 & -.075 & .014 & .510 & .366 & .426 & .606 & .287 & .817 & 2.349 & .177 \\
Two-variable, flip-augmented & .901 & .436 & .176 & .0513 & .932 & -.109 & .006 & .490 & .352 & .396 & .563 & .277 & .789 & 2.525 & .202 \\
\bottomrule
\end{tabular}}
\caption{Full metric suite, source cohort(s). P@$b$ and R@$b$: precision and recall when the $b$\% highest-scored units are reviewed; NNR@10: units reviewed per event found at the 10\% budget; Share@R.8: share of units that must be reviewed to reach recall 0.8 (workload reduction is one minus this share). Calibration slope and intercept are from a logistic regression of the outcome on the logit of the calibrated probability; ECE uses ten equal-width bins.}
\label{tab:metrics_source}
\end{table*}

\begin{table*}[!t]\centering\scriptsize
\setlength{\tabcolsep}{2.5pt}
\resizebox{\textwidth}{!}{%
\begin{tabular}{@{}lrrrrrrrrrrrrrrr@{}}
\toprule
Score & AUROC & AUPRC & Log loss & Brier & Cal.\ slope & Cal.\ int. & ECE & P@5 & R@5 & P@10 & R@10 & P@20 & R@20 & NNR@10 & Share@R.8 \\
\midrule
\multicolumn{16}{@{}l}{\emph{Calendar cohort, 3 y, 756 units, 58 events}} \\
Margin rule $1-|p_t|$, untrained & .907 & .405 & .193 & .0573 & 1.151 & .272 & .026 & .447 & .293 & .368 & .483 & .298 & .776 & 2.714 & .208 \\
Stationary Beta-Binomial, count-predicted (operational) & .920 & .448 & .184 & .0549 & .674 & .145 & .028 & .500 & .328 & .408 & .534 & .318 & .828 & 2.452 & .185 \\
Two-variable logistic ($p$, $a$) & .913 & .351 & .179 & .0554 & 1.121 & .117 & .015 & .395 & .259 & .408 & .534 & .305 & .793 & 2.452 & .206 \\
\compact logistic & .890 & .385 & .203 & .0575 & .778 & -.332 & .020 & .474 & .310 & .382 & .500 & .305 & .793 & 2.621 & .218 \\
\midrule
\multicolumn{16}{@{}l}{\emph{Calendar cohort, 5 y, 625 units, 45 events}} \\
Margin rule $1-|p_t|$, untrained & .897 & .372 & .195 & .0572 & .912 & -.021 & .022 & .419 & .289 & .323 & .444 & .256 & .711 & 3.100 & .246 \\
Stationary Beta-Binomial, count-predicted (operational) & .901 & .394 & .208 & .0550 & .485 & -.172 & .031 & .484 & .333 & .355 & .489 & .256 & .711 & 2.818 & .227 \\
Two-variable logistic ($p$, $a$) & .920 & .380 & .169 & .0522 & 1.168 & .385 & .014 & .387 & .267 & .403 & .556 & .288 & .800 & 2.480 & .190 \\
\compact logistic & .893 & .364 & .198 & .0590 & .674 & -.629 & .041 & .323 & .222 & .387 & .533 & .256 & .711 & 2.583 & .218 \\
\midrule
\multicolumn{16}{@{}l}{\emph{Calendar cohort, 10 y, 492 units, 35 events}} \\
Margin rule $1-|p_t|$, untrained & .867 & .285 & .210 & .0617 & .749 & -.386 & .028 & .320 & .229 & .286 & .400 & .224 & .629 & 3.500 & .256 \\
Stationary Beta-Binomial, count-predicted (operational) & .863 & .309 & .298 & .0598 & .314 & -.468 & .041 & .440 & .314 & .286 & .400 & .235 & .657 & 3.500 & .293 \\
Two-variable logistic ($p$, $a$) & .876 & .303 & .190 & .0564 & .947 & -.068 & .010 & .360 & .257 & .347 & .486 & .255 & .714 & 2.882 & .234 \\
\compact logistic & .797 & .218 & .381 & .0942 & .241 & -1.820 & .090 & .240 & .171 & .265 & .371 & .184 & .514 & 3.769 & .392 \\
\bottomrule
\end{tabular}}
\caption{Full metric suite, forward cohort(s). P@$b$ and R@$b$: precision and recall when the $b$\% highest-scored units are reviewed; NNR@10: units reviewed per event found at the 10\% budget; Share@R.8: share of units that must be reviewed to reach recall 0.8 (workload reduction is one minus this share). Calibration slope and intercept are from a logistic regression of the outcome on the logit of the calibrated probability; ECE uses ten equal-width bins.}
\label{tab:metrics_forward}
\end{table*}

\begin{table*}[!t]\centering\scriptsize
\setlength{\tabcolsep}{2.5pt}
\resizebox{\textwidth}{!}{%
\begin{tabular}{@{}lrrrrrrrrrrrrrrr@{}}
\toprule
Score & AUROC & AUPRC & Log loss & Brier & Cal.\ slope & Cal.\ int. & ECE & P@5 & R@5 & P@10 & R@10 & P@20 & R@20 & NNR@10 & Share@R.8 \\
\midrule
\multicolumn{16}{@{}l}{\emph{External-1, frozen, 322 units, 19 events}} \\
Margin rule $1-|p_t|$, untrained & .893 & .505 & .157 & .0415 & 1.192 & -.076 & .033 & .438 & .368 & .281 & .474 & .219 & .737 & 3.556 & .261 \\
Stationary Beta-Binomial, count-predicted (operational) & .907 & .507 & .160 & .0442 & .609 & -.397 & .045 & .438 & .368 & .281 & .474 & .234 & .789 & 3.556 & .211 \\
Two-variable logistic ($p$, $a$) & .948 & .679 & .114 & .0330 & 1.073 & .036 & .019 & .563 & .474 & .469 & .789 & .281 & .947 & 2.133 & .109 \\
\compact logistic & .951 & .662 & .112 & .0326 & 1.188 & .030 & .016 & .563 & .474 & .438 & .737 & .281 & .947 & 2.286 & .109 \\
\midrule
\multicolumn{16}{@{}l}{\emph{External-2, frozen, 1257 units, 95 events}} \\
Margin rule $1-|p_t|$, untrained & .896 & .487 & .183 & .0522 & 1.252 & .512 & .017 & .524 & .347 & .421 & .558 & .283 & .747 & 2.377 & .231 \\
Stationary Beta-Binomial, count-predicted (operational) & .898 & .486 & .207 & .0543 & .626 & .148 & .035 & .540 & .358 & .429 & .568 & .267 & .705 & 2.333 & .242 \\
Two-variable logistic ($p$, $a$) & .823 & .267 & .265 & .0742 & .439 & -1.343 & .052 & .333 & .221 & .341 & .453 & .243 & .642 & 2.930 & .345 \\
\compact logistic & .819 & .257 & .260 & .0739 & .470 & -1.312 & .051 & .302 & .200 & .333 & .442 & .223 & .589 & 3.000 & .353 \\
\bottomrule
\end{tabular}}
\caption{Full metric suite, external cohort(s). P@$b$ and R@$b$: precision and recall when the $b$\% highest-scored units are reviewed; NNR@10: units reviewed per event found at the 10\% budget; Share@R.8: share of units that must be reviewed to reach recall 0.8 (workload reduction is one minus this share). Calibration slope and intercept are from a logistic regression of the outcome on the logit of the calibrated probability; ECE uses ten equal-width bins.}
\label{tab:metrics_external}
\end{table*}


\subsection{Origin-defined calendar cohort}
\label{app:cohort}
\Cref{tab:census} is the census of the cohort of \cref{sec:task}: every claim with at least
40 records before the origin is a candidate; accrual is at least ten signed records in the window;
windows ending after 2025 are excluded rather than evaluated on a partially observed year (the
stored corpus holds 5{,}818 records dated 2026). The accrual model is a class-balanced logistic
regression on $\log(1+n_t)$, $\log(1+\text{accrual}_t)$ and $\log(1+n_{\text{pre}})$ fitted on
all earlier cohort units of the other claims. Training for the change model uses every earlier
unit with an observable outcome on a two-year cutoff grid, with the test claim's own units removed;
one fit per test claim and origin. \Cref{fig:calendar} shows the per-origin values.

\paragraph{Accrual-weighted sensitivity.} Accrued and censored candidates differ on pre-origin
features, most in recent accrual (up to 1.13 standard deviations). Weighting each accrued unit by the
inverse of its predicted probability of accrual, estimated from pre-origin features on claim-disjoint
units, moves no paired difference by more than 0.014, and every weighted comparison reaches the same
conclusion as \cref{tab:protocol}. At ten years no candidate is censored, so no weighting is needed.

\begin{table*}[!t]\centering\scriptsize
\setlength{\tabcolsep}{3pt}
\resizebox{\textwidth}{!}{%
\begin{tabular}{@{}llrrrrrll@{}}
\toprule
Origin & Window & Candidates & Accrued & Censored & Events & Rate & AUROC among accrued: stationary / margin / two-var / \compact & Status \\
\midrule
\multicolumn{9}{@{}l}{\emph{Horizon 3 years}} \\
2004 & [2004, 2007) & 68 & 65 & 3 & 4 & .062 & .975 / .939 / .943 / .930 & complete \\
2007 & [2007, 2010) & 88 & 82 & 6 & 10 & .122 & .926 / .926 / .933 / .924 & complete \\
2010 & [2010, 2013) & 99 & 92 & 7 & 7 & .076 & .940 / .916 / .887 / .903 & complete \\
2013 & [2013, 2016) & 115 & 112 & 3 & 11 & .098 & .908 / .904 / .897 / .820 & complete \\
2016 & [2016, 2019) & 122 & 117 & 5 & 7 & .060 & .916 / .929 / .933 / .935 & complete \\
2019 & [2019, 2022) & 143 & 136 & 7 & 9 & .066 & .907 / .877 / .901 / .893 & complete \\
2022 & [2022, 2025) & 165 & 152 & 13 & 10 & .066 & .928 / .914 / .934 / .929 & complete \\
\midrule
\multicolumn{9}{@{}l}{\emph{Horizon 5 years}} \\
2004 & [2004, 2009) & 68 & 66 & 2 & 4 & .061 & .964 / .927 / .932 / .875 & complete \\
2007 & [2007, 2012) & 88 & 86 & 2 & 10 & .116 & .882 / .890 / .897 / .842 & complete \\
2010 & [2010, 2015) & 99 & 97 & 2 & 8 & .083 & .898 / .882 / .930 / .892 & complete \\
2013 & [2013, 2018) & 115 & 113 & 2 & 9 & .080 & .933 / .928 / .937 / .922 & complete \\
2016 & [2016, 2021) & 122 & 121 & 1 & 7 & .058 & .902 / .925 / .946 / .917 & complete \\
2019 & [2019, 2024) & 143 & 142 & 1 & 7 & .049 & .892 / .889 / .921 / .941 & complete \\
2022 & [2022, 2027) & 165 & -- & -- & -- & -- & -- & window extends past 2025: not evaluated \\
\midrule
\multicolumn{9}{@{}l}{\emph{Horizon 10 years}} \\
2004 & [2004, 2014) & 68 & 68 & 0 & 5 & .074 & .924 / .844 / .740 / .816 & complete \\
2007 & [2007, 2017) & 88 & 88 & 0 & 8 & .091 & .875 / .905 / .900 / .798 & complete \\
2010 & [2010, 2020) & 99 & 99 & 0 & 5 & .051 & .832 / .860 / .892 / .753 & complete \\
2013 & [2013, 2023) & 115 & 115 & 0 & 8 & .070 & .832 / .838 / .851 / .771 & complete \\
2016 & [2016, 2026) & 122 & 122 & 0 & 9 & .074 & .884 / .883 / .920 / .898 & complete \\
2019 & [2019, 2029) & 143 & -- & -- & -- & -- & -- & window extends past 2025: not evaluated \\
2022 & [2022, 2032) & 165 & -- & -- & -- & -- & -- & window extends past 2025: not evaluated \\
\bottomrule
\end{tabular}}
\caption{Census of the origin-defined calendar cohort. Candidates are claims with at least 40 records before the origin; censored units have fewer than ten signed records in the window; the rate is the event rate among accrued units. Per-origin AUROCs are claim-disjoint and rest on 4-11 events. Accrual and events are not reported for windows that extend past 2025.}
\label{tab:census}
\end{table*}


\begin{figure*}[!t]
\centering
\includegraphics[width=\textwidth]{figures/fig5_calendar_cohort.pdf}
\caption{Origin-defined calendar cohort: AUROC among accrued units by origin and horizon for the
stationary Beta-Binomial baseline, the margin rule, the two-variable score and \compact (claim-disjoint
training, complete windows only). Labels give events/accrued units per origin; per-origin values
rest on 4-11 events and are indicative only.}
\label{fig:calendar}
\end{figure*}

\section{Frozen External Transfer and the Two-Variable Score}
\label{app:crossdomain}
\begin{table}[!t]\centering\scriptsize
\setlength{\tabcolsep}{3pt}
\resizebox{\columnwidth}{!}{%
\begin{tabular}{@{}llll@{}}
\toprule
Domain & Score & AUROC [95\% CI] & AUPRC [95\% CI] \\
\midrule
External-1 & \compact, frozen & \textbf{.951} [.906, .993] & \textbf{.662} [.417, .874] \\
(322 units, 19 events) & $1-|p_t|$ & .893 [.838, .977] & .505 [.309, .811] \\
 & Low agreement & .856 [.796, .965] & .437 [.200, .693] \\
 & Recent dissent & .737 [.371, .963] & .341 [.052, .662] \\
\midrule
External-2 & \compact, frozen & .819 [.758, .871] & .257 [.167, .387] \\
(1{,}257 units, 95 events) & $1-|p_t|$ & \textbf{.896} [.856, .932] & \textbf{.487} [.358, .599] \\
 & Low agreement & .882 [.839, .916] & .392 [.257, .529] \\
 & Recent dissent & .887 [.852, .917] & .349 [.253, .454] \\
\bottomrule
\end{tabular}}
\caption{Frozen external categorical validation under the tie convention of
\cref{eq:target}. \compact is fitted once on the 1{,}012 source future-eligible units and never
refitted on target-domain labels; the three unfitted signals need no fitting. Intervals are
claim-clustered over the external claims.}
\label{tab:external}
\end{table}

Refitting a balanced logistic model on eight named signals separately in each domain keeps the
signs of agreement (negative), heterogeneity $I^2$ (negative) and volume (positive) in all three
domains, while recent dissent, design mix and the decline signal change sign; a refitted
coefficient's sign is weak evidence about the transport of a frozen score, so this is recorded
only as context for \cref{tab:external} (the full coefficient table is in the released
results).

\paragraph{The two-variable score.} It is the same class-balanced $L_2$ logistic regression as
\compact, fitted on standardised $[p,a]_t$ only. \cref{tab:astra5} reports it
leave-one-claim-out on the source under each agreement definition of
\cref{app:agreement}, with paired differences from \compact, and frozen (fitted once on
the 1{,}012 source units, never refitted) on External-1 and External-2, where agreement is rebuilt
from the released label counts under the same definition. Under the shipped definition A it gives
0.948 against 0.946 on the source ($\Delta$ $+0.002$ \ci{-0.006}{0.011}), 0.948 against 0.951 on
External-1 ($-0.004$ \ci{-0.017}{0.002}) and 0.823 against 0.819 on External-2 ($+0.004$
\ci{-0.004}{0.016}); on External-2 the unfitted $1-|p_t|$ (0.896) is above both frozen scores
in AUROC ($-0.074$ \ci{-0.123}{-0.024} for the two-variable score) and in AUPRC ($-0.220$
\ci{-0.308}{-0.104}), and the frozen scores are above it on External-1 without an established gain
($+0.055$ \ci{-0.052}{0.106}), so dropping the count and effect variables does not remove the
transfer gap. The gap is a transfer loss rather than an inability of the two inputs to rank the
domain: refitted inside External-2 leave-one-claim-out, the two-variable score's difference from
the margin rule is $-0.017$ \ci{-0.047}{0.009} (External-1: $+0.054$ \ci{-0.068}{0.116}), so
letting the model adapt to the target domain recovers most of what the frozen check measures.

\begin{table*}[!t]\centering\scriptsize
\setlength{\tabcolsep}{3pt}
\begin{tabular}{@{}llrrlllll@{}}
\toprule
Def. & Domain & Units & Events & Two-var AUROC [95\% CI] & AUPRC & C-7, same def. & $\Delta$ vs C-7 [95\% CI] & $1-|p_t|$ \\
\midrule
A & source & 1012 & 71 & .948 [.921, .971] & .535 & .946 & $+$.002 (no interval) & .881 \\
A & External-1 & 322 & 19 & .948 [.892, .994] & .679 & .951 & $-$.004 [$-$.017, .002] & .893 \\
A & External-2 & 1257 & 95 & .823 [.766, .874] & .267 & .819 & $+$.004 [$-$.004, .016] & .896 \\
B & source & 1012 & 71 & .890 [.830, .939] & .387 & .899 & $-$.009 (no interval) & .881 \\
B & External-1 & 322 & 19 & .895 [.812, .990] & .478 & .893 & $+$.002 [$-$.020, .025] & .893 \\
B & External-2 & 1257 & 95 & .670 [.608, .733] & .115 & .680 & $-$.010 [$-$.032, .015] & .896 \\
C & source & 1012 & 71 & .925 [.890, .959] & .500 & .925 & $-$.001 (no interval) & .881 \\
C & External-1 & 322 & 19 & .964 [.932, .994] & .657 & .961 & $+$.003 [$-$.004, .008] & .893 \\
C & External-2 & 1257 & 95 & .841 [.785, .894] & .398 & .833 & $+$.009 [$-$.006, .029] & .896 \\
D & source & 1012 & 71 & .849 [.758, .922] & .349 & .858 & $-$.009 (no interval) & .881 \\
D & External-1 & 322 & 19 & .895 [.794, .994] & .514 & .876 & $+$.018 [$-$.023, .065] & .893 \\
D & External-2 & 1257 & 95 & .601 [.540, .667] & .085 & .646 & $-$.045 [$-$.081, $-$.006] & .896 \\
\bottomrule
\end{tabular}
\caption{Two-variable score under each agreement definition (A shipped; B three-way maximum share;
C signed-only share; D entropy consensus), leave-one-claim-out on the source and frozen on the
external corpora with agreement rebuilt from the released label counts. The comparator is
\compact rebuilt under the \emph{same} definition (external rows: paired interval on identical
claim resamples; source rows: point difference only). The paired source differences of each
definition's two-variable score from the definition-A score ($-0.058$, $-0.023$, $-0.099$ for B, C,
D) are in \cref{tab:astra2}. The margin rule is definition-independent. On External-2 the
branch-free score (0.841) remains below the margin rule (0.896).}
\label{tab:astra5}
\end{table*}

\section{Stationary Null Models}
\label{app:null}
Two null models test whether the observed ranking could arise from a claim's stationary label
composition without any temporal organisation. Null A permutes each claim's direction labels
across its records (publication years fixed); Null B draws each record's label i.i.d.\ from the
claim's own label frequencies, so Null A preserves each claim's realised label counts exactly while
Null B preserves the generating probabilities and expected frequencies, not the exact counts of
every draw. For every draw the direction block, \compact's
direction inputs and the target are rebuilt from the permuted or simulated stream, the text-derived
effect-estimate inputs are kept, and the full leave-one-claim-out pipeline is rerun for \compact
and the two-variable model; the pipeline reproduces the observed 0.946 exactly when run on the
real stream. Each null is run for 1{,}000 draws split over eight seeds. All 1{,}000 draws of each null were evaluable. Null A: \compact null mean 0.936 (SD 0.023,
95\% null interval \ci{0.882}{0.967}), two-variable 0.948 \ci{0.907}{0.974}, unfitted $1-|p_t|$
0.957 \ci{0.929}{0.977}, \compact AUPRC 0.351 \ci{0.192}{0.519}, events 35.3 \ci{24}{47},
prevalence 0.035. Null B: 0.935 \ci{0.873}{0.970},
0.947 \ci{0.898}{0.975}, 0.957 \ci{0.924}{0.979}, AUPRC 0.362 \ci{0.171}{0.570}, events 36.1
\ci{21}{52}, prevalence 0.036. Observed: 0.946, 0.948, 0.881, 0.527,
71, 0.070. Empirical upper-tail $p$-values: \compact AUROC 0.36 and 0.40,
two-variable AUROC 0.57 and 0.55, \compact AUPRC 0.024 and 0.067, events and prevalence 0.001 under both
nulls; the unfitted $1-|p_t|$ is
\emph{below} its null (lower-tail $p=0.002$ and $0.001$). The observed fitted-score AUROCs are not distinguishable from these null distributions, which
is not an equivalence test and does not show that the histories contain no temporal information;
it shows that high AUROC alone does not establish incremental temporal-order information. The
nulls are rejected on the number of events: the observed annotation streams exhibit sign-change events about twice as often as a stationary process with the same
composition, and those changes are less tied to $|p_t|$, which is why the margin rule ranks the
observed outcomes (0.881) worse than the null outcomes (0.957) while the fitted scores do not.
\cref{tab:astra7} and \cref{fig:nullgrid} report the null distributions, the observed
values and the empirical upper-tail $p$-values, $(\#\{\text{draws}\ge\text{observed}\}+1)/(n+1)$.
Because the nulls keep every claim's label composition, a rejected null would show only that the
ranking depends on the temporal ordering of the labels, not that a causal temporal mechanism
exists. The quantities that are distinguishable are the event count and the prevalence. The difference
between the two-variable and margin-rule null means is $-0.009$ (Null A) and $-0.010$ (Null B),
against $+0.067$ observed; paired within each draw, the observed difference lies above the null distribution
($p=0.001$). AUPRC is treated separately: it depends on both ranking and prevalence,
the prevalence differs between observed and null streams, one \compact contrast ($p=0.067$) and
both two-variable contrasts (0.112 and 0.189) do not cross 0.05, and no common-prevalence
comparison is available, so no temporal-discrimination gain is inferred from AUPRC.

\begin{table}[!t]\centering\scriptsize
\setlength{\tabcolsep}{3pt}
\resizebox{\columnwidth}{!}{%
\begin{tabular}{@{}lrrrrrr@{}}
\toprule
Null & Metric & Observed & Null mean & Null 95\% interval & $p$ (observed $\ge$ null) & Draws \\
\midrule
Null A: within-claim permutation & Compact-7 AUROC & .946 & .936 & [0.882,\,0.967] & .365 & 1000 \\
Null A: within-claim permutation & Compact-7 AUPRC & .527 & .351 & [0.192,\,0.519] & .024 & 1000 \\
Null A: within-claim permutation & 2-variable AUROC & .948 & .948 & [0.907,\,0.974] & .569 & 1000 \\
Null A: within-claim permutation & 2-variable AUPRC & .535 & .424 & [0.252,\,0.600] & .112 & 1000 \\
Null A: within-claim permutation & 1-|p\_t| AUROC & .881 & .957 & [0.929,\,0.977] & .999 & 1000 \\
Null A: within-claim permutation & event prevalence & .070 & .035 & [0.024,\,0.046] & .001 & 1000 \\
Null A: within-claim permutation & events & 71.000 & 35.299 & [24.000,\,47.000] & .001 & 1000 \\
Null B: stationary voting & Compact-7 AUROC & .946 & .935 & [0.873,\,0.970] & .396 & 1000 \\
Null B: stationary voting & Compact-7 AUPRC & .527 & .362 & [0.171,\,0.570] & .067 & 1000 \\
Null B: stationary voting & 2-variable AUROC & .948 & .947 & [0.898,\,0.975] & .547 & 1000 \\
Null B: stationary voting & 2-variable AUPRC & .535 & .432 & [0.223,\,0.634] & .189 & 1000 \\
Null B: stationary voting & 1-|p\_t| AUROC & .881 & .957 & [0.924,\,0.979] & 1.000 & 1000 \\
Null B: stationary voting & event prevalence & .070 & .036 & [0.021,\,0.051] & .001 & 1000 \\
Null B: stationary voting & events & 71.000 & 36.099 & [21.000,\,52.000] & .001 & 1000 \\
\bottomrule
\end{tabular}}
\caption{Stationary-voting and permutation nulls. Null A permutes each claim's direction labels across its publication years; Null B draws labels i.i.d.\ from each claim's stationary label distribution over the observed publication times. Every draw rebuilds the benchmark and reruns the leave-one-claim-out pipeline.}
\label{tab:astra7}
\end{table}

\begin{figure*}[!t]
\centering
\includegraphics[width=\textwidth]{figures/fig6_null_audit.pdf}
\caption{Null~A and Null~B distributions under the primary annotation for the two-variable AUROC,
the margin-rule AUROC and the event count (1{,}000 draws per null), with the observed values marked.
Nulls~C and~D and the second and intersection annotations are summarised in \cref{tab:null,tab:nullann,tab:nullblocks}.}
\label{fig:nullgrid}
\end{figure*}

\paragraph{Every annotation, and drift-preserving controls.} A reimplementation that rebuilds the
target and reruns the leave-one-claim-out two-variable model (not \compact) extends both nulls to
the second annotation and intersection labels, and adds Null~C, which permutes labels within claim and
decade so that the era-level label mix of every claim is preserved, 1{,}000 draws each
(\cref{tab:nullann}). On the primary annotation it gives 35.6
\ci{26}{46} and 36.5 \ci{20}{52} null events, against the archived 35.3 and 36.1. The excess of events
over Nulls A and B holds under the second annotation for both nulls and under intersection labels for
Null~A only (Null~B: $p=0.054$). Null~C absorbs about 72\% of the primary excess
(71 against 61.2 \ci{53}{69}, $p=0.008$); under the second annotation and intersection labels the
observed counts fall inside its range ($p=0.475$ and $0.627$ in this reimplementation run;
$0.489$ and $0.641$ in the independent partition-grid run of \cref{tab:nullblocks}, a Monte Carlo
difference between two draws of 1{,}000). Null~C also reproduces the signed
score's paired gain over the margin rule under the primary annotation ($+0.067$ observed against a
null mean of $+0.065$ \ci{0.047}{0.083}, $p=0.427$), whereas under Nulls A and B the same difference
has $p=0.001$: preserving the era-level label mix suffices to reproduce the gain under this scheme, which does not identify the drift's contribution. The
paired signed-minus-margin difference follows the same ordering where it is negative: under the
second annotation the observed $-0.096$ lies above its Null~A mean of $-0.204$ ($p=0.004$), so the
signed score loses to the margin rule, but by less than on exchangeable streams.

\paragraph{What Null~C conditions on.} Null~C permutes the labels of a claim's records, of every
label type including \textsc{null} and \textsc{unclear}, among the record positions that fall in
the same claim-block cell, so that each cell keeps its label totals and every record keeps its
date. Labels in complete blocks on either side of a cutoff therefore cannot cross it; only the
allocation within the block that straddles the cutoff can change a unit's pre- and post-cutoff
counts, and a cutoff on a block boundary leaves them, and hence the endpoint and the two
count-derived inputs, unchanged. The control tests what remains after conditioning on the observed
claim-block label totals and does not identify the causal contribution of temporal drift. An
invariance audit counts the units whose invariance is guaranteed by the date partition alone:
every record of the straddling block, whatever its label, lies on one side of the cutoff, or the
cutoff falls on a block boundary (a condition on signed records alone would not suffice, since a
permutation can move a signed label onto the position of a \textsc{null} or \textsc{unclear}
record on the other side). The condition is sufficient, not necessary: a straddling block whose
records all carry one label type, or in which the same label appears on both sides of the cutoff
with no other label present, also cannot change its counts under any permutation, and such
homogeneous crossing blocks are not counted, so the reported shares are lower bounds on the
fraction of units whose counts cannot change. Under the ten-year partition 9.6\% of units,
holding 3 of the 71 primary events, are invariant by the partition alone, so the reproduction of
the event count is not largely guaranteed by that invariance; five-year blocks condition on more
temporal structure and fix 31.2\% of units and 20 events. A unit that is not invariant may still keep its
endpoint in most draws, so ``not invariant'' is not ``substantially perturbed''.
\paragraph{Partition sensitivity.} The decision convention throughout, as in \cref{tab:null},
is the central 95\% null range: an observed value is outside it when its upper-tail probability
is below 0.025, not 0.05. \Cref{tab:nullblocks} varies the block for the primary annotation: with
five-year blocks the observed 71 events fall inside the null range (64.9 \ci{59}{71}, $p=0.027$),
as they do with decade boundaries shifted by five years (64.5 \ci{58}{71}, $p=0.038$), both
probabilities below 0.05 but above the 0.025 threshold, while twenty-year blocks widen the gap
(56.3 \ci{48}{64}, $p=0.002$). Under the second annotation and intersection labels the observed
counts (52 and 55) lie inside the null range in every one of the four partitions. The excess over
era-conditioned controls is therefore confined to the primary annotation and depends on the
partition in both its magnitude and its decision at the 0.025 threshold; this restriction concerns the era-conditioned comparisons and does not remove the excess
over Nulls A and B that the second and intersection annotations also show (\cref{tab:nullann}).
The grid evaluates event counts under all three annotations; the paired signed-minus-margin
difference was evaluated across partitions under the primary annotation only, where it is
reproduced under every block ($p$ from 0.087 to 0.427), and no partition robustness of the learned
advantage is claimed under the other annotations.
\paragraph{Null~D.} Null~D replaces the block conditioning by a parametric trend. For each claim a
logistic regression of the probability that a signed record is \textsc{harmful} on publication year
is fitted to that claim's signed labels over the archived history by maximum likelihood, with a
numerical ridge of $10^{-6}$ added to the information matrix as optimizer damping only, not as a
penalty in the objective. For a claim with a single signed class or complete separation the
unpenalised likelihood has no finite maximiser, and the fit used is the numerically truncated
iterate at the optimizer's stopping rule, whose fitted probabilities lie near 0 or 1; those
stopping choices therefore help define the generated probabilities for such claims, and the
control is not a penalised estimator. The fitted parameters are then held fixed, and in each of 1{,}000 draws every signed record's label is
redrawn from the fitted probability at its own date while \textsc{null} and \textsc{unclear}
records keep their labels and positions, after which the target is rebuilt and the score refitted.
% AUTHORS: insert the optimizer's stopping rule (tolerance and maximum iterations), any probability
% clipping, and the number of single-class or non-converged claims, from the Null D script.
The control therefore carries no parameter uncertainty, is fitted on the archived history rather
than on origin-available data, and is a retrospective diagnostic of whether a smooth per-claim
drift in the signed label mix accounts for the observed quantities, not an origin-available
reference and not a generative model of conclusion change. It is not rejected on the event count
(70.1 \ci{54}{87}, $p=0.458$) or on the gain ($+0.040$ against $+0.067$, $p=0.099$); a failure to
reject does not establish it as the true generating process. The residual of about ten reversals
reported for the ten-year Null~C is therefore a conditional, partition-specific contrast, not an
identified set of unexplained reversals, and is not carried in the main text.
\begin{table}[!t]\centering\scriptsize
\setlength{\tabcolsep}{3pt}
\begin{tabular}{@{}llrrr@{}}
\toprule
Control & Events [range] & $p$ & Paired $\Delta$ & $p$ \\
\midrule
C, 10-year blocks & 61.2 [53, 69] & .008 & $+$.065 & .427 \\
C, 5-year blocks & 64.9 [59, 71] & .027 & $+$.065 & .427 \\
C, shifted 5 y & 64.5 [58, 71] & .038 & $+$.059 & .131 \\
C, 20-year blocks & 56.3 [48, 64] & .002 & $+$.051 & .087 \\
D, fitted trend & 70.1 [54, 87] & .458 & $+$.040 & .099 \\
\midrule
\multicolumn{5}{@{}l}{\emph{Second annotation (52 events; counts only)}} \\
C, 10-year blocks & 51.3 [45, 57] & .489 & -- & -- \\
C, 5-year blocks & 52.3 [48, 57] & .645 & -- & -- \\
C, shifted 5 y & 52.5 [47, 58] & .633 & -- & -- \\
C, 20-year blocks & 48.4 [41, 55] & .202 & -- & -- \\
\midrule
\multicolumn{5}{@{}l}{\emph{Intersection labels (55 events; counts only)}} \\
C, 10-year blocks & 55.7 [49, 62] & .641 & -- & -- \\
C, 5-year blocks & 54.7 [50, 60] & .541 & -- & -- \\
C, shifted 5 y & 53.8 [48, 60] & .401 & -- & -- \\
C, 20-year blocks & 55.9 [48, 64] & .661 & -- & -- \\
\bottomrule
\end{tabular}
\caption{Partition sensitivity of Null~C and the trend-preserving Null~D (1{,}000 draws each;
$p$: upper-tail probability of the observed count or paired gain; an observed value lies outside
the central 95\% null range when $p<0.025$). Primary annotation: 71
events, $\Delta=+0.067$; invariant units / events under the all-records condition of the text:
9.6\% / 3 (10-year), 31.2\% / 20 (5-year), lower bounds that exclude homogeneous crossing blocks;
not computed for the shifted and 20-year partitions. Learned-score differences were evaluated across partitions under the primary annotation only.}
\label{tab:nullblocks}
\end{table}

\begin{table}[!t]\centering\footnotesize
\setlength{\tabcolsep}{3pt}
\resizebox{\columnwidth}{!}{%
\begin{tabular}{@{}lrlll@{}}
\toprule
Annotation & Observed & Null~A & Null~B & Null~C \\
\midrule
Primary & 71 & 35.6 [26, 46] & 36.5 [20, 52] & 61.2 [53, 69] \\
Second & 52 & 34.3 [24, 45] & 36.0 [21, 51] & 51.3 [45, 57] \\
Intersection & 55 & 40.2 [29, 51] & 40.9 [25, 56] & 55.7 [49, 62] \\
\bottomrule
\end{tabular}}
\caption{Events under the exchangeable controls for every annotation (null mean [95\% range];
reimplementation, 1{,}000 draws per null). The archived primary values of \cref{tab:astra7} are 35.3 and 36.1.}
\label{tab:nullann}
\end{table}

\section{Complete Target Robustness}
\label{app:robust}
\paragraph{Zero-direction convention.} Eleven future-eligible source units have a zero direction on
exactly one side of the target: eight with $p_t=0$, of which seven are supported exact ties (15/15,
16/16, 18/18, 13/13 and 57/57 signed pre-cutoff records among them) and one is the sparse-history
fallback (fewer than eight signed records), and three with $p^{+}_t=0$ (47/47, 18/18 and 135/135 signed post-cutoff
records). Under \cref{eq:target} a unit with a zero direction on exactly one side is an event, because a resolved
direction differs from no direction, whereas a unit with zero direction on both sides would not be; all eleven
observed zero-direction units are zero on one side only and are therefore events: the source benchmark has 71 events in 23 claims, the 145-claim subset 65 in 21, External-1 19 in 7 and External-2 95 in 40. Under an alternative binary tie mapping
that assigns a tie to the protective sign, the source count is 63 events in 19 claims; the tie-free
sensitivity (\cref{tab:margin}) avoids dependence on either mapping.

\paragraph{Decisive reversals.} Requiring $|p_t|$ and $|p^{+}_t|$ to exceed a margin removes
indecisive units and events alike. At $\delta=0.2$ only seven events remain; the signed score's
difference from the margin rule moves to $+0.194$ under the primary annotation and to $-0.466$ under
the second, each with an interval too wide to rank methods, so the endpoint amplifies each
annotation's direction without deciding anything on this benchmark.

\paragraph{Margins and quintiles.} \cref{tab:margin} gives the full margin sweep. The lowest
quintile of $|p_t|$ ($|p_t|\le0.31$) holds 54 of the 71 events at a rate of 0.267 \ci{0.164}{0.376};
the next quintiles hold 11, 3, 2 and 1 events at rates 0.054, 0.015, 0.010 and 0.005.
\Cref{tab:classquint} splits each quintile by pre-cutoff majority class: \textsc{protective}-majority
events occur only in the lowest quintile, \textsc{harmful}-majority events in every quintile, and the
class effect holds within strata (crude odds ratio 6.11, Mantel-Haenszel over quintiles 4.35,
adjusted for $|p_t|$ in a claim-clustered class-balanced logistic regression 7.80 \ci{2.01}{33.82}).

\begin{table}[!t]\centering\scriptsize
\setlength{\tabcolsep}{4pt}
\begin{tabular}{@{}lll@{}}
\toprule
$|p_t|$ quintile & \textsc{harm}-majority $n$ / events / rate & \textsc{prot}-majority $n$ / events / rate \\
\midrule
Q1 (lowest) & 59 / 26 / .441 & 137 / 21 / .153 \\
Q2 & 49 / 10 / .204 & 152 / 0 / .000 \\
Q3 & 72 / 3 / .042 & 130 / 0 / .000 \\
Q4 & 68 / 2 / .029 & 134 / 0 / .000 \\
Q5 (highest) & 26 / 1 / .038 & 177 / 0 / .000 \\
\bottomrule
\end{tabular}
\caption{Event rate by pre-cutoff majority class within $|p_t|$ quintiles (1{,}004 units with
$p_t\neq0$). Median $|p_t|$: 0.588 for \textsc{harmful} majorities, 0.612 for \textsc{protective}.}
\label{tab:classquint}
\end{table}

\paragraph{Coverage of selected historical cases.} \Cref{tab:reversals} lists eight claims among
the 165 whose evidence is generally described as having reversed, with the year and primary source
of the trial or guidance usually cited for the reversal, the first year of the archived stream, and
whether the subsequent-stream endpoint of \cref{eq:target} is positive at any cutoff under the
primary annotation. The table describes archive coverage and endpoint values; it supports no
claim about detection timing. The endpoint compares the pooled direction of all records before a
cutoff with that of all records at or after it, so the cutoff is the origin at which the forecast
is posed, not the date at which any change occurs: a positive endpoint at the 1990 cutoff of the
folic acid claim is compatible with a change driven by evidence published in 2007, and no lag can
be read from the difference between a cutoff year and a pivotal year. Seven of the eight claims
have archives that begin before the pivotal year; the beta-carotene archive begins in 2012 because
of the retrieval cap (\cref{app:snapshot}) and is unobservable. An archive that begins before the
pivotal year does not by itself show that an eligible cutoff with adequate pre-pivotal evidence
exists or that the pivotal publication is among the archived records, so the table gives both:
the two hormone-replacement claims and the vitamin E claim have no eligible cutoff before their
pivotal trials and their archives do not contain the trial reports, so no forecast could have
been posed for them; the selenium, folic acid, niacin and dietary cholesterol claims have two to
four eligible cutoffs before their pivotal year, and the three trial-anchored ones contain the
trial reports; and hormone replacement and
cardiovascular disease holds $p_t$ between $-0.04$ and $-0.09$ at every one of its seven cutoffs
with no positive endpoint. A timing claim would require a time-localised change statistic, which this benchmark
does not define.

\begin{table*}[!t]\centering\scriptsize
\setlength{\tabcolsep}{3pt}
\begin{tabular}{@{}p{0.26\textwidth}p{0.17\textwidth}p{0.07\textwidth}p{0.09\textwidth}p{0.10\textwidth}p{0.09\textwidth}p{0.14\textwidth}@{}}
\toprule
Claim & Usually cited source & Year & Archive from & Eligible cutoffs before & Source in archive & Endpoint positive (earliest cutoff) \\
\midrule
hormone replacement $\rightarrow$ cardiovascular disease & WHI trial & 2002 & 1998 & 0 & no & no \\
hormone replacement $\rightarrow$ breast cancer & WHI trial & 2002 & 1998 & 0 & no & yes (2022) \\
beta-carotene $\rightarrow$ cancer & ATBC, CARET trials & 1994, 1996 & 2012 (capped) & 0 & no & unobservable \\
vitamin E $\rightarrow$ cardiovascular disease & HOPE, HPS trials & 2000, 2002 & 1998 & 0 & no & no \\
selenium $\rightarrow$ prostate cancer & SELECT trial & 2009 & 1983 & 2 & yes & no \\
folic acid $\rightarrow$ colorectal cancer & Cole et al.\ trial & 2007 & 1981 & 3 & yes & yes (1990) \\
niacin $\rightarrow$ cardiovascular disease & AIM-HIGH, HPS2-THRIVE trials & 2011, 2014 & 1980 & 4 & yes & no \\
dietary cholesterol $\rightarrow$ cardiovascular disease & 2015 DGAC report & 2015 & 1984 & 4 & n/a & no \\
\bottomrule
\end{tabular}
\caption{Coverage of selected historical cases: the source usually cited for the reversal (WHI:
\citealp{WHI2002}; ATBC: \citealp{ATBC1994}; CARET: \citealp{CARET1996}; HOPE: \citealp{HOPE2000};
HPS: \citealp{HPS2002}; SELECT: \citealp{SELECT2009}; \citealp{Cole2007}; AIM-HIGH:
\citealp{AIMHIGH2011}; HPS2-THRIVE: \citealp{HPS2THRIVE2014}; DGAC: \citealp{DGAC2015}) and its
year, the first year of the archived stream, the number of eligible cutoffs before the pivotal
year, whether the cited source's PMID is among the archived records (the 2015 guidance has no
single PMID), and whether the subsequent-stream endpoint is positive at any cutoff under the
primary annotation (earliest such cutoff in parentheses); eligible cutoffs are the claim's
future-eligible benchmark cutoffs strictly before the pivotal year. Three of the seven observable
cases, and the unobservable one, have no eligible cutoff before the pivotal year, so no forecast
could have been posed for them. The
cutoff is the forecast origin, not the date of change, so no detection lag is defined;
\citet{Prasad2013} is the general reference for reversal as a phenomenon, not the source of these
dates.}
\label{tab:reversals}
\end{table*}

\begin{table}[!t]\centering\scriptsize
\setlength{\tabcolsep}{3.5pt}
\resizebox{\columnwidth}{!}{%
\begin{tabular}{@{}lrrrlrl@{}}
\toprule
Units kept & $n$ & Events & Rate & AUROC [95\% CI] & AUPRC & $1-|p_t|$ \\
\midrule
All future-eligible & 1{,}012 & 71 & 0.070 & 0.946 [0.918, 0.969] & 0.527 & 0.881 \\
$|p_t|\ge0.05$ & 976 & 52 & 0.053 & 0.945 [0.913, 0.969] & 0.456 & 0.855 \\
$|p_t|\ge0.10$ & 949 & 43 & 0.045 & 0.951 [0.919, 0.976] & 0.477 & 0.842 \\
$|p_t|\ge0.15$ & 921 & 36 & 0.039 & 0.949 [0.916, 0.976] & 0.433 & 0.833 \\
$|p_t|\ge0.20$ & 904 & 30 & 0.033 & 0.940 [0.899, 0.974] & 0.364 & 0.811 \\
$|p_t|\ge0.25$ & 860 & 24 & 0.028 & 0.963 [0.934, 0.989] & 0.400 & 0.804 \\
$|p_t|\ge0.30$ & 816 & 17 & 0.021 & 0.959 [0.926, 0.991] & 0.340 & 0.765 \\
\midrule
No zero-direction units & 1{,}001 & 60 & 0.060 & 0.938 [0.907, 0.964] & 0.417 & 0.862 \\
$n^{+}_{\text{signed}}\ge10$ & 1{,}012 & 71 & 0.070 & 0.946 [0.918, 0.969] & 0.527 & 0.881 \\
$n^{+}_{\text{signed}}\ge20$ & 1{,}000 & 71 & 0.071 & 0.945 [0.917, 0.970] & 0.524 & 0.881 \\
\bottomrule
\end{tabular}}
\caption{Target-margin, tie-exclusion and signed-post-evidence robustness. \compact is refitted
leave-one-claim-out inside each retained sample; the last column is the AUROC of the unfitted score
$1-|p_t|$ on the same units. Pre-cutoff margin filtering removes the eight pre-cutoff
zero-direction units but does not necessarily remove the three post-cutoff ties: two zero-direction
units remain at $|p_t|\ge0.05$, one from $|p_t|\ge0.10$ through $|p_t|\ge0.20$, and none at
$|p_t|\ge0.25$. The surviving post-cutoff ties are two cutoffs of one claim (butter and
cardiovascular disease, 2011 with $|p_t|=0.21$ and 2022 with $|p_t|=0.08$), both with a
\textsc{harmful} pre-cutoff sign, so the margin-filtered rows are unchanged under the alternative
\textsc{protective} tie mapping, as is the zero-direction-free row; this equivalence does not
apply to the unfiltered row or to the signed-post-evidence rows, which retain all eleven
zero-direction events (63 events under the alternative mapping). The signed-post-evidence rows
require at least $m$ signed (non-null, non-unclear) post-cutoff records to define the target; the
minimum over the future-eligible sample is 11 and the median 230, so $m=10$ excludes nothing.}
\label{tab:margin}
\end{table}

\paragraph{Zero-direction and signed-future-evidence sensitivities.} Excluding the eleven zero-direction units
and refitting gives 0.938 \ci{0.907}{0.964} and AUPRC 0.417 \ci{0.263}{0.570} on 60 events in 19
claims; the primary model restricted to the same units gives 0.938 (paired $\Delta$ $-0.001$
\ci{-0.007}{0.006}). On the zero-direction-free units \compact stays clearly above the unfitted signals
($+0.076$ \ci{0.030}{0.142} over $1-|p_t|$, $+0.043$ \ci{0.005}{0.089} over low agreement, $+0.054$
\ci{-0.001}{0.114} over recent dissent). Requiring at least 20 signed post-cutoff records to define
the target leaves 1{,}000 units and all 71 events, with 0.945 \ci{0.917}{0.970}: the future target
remains predictable when it must be supported by a minimum number of signed post-cutoff
observations.

\paragraph{Fixed horizons.} Within fixed windows the ranking is 0.908 \ci{0.871}{0.940} at three
years (827 units, 57 events in 26 claims), 0.916 \ci{0.874}{0.948} at five years (783 units, 51
events) and 0.916 \ci{0.870}{0.955} at ten years (526 units, 39 events), with AUPRC 0.391, 0.409
and 0.479. These rows require at least ten \emph{non-null} records in the window. The horizon rows
of \cref{app:cumulative} (821/56, 777/50 and 526/39 units/events; 0.908, 0.930 and 0.916)
instead require at least ten \emph{signed} records in the window, which is why the three- and
five-year cohorts and their AUROCs differ between the two tables; they are distinct eligibility
rules, not the same analysis.

\section{Cumulative and Three-State Endpoints}
\label{app:cumulative}
With $p$ the mean of the $-1/0/+1$ codes over signed records and $\sgn(0)=0$, the primary endpoint
is $R_{\mathrm{sub}}=\mathbf{1}[\sgn(p_{\mathrm{post}})\neq\sgn(p_{\mathrm{pre}})]$, the cumulative
endpoint $R_{\mathrm{cum}}=\mathbf{1}[\sgn(p_{\mathrm{cum}})\neq\sgn(p_{\mathrm{pre}})]$ with
$p_{\mathrm{cum}}$ over all records of the claim, and the three-state endpoint counts a change of
the state $\{-1,0,+1\}$ with $0$ for $|p|<\delta$, for the subsequent and the cumulative
comparison. Horizon variants use the window $[t,t+H)$ with at least ten signed records in the
window and a claim last year of at least $t+H-1$. \cref{tab:astra3} gives the endpoint variants on the source future-eligible population, with
the fixed-horizon rows subject to their stated horizon-specific eligibility, and the frozen models
refitted leave-one-claim-out. The event sets
nest: 34 units are events under both $R_{\mathrm{sub}}$ and $R_{\mathrm{cum}}$, 37 under
$R_{\mathrm{sub}}$ only (in twelve claims, e.g.\ \texttt{eggs\_cvd\_events} at 1988: $p_{\mathrm{pre}}=0.87$ on 30
signed records, $p_{\mathrm{post}}=-0.01$ on 539, $p_{\mathrm{cum}}=0.04$), and none under
$R_{\mathrm{cum}}$ only. The cumulative endpoint is rarer; on it \compact gives 0.921 \ci{0.862}{0.984} (0.9205 before
rounding), the two-variable score 0.941 and the unfitted margin rule 0.932, so the untrained rule
exceeds \compact's point estimate and no score ``improves'' on the cumulative endpoint in an
established sense. On the cumulative fixed-horizon variants the margin rule alone reaches 0.992,
0.991 and 0.971 (14, 17 and 20 events): some endpoint constructions are largely predictable from
the initial margin, and these are rare events on different samples, not learned-model
performance. The three-state endpoint at $\delta=0.20$ counts many small moves into and out of
the signed-balance band and every score falls to about 0.90 on it. The horizon rows here require
at least ten signed records in the window and a claim last year of at least $t+H-1$, a different
eligibility rule from the fixed-horizon rows of \cref{app:robust} (ten non-null records),
which is why the cohorts differ.

\begin{table*}[!t]\centering\scriptsize
\setlength{\tabcolsep}{3pt}
\resizebox{\textwidth}{!}{%
\begin{tabular}{@{}lrrrrrrrrrr@{}}
\toprule
Endpoint & Units & Events & Pos.\ claims & Prev. & Compact-7 AUROC & [95\% CI] & AUPRC & 2-var AUROC & $1-|p_t|$ & Claim AUROC \\
\midrule
R\_sub (primary) & 1012 & 71 & 23 & .070 & .946 & [0.918,\,0.969] & .527 & .948 & .881 & .974 \\
R\_cum & 1012 & 34 & 15 & .034 & .921 & [0.862,\,0.984] & .424 & .941 & .932 & .976 \\
R3\_sub delta=0.10 & 1012 & 100 & 29 & .099 & .952 & [0.925,\,0.973] & .597 & .946 & .889 & .980 \\
R3\_cum delta=0.10 & 1012 & 58 & 19 & .057 & .944 & [0.908,\,0.974] & .479 & .946 & .881 & .966 \\
R3\_sub delta=0.20 & 1012 & 117 & 37 & .116 & .896 & [0.857,\,0.929] & .438 & .892 & .854 & .902 \\
R3\_cum delta=0.20 & 1012 & 72 & 24 & .071 & .902 & [0.850,\,0.946] & .383 & .909 & .853 & .923 \\
R\_sub 3y & 821 & 56 & 26 & .068 & .907 & [0.871,\,0.940] & .390 & .913 & .913 & .943 \\
R\_cum 3y & 821 & 14 & 12 & .017 & .974 & [0.944,\,0.996] & .436 & .974 & .992 & .966 \\
R\_sub 5y & 777 & 50 & 24 & .064 & .930 & [0.897,\,0.957] & .421 & .929 & .908 & .948 \\
R\_cum 5y & 777 & 17 & 14 & .022 & .939 & [0.873,\,0.981] & .239 & .967 & .991 & .967 \\
R\_sub 10y & 526 & 39 & 19 & .074 & .916 & [0.870,\,0.955] & .479 & .934 & .899 & .948 \\
R\_cum 10y & 526 & 20 & 13 & .038 & .937 & [0.882,\,0.982] & .431 & .956 & .971 & .980 \\
\bottomrule
\end{tabular}}
\caption{Subsequent-evidence endpoint ($R_{\mathrm{sub}}$, the primary target), cumulative updated-direction endpoint ($R_{\mathrm{cum}}$), three-state endpoints with a no-predominant-direction band $|p|<\delta$, and fixed-horizon variants. Non-horizon variants use the 1{,}012 future-eligible source units; fixed-horizon rows use their stated horizon-specific eligibility. Models are fitted leave-one-claim-out within each evaluated sample.}
\label{tab:astra3}
\end{table*}

\section{Numeric Effect-Reversal Endpoint: A Negative Scope Result}
\label{app:numeric}
Defining an event as a sign change of the DerSimonian-Laird pooled log-ratio effect between the
pre-cutoff estimates and those inside a ten-year window, which uses no direction label, \compact
reaches AUROC 0.426 \ci{0.343}{0.512} pooled over all domains (587 units, 90 events), against 0.643
\ci{0.562}{0.718} for the unfitted magnitude $-|\mu_{\text{pre}}|$ ($\Delta$ $-0.217$
\ci{-0.356}{-0.071}); the hypothesised verdict is not supported in any domain. Only 7.5\% of
records carry a parseable ratio with a confidence interval, so the endpoint rests on 587 units.
\compact does not predict generic pooled-effect reversal; it predicts later change of the
categorical direction, which is what it was built for, and this falsification boundary is retained
deliberately.

\section{Agreement Definition Audit and Sensitivity}
\label{app:agreement}
The shipped agreement $a$ is computed by \texttt{agg\_verdict}: with $n_N$ \textsc{null} records
among $n$ resolved records, $a=n_N/n$ when $n_N/n\ge0.5$ (the \textsc{null}-dominant branch) and
otherwise the majority share among the signed records, which equals $(1+|p_t|)/2$. Rebuilding the
direction block from the record streams reproduces every shipped direction column
exactly (maximum absolute difference 0). Of the 1{,}078 units, 1{,}035 are on the signed
branch, where the identity holds on every unit (Pearson 1.000 with $|p_t|$), 42 are on the
\textsc{null} branch, where $a$ equals the \textsc{null} share, and one has fewer than eight signed
records and $a=0$; 22 units have a \textsc{null} share in $[0.45,0.50)$ and 19 in $[0.50,0.55)$.
Ten pairs of consecutive cutoffs of one claim cross the branch, with a mean absolute jump in $a$ of
0.168 (0.046 under definition B); for example \texttt{late\_night\_eating\_obesity} moves from
$a=0.50$ at 2013 to $0.857$ at 2015 as its \textsc{null} share moves from 0.50 to 0.44. Over all
units the correlation of $a$ with $|p_t|$ is 0.944 (Spearman 0.959).

Four definitions were then substituted upstream of both models, with everything
else rebuilt: A the shipped implementation; B the maximum share among
\textsc{protective}/\textsc{harmful}/\textsc{null}; C the signed-only share $(1+|p_t|)/2$ with no
branch; D one minus the normalised entropy of the three-way distribution; and C+null, which adds the
\textsc{null} share as an explicit third input of the two-variable model. \cref{tab:astra2}
reports the results with paired claim-clustered differences under the primary annotation. B and D
are clearly worse; C is marginally worse ($-0.023$, upper limit 0.002) and C+null recovers most of
the gap ($-0.015$, interval covering zero). That primary-annotation comparison is not a
justification for A across annotations: under the second annotation the substitution of C for A
repairs the shipped score's loss ($0.758$ to $0.940$ on common support; \cref{tab:branchfree}),
and on the primary annotation it removes the established gain over the margin rule. A is retained
as the archived reference specification because every protocol cell, control and re-annotation was
computed on it; C is the matched diagnostic intervention, reported alongside A and never merged
with it into one method row, and C+null is the recommended agreement input for a subsequent
version. On histories with fewer than eight signed records, where $p_t$ is gated to 0, definition
A returns $a=0$ (the archived fallback); the rebuild script for definition C applies the same
gate and returns $a=0$ on those units rather than the formula value $(1+0)/2=0.5$ (verified
against the frozen builder with no mismatches), so that the two definitions differ only on the
\textsc{null}-dominant branch (one unit under the primary annotation, ten under the second,
fourteen under intersection labels are affected by the gate).
The manuscript describes A's two-branch form and its dependence on $|p_t|$ exactly instead of
treating direction and agreement as independent dimensions.
\begin{table*}[!t]\centering\scriptsize
\setlength{\tabcolsep}{3pt}
\resizebox{\textwidth}{!}{%
\begin{tabular}{@{}lrrrrrr@{}}
\toprule
Def. & 2-var AUROC & [95\% CI] & AUPRC & $\Delta$ vs A & [95\% CI] & Compact-7 AUROC \\
\midrule
A & .948 & [0.921,\,0.971] & .535 &  &  & .946 \\
B & .890 & [0.830,\,0.939] & .387 & -0.058 & [-0.114,\,-0.010] & .899 \\
C & .925 & [0.890,\,0.959] & .500 & -0.023 & [-0.050,\,0.002] & .925 \\
D & .849 & [0.758,\,0.922] & .349 & -0.099 & [-0.191,\,-0.026] & .858 \\
C+null & .932 & [0.901,\,0.960] & .477 & -0.015 & [-0.036,\,0.006] & \\
p\_t only & .766 & [0.696,\,0.826] & .128 &  &  & \\
1-|p\_t| & .881 & [0.824,\,0.936] & .468 &  &  & \\
\bottomrule
\end{tabular}}
\caption{Agreement-definition sensitivity on the 1,012 future-eligible units. A: shipped implementation; B: max share of \textsc{protective}/\textsc{harmful}/\textsc{null}; C: signed-only majority share $(1+|p_t|)/2$; D: normalised-entropy consensus; C+null adds the \textsc{null} share as a third input. Paired differences on identical claim resamples.}
\label{tab:astra2}
\end{table*}

\section{Semantic-Orientation Audit}
\label{app:orientation}
Deterministic rules over the claim identifier, exposure term and outcome term assign each of the
165 claims to zero or more orientation categories: deficiency exposure (12 claims; \texttt{deficien},
\texttt{malnutrition}, \texttt{intolerance}); therapeutic or preventive use (24; prevention,
treatment, therapy, remission, symptom, response, replacement, probiotic, multivitamin and similar
terms); restriction or inversion (15; skipping, replacing, low-fat, energy deficit, restricted,
fasting, late-night, alternate-day, low intake, free diet); benefit-valence outcome (27; weight
loss, cognitive function, remission, growth, diversity, fertility, bioavailability, absorption,
response, control, prevention, neurodevelopment, metabolism, haemoglobin); and non-disease outcome
term (19; duration, intake, risk, progression, and the benefit-valence terms). The rules are the
whole method; no claim is re-labelled by hand, and the per-claim categories, reasons and affected
record counts are released. The target is invariant to flipping the sign of every label of a claim, so canonicalising a
claim's orientation cannot change any event; the predictors are not invariant, because a score
with a signed $p$ term changes under $p\to-p$ (on the signed branch a logit
$b_0+b_1 p+b_2|p|$ changes by $-2b_1 p$), so claim-specific recoding can alter
predictions and fitted results. A mixture of orientations within a claim cannot be resolved
without reading records. The sensitivity
is therefore exclusion: the core tier removes the 53 claims in the four categories named by the
review (deficiency, therapeutic or preventive, restriction or inversion, non-disease outcome term),
and the strict tier also removes benefit-valence outcomes (63 claims). \cref{tab:astra4} shows
that removing a third to two-fifths of the benchmark leaves the benchmark signal unchanged; the events lost with the flagged claims (20 and 26 of 71) are not the source of the
result. The retained claims are lexically filtered, not an adjudicated orientation-clean cohort.

\begin{table*}[!t]\centering\scriptsize
\setlength{\tabcolsep}{3pt}
\resizebox{\textwidth}{!}{%
\begin{tabular}{@{}lrrrrrrr@{}}
\toprule
Benchmark & Units & Claims & Events & Compact-7 AUROC & [95\% CI] & 2-var AUROC & Claim AUROC \\
\midrule
full benchmark & 1012 & 160 & 71 & .946 & [0.918,\,0.969] & .948 & .974 \\
orientation-clean, core & 681 & 107 & 51 & .945 & [0.909,\,0.977] & .953 & .985 \\
orientation-clean, strict & 618 & 97 & 45 & .952 & [0.915,\,0.984] & .964 & .990 \\
\bottomrule
\end{tabular}}
\caption{Semantic-orientation sensitivity: the benchmark with every claim flagged by the deterministic orientation rules removed (core: deficiency, therapeutic or preventive, restriction or inversion, non-disease outcome term; strict: core plus benefit-valence outcomes), models refitted leave-one-claim-out inside the retained claims.}
\label{tab:astra4}
\end{table*}

\section{Complete Dependence Analyses}
\label{app:dependence}

\paragraph{Cutoff designs and claim-level scoring.} \cref{tab:designs} gives every cutoff
design, scored both with the all-cutoff model's out-of-fold predictions and after refitting inside
the design, together with claim-level scoring and the publication-disjoint protocol. These
sensitivities reuse the same events and are not independent confirmations.

\begin{table}[!t]\centering\scriptsize
\setlength{\tabcolsep}{3pt}
\resizebox{\columnwidth}{!}{%
\begin{tabular}{@{}lrrll@{}}
\toprule
Design & Units & Events & AUROC, all-cutoff model & AUROC, refit \\
\midrule
All cutoffs & 1{,}012 & 71 & 0.946 [0.918, 0.969] & -- \\
Claim level, maximum score over cutoffs & 160 & 23 & 0.974 [0.950, 0.993] & -- \\
PMID-overlap-disjoint leave-one-claim-out & 1{,}012 & 71 & -- & 0.946 [0.918, 0.969] \\
One random cutoff per claim$^{\ast}$ & 160 & 10.8 & 0.945 [0.908, 0.976] & 0.921 [0.850, 0.966] \\
Earliest eligible cutoff & 160 & 17 & 0.965 [0.933, 0.990] & 0.943 [0.896, 0.982] \\
Middle eligible cutoff & 160 & 11 & 0.961 [0.913, 0.995] & 0.932 [0.870, 0.983] \\
Latest eligible cutoff & 160 & 12 & 0.971 [0.941, 0.992] & 0.964 [0.925, 0.990] \\
Equalised cutoff count$^{\ast}$ & 160 & 10.7 & 0.944 [0.906, 0.971] & 0.916 [0.845, 0.964] \\
Seven cutoffs, claims with fewer dropped & 770 & 59 & 0.940 & 0.937 \\
Cutoff positions 1-2 & 318 & 25 & 0.938 [0.908, 0.965] & -- \\
Cutoff positions 3+ & 694 & 46 & 0.949 [0.918, 0.976] & -- \\
\bottomrule
\end{tabular}}
\caption{Cutoff designs, claim-level scoring and the publication-disjoint protocol.
$^{\ast}$Mean over 500 (random) or 200 (equalised) draws; the bracket is the 2.5 to 97.5 percentile
range across draws, a design-variability range rather than a confidence interval. Other brackets
are claim-clustered 95\% intervals. Claim-level AUPRC is 0.797 \ci{0.617}{0.957}; the
publication-disjoint AUPRC is 0.525 \ci{0.377}{0.666}. Leave-one-positive-claim-out refits change
AUROC by at most 0.0052 and AUPRC by at most 0.054.}
\label{tab:designs}
\end{table}

\paragraph{Positive-claim influence and event concentration.} Dropping each of the 23 positive
claims in turn and refitting changes AUROC by between $-0.0037$ (eggs $\rightarrow$ cardiovascular
events, 6 events) and $+0.0052$ (vitamin B12 $\rightarrow$ cancer, 7 events) and AUPRC by at most
0.054; the four heaviest claims hold 35\% of the 71 events and eleven claims contribute three or
more.

\paragraph{PMID-overlap-disjoint evaluation.} Of the 96{,}091 distinct identifiers, 12{,}411
(12.9\%) appear in two or more claims' histories, and 25.6\% of the 112{,}453 associations use such
a shared identifier (2.3 claims per shared identifier on average). Leave-one-claim-out removes the
held-out claim but not its publications from the training claims. The disjoint protocol therefore,
for each held-out claim, removes every identifier of its labelled history from every other claim,
rebuilds the \compact features, the target and the future eligibility of every training unit
from the reduced histories (a rebuild that reproduces the benchmark exactly when nothing is
removed), refits \compact, and scores the held-out claim's unchanged units. The rebuilt training sets lose on average 0.25\% of their associations (at most 2.3\%) and 0.6
units per fold; the disjoint protocol gives AUROC 0.946 \ci{0.918}{0.969} and AUPRC 0.525
\ci{0.377}{0.666}, a paired change of $-0.0001$ \ci{-0.001}{0.001} and $-0.002$ \ci{-0.007}{0.002}
against standard leave-one-claim-out. The per-fold removal is small by construction: each held-out
claim's shared publications are a small fraction of the other 164 claims' histories. The claim
graph in which two claims are linked by any shared identifier has a single connected component
containing all 165 claims, so a leave-one-component-out split is degenerate and is not reported;
instead, a stricter unique-publication check removes every one of the 12{,}411 shared identifiers
(25.6\% of all associations) from every claim, rebuilds features, targets and eligibility for all
units, and runs standard leave-one-claim-out on what remains: 880 future-eligible units in 150 claims
with 57 events in 20 claims, AUROC 0.925 \ci{0.886}{0.960} and AUPRC 0.384 \ci{0.255}{0.621}. The
signal therefore does not depend on publications that the claims share. Repeated estimates from the same cohorts published under different identifiers
remain a possible dependence that no identifier-level protocol removes.

\section{Full Feature and Ablation Results}
\label{app:ablation}
\paragraph{Learned models.} \cref{tab:learnedfull} gives every learned model on the same
units, folds and resamples, together with the 12 rules scored for later change
(\cref{tab:rules}). On the 145-claim subset the same picture holds (\compact 0.950;
random forest 0.944, elastic net 0.944, gradient boosting 0.933, GAM 0.934).
\label{app:learned}

\begin{sidewaystable*}\centering\scriptsize
\setlength{\tabcolsep}{3.5pt}
\resizebox{\textheight}{!}{%
\begin{tabular}{@{}llrlrlrrlr@{}}
\toprule
Family & Method (inputs) & AUROC & 95\% CI & AUPRC & 95\% CI & F1 & Claim AUROC & $\Delta$AUROC [95\% CI] & $\Delta$AUPRC [95\% CI] \\
\midrule
Ours & \textbf{\compact (7)} & \textbf{0.946} & [0.918, 0.969] & 0.527 & [0.378, 0.669] & 0.494 & \textbf{0.974} & -- & -- \\
\midrule
Same 7 inputs & Generalised additive model (7) & 0.925 & [0.880, 0.964] & 0.491 & [0.329, 0.645] & 0.460 & 0.967 & $+0.020$ [0.000, 0.050] & $+0.035$ [$-$0.019, 0.111] \\
 & Gradient boosting (7) & 0.911 & [0.861, 0.957] & 0.447 & [0.311, 0.612] & 0.483 & 0.974 & $+0.035$ [0.009, 0.066] & $+0.080$ [0.017, 0.134] \\
 & Random forest (7) & 0.906 & [0.865, 0.946] & 0.439 & [0.318, 0.581] & 0.413 & 0.969 & $+0.040$ [0.015, 0.063] & $+0.088$ [0.008, 0.170] \\
\midrule
All pre-cutoff variables & Random forest & 0.938 & [0.908, 0.965] & \textbf{0.545} & [0.394, 0.688] & 0.545 & 0.969 & $+0.007$ [$-$0.011, 0.027] & $-0.018$ [$-$0.121, 0.089] \\
 & Elastic net & 0.938 & [0.909, 0.963] & 0.450 & [0.304, 0.605] & 0.517 & 0.932 & $+0.008$ [$-$0.013, 0.034] & $+0.077$ [$-$0.032, 0.198] \\
 & Gradient boosting (XGBoost) & 0.932 & [0.888, 0.967] & 0.543 & [0.359, 0.719] & \textbf{0.578} & 0.971 & $+0.014$ [$-$0.013, 0.047] & $-0.016$ [$-$0.125, 0.101] \\
 & SVM, RBF kernel & 0.925 & [0.894, 0.955] & 0.458 & [0.287, 0.625] & 0.365 & 0.941 & $+0.021$ [$-$0.005, 0.051] & $+0.069$ [$-$0.050, 0.217] \\
 & $k$-nearest neighbours & 0.909 & [0.858, 0.954] & 0.382 & [0.251, 0.546] & 0.497 & 0.929 & $+0.037$ [$-$0.004, 0.083] & $+0.145$ [0.009, 0.270] \\
 & Best single variable & 0.901 & [0.849, 0.948] & 0.509 & [0.320, 0.668] & 0.447 & 0.947 & $+0.044$ [0.003, 0.087] & $+0.017$ [$-$0.086, 0.149] \\
 & Logistic regression, unpenalised & 0.897 & [0.847, 0.945] & 0.373 & [0.231, 0.555] & 0.417 & 0.929 & $+0.048$ [0.005, 0.100] & $+0.153$ [0.025, 0.266] \\
 & Balanced logistic (the \compact estimator, all inputs) & 0.893 & [0.839, 0.944] & 0.374 & [0.235, 0.549] & 0.443 & 0.929 & $+0.053$ [0.009, 0.108] & $+0.153$ [0.036, 0.262] \\
 & Neural network (MLP, 32 units) & 0.889 & [0.841, 0.936] & 0.360 & [0.220, 0.549] & 0.397 & 0.928 & $+0.057$ [0.014, 0.105] & $+0.167$ [0.009, 0.302] \\
 & Gaussian naive Bayes & 0.778 & [0.671, 0.877] & 0.289 & [0.170, 0.464] & 0.347 & 0.839 & $+0.167$ [0.071, 0.272] & $+0.237$ [0.082, 0.372] \\
 & Decision tree & 0.674 & [0.582, 0.775] & 0.315 & [0.214, 0.446] & 0.507 & 0.954 & $+0.272$ [0.167, 0.369] & $+0.211$ [0.055, 0.347] \\
\midrule
Nested selection & Forward selection over all variables (median 13) & 0.935 & [0.902, 0.962] & 0.481 & [0.323, 0.648] & -- & -- & $+0.011$ [$-$0.011, 0.035] & $+0.046$ [$-$0.069, 0.156] \\
 & Backward elimination from all variables (median 12) & 0.923 & [0.883, 0.961] & 0.485 & [0.340, 0.638] & -- & -- & $+0.023$ [$-$0.012, 0.065] & $+0.041$ [$-$0.052, 0.139] \\
\bottomrule
\end{tabular}}
\caption{Every learned model for later change on the 1{,}012 future-eligible units (71 events),
leave-one-claim-out with nested tuning (\cref{tab:appgrids}), claim-clustered 95\% intervals
from the same 2{,}000 resamples. $\Delta$ is \compact minus the model on identical resamples. Bold
marks the best value in a column among the learned models listed; the two-variable score
(\cref{tab:rq3}) has the highest AUROC point estimate of any fitted score, 0.948. A paired
interval that covers zero is not a difference. Random forest and gradient boosting on the full
representation have the higher AUPRC point estimates and gradient boosting the higher F1; they are
retained as valid competitive baselines. Nested selection
chooses the variable subset inside each outer fold.}
\label{tab:learnedfull}
\end{sidewaystable*}

\begin{sidewaystable*}\centering\scriptsize
\setlength{\tabcolsep}{3.5pt}
\resizebox{\textheight}{!}{%
\begin{tabular}{@{}llrlrlrrrl@{}}
\toprule
Family & Method & AUROC & 95\% CI & AUPRC & 95\% CI & F1 & Prec.$_{\text{m}}$ & Claim AUROC & $\Delta$AUROC [95\% CI] \\
\midrule
Reference & \compact & \textbf{0.946} & [0.918, 0.969] & \textbf{0.527} & [0.378, 0.669] & 0.494 & 0.507 & \textbf{0.974} & -- \\
\midrule
Rule & Cumulative-agreement vote & 0.859 & [0.793, 0.920] & 0.285 & [0.177, 0.403] & 0.422 & 0.408 & 0.924 & $+0.086$ [0.034, 0.141] \\
 & Conflict-aware rule & 0.846 & [0.786, 0.900] & 0.241 & [0.149, 0.346] & 0.198 & 0.325 & 0.939 & $+0.100$ [0.051, 0.150] \\
 & Replacement-time retrieval & 0.808 & [0.754, 0.861] & 0.186 & [0.114, 0.265] & 0.330 & 0.203 & 0.803 & $+0.138$ [0.088, 0.188] \\
 & Synthesis-presence workflow & 0.741 & [0.707, 0.775] & 0.128 & [0.076, 0.183] & 0.227 & 0.128 & 0.668 & $+0.205$ [0.167, 0.242] \\
 & Decline-triggered workflow & 0.730 & [0.672, 0.784] & 0.148 & [0.091, 0.213] & 0.280 & 0.036 & 0.612 & $+0.216$ [0.168, 0.265] \\
 & Scoping-review workflow & 0.704 & [0.627, 0.767] & 0.118 & [0.068, 0.173] & 0.219 & 0.126 & 0.650 & $+0.242$ [0.184, 0.312] \\
 & Recency only & 0.644 & [0.559, 0.720] & 0.105 & [0.054, 0.168] & 0.214 & 0.131 & 0.665 & $+0.301$ [0.221, 0.390] \\
 & Conclusiveness proxy (TSA-style) & 0.620 & [0.527, 0.714] & 0.102 & [0.056, 0.161] & 0.216 & 0.144 & 0.629 & $+0.326$ [0.234, 0.416] \\
 & Evidence-grade workflow & 0.582 & [0.558, 0.607] & 0.083 & [0.049, 0.120] & 0.153 & 0.083 & 0.518 & $+0.363$ [0.330, 0.396] \\
 & Volume threshold & 0.565 & [0.472, 0.665] & 0.081 & [0.046, 0.127] & 0.156 & 0.091 & 0.544 & $+0.381$ [0.286, 0.472] \\
 & Living-review workflow & 0.564 & [0.513, 0.604] & 0.080 & [0.046, 0.117] & 0.148 & 0.081 & 0.500 & $+0.382$ [0.343, 0.427] \\
 & Pooled-estimate workflow & 0.427 & [0.344, 0.527] & 0.064 & [0.037, 0.096] & 0.133 & 0.047 & 0.502 & $+0.518$ [0.410, 0.611] \\
\midrule
Signal & Low agreement & 0.908 & [0.858, 0.950] & 0.447 & [0.299, 0.599] & 0.478 & 0.423 & 0.945 & $+0.038$ [0.003, 0.078] \\
 & Recent dissent & 0.883 & [0.831, 0.936] & 0.523 & [0.377, 0.676] & \textbf{0.566} & \textbf{0.578} & 0.973 & $+0.063$ [0.006, 0.125] \\
 & $1-|p_t|$ (direction margin) & 0.881 & [0.824, 0.936] & 0.468 & [0.359, 0.589] & 0.398 & 0.423 & 0.971 & $+0.064$ [0.011, 0.122] \\
 & Absent consensus & 0.786 & [0.702, 0.864] & 0.259 & [0.151, 0.365] & 0.463 & 0.356 & 0.880 & $+0.160$ [0.086, 0.238] \\
 & Low conclusiveness & 0.632 & [0.534, 0.733] & 0.112 & [0.060, 0.184] & 0.229 & 0.160 & 0.680 & $+0.313$ [0.217, 0.410] \\
 & Low evidence volume & 0.551 & [0.432, 0.677] & 0.100 & [0.056, 0.175] & 0.164 & 0.155 & 0.511 & $+0.394$ [0.271, 0.509] \\
 & Decline signal & 0.504 & [0.443, 0.561] & 0.073 & [0.043, 0.112] & 0.147 & 0.042 & 0.496 & $+0.442$ [0.377, 0.511] \\
 & Heterogeneity $I^2$ & 0.385 & [0.285, 0.491] & 0.067 & [0.038, 0.107] & 0.131 & 0.042 & 0.617 & $+0.561$ [0.443, 0.670] \\
 & Imprecision & 0.338 & [0.246, 0.429] & 0.053 & [0.031, 0.081] & 0.131 & 0.014 & 0.448 & $+0.608$ [0.507, 0.709] \\
 & Recent accrual share & 0.330 & [0.241, 0.433] & 0.049 & [0.030, 0.076] & 0.127 & 0.023 & 0.448 & $+0.616$ [0.506, 0.710] \\
 & Agreement contraction & 0.318 & [0.236, 0.411] & 0.052 & [0.031, 0.093] & 0.131 & 0.042 & 0.598 & $+0.628$ [0.536, 0.710] \\
\midrule
Trivial & Event prevalence (constant) & 0.500 & -- & 0.070 & -- & -- & 0.070 & 0.500 & -- \\
\bottomrule
\end{tabular}}
\caption{The 12 rule comparators (scored \stable $<$ \forming $<$ \unstable) and the single oriented signals, for later change on the 1{,}012 future-eligible units
with the same folds, threshold rule and claim resamples as \compact. F1 is taken at the
training-claim threshold; Prec.$_{\text{m}}$ is the precision when as many units are flagged as
there are events; claim AUROC scores a claim by its maximum over cutoffs.}
\label{tab:rules}
\end{sidewaystable*}

\paragraph{Within \compact.} The standardised coefficients of the fit on all 1{,}012 source units are agreement
$-3.32$, pooled direction $+2.97$, $n_{\text{pre}}$ $+1.32$, $n_{\text{directional}}$ $-1.07$,
$n_{RR}$ $-0.70$, $\sigma_{\log RR}$ $-0.06$ and $\mu_{\log RR}$ $+0.05$. Agreement alone ranks later
change at 0.901 and pooled direction alone at 0.766; the two together reach 0.948 \ci{0.921}{0.971},
no clear difference from the seven ($+0.002$ \ci{-0.006}{0.011}); the volume pair alone gives 0.603
and the effect-estimate triple 0.635. Removing $p_t$ from \compact costs 0.045 and leaves 0.900 (\cref{tab:rq3}); because the
count and effect inputs are correlated with the direction inputs, coefficient magnitudes are not a
reliable attribution of their value.

\paragraph{Adding blocks and nested selection.} The eleven direction variables alone match \compact
(0.947); adding the temporal-accrual block changes nothing (0.946) and adding the decline block
lowers it (0.916). None of the other 48 variables gives a clear gain when added to the seven; the
best, agreement contraction, adds $+0.003$ \ci{-0.008}{0.014}. Greedy selection over the full
representation, nested inside the folds, reaches 0.935 \ci{0.902}{0.962} forward (median 13
variables, chosen in 99\% of folds: null share, pooled direction and agreement) and 0.923
\ci{0.883}{0.961} backward, neither clearly different from the seven-variable set.

\begin{table}[!t]\centering\scriptsize
\setlength{\tabcolsep}{3pt}
\resizebox{\columnwidth}{!}{%
\begin{tabular}{@{}lrrr@{}}
\toprule
Predictor & AUROC & AUPRC & $\Delta$AUROC vs \compact [95\% CI] \\
\midrule
\textbf{\compact (7)} & \textbf{0.946} & \textbf{0.527} & -- \\
Pooled direction + agreement (2) & 0.948 & 0.535 & $+0.002$ [$-$0.006, 0.011] \\
Direction block alone (11) & 0.947 & 0.533 & $+0.001$ \\
\compact + temporal accrual (11) & 0.946 & 0.526 & $+0.000$ \\
\compact + decline (11) & 0.916 & 0.494 & $-0.030$ \\
Agreement alone (1) & 0.901 & 0.436 & $-0.045$ [$-$0.093, $-$0.006] \\
\compact without $p_t$ (6) & 0.900 & 0.429 & $-0.045$ \\
Temporal accrual alone (4) & 0.898 & 0.450 & $-0.048$ \\
All 55 variables, balanced logistic & 0.893 & 0.374 & $-0.053$ [$-$0.108, $-$0.009] \\
Pooled direction alone (1) & 0.766 & 0.128 & $-0.179$ [$-$0.249, $-$0.122] \\
Decline block alone (4) & 0.537 & 0.084 & clearly lower \\
Heterogeneity block alone (6) & 0.486 & 0.077 & clearly lower \\
\bottomrule
\end{tabular}}
\caption{Which pre-cutoff signals carry the prediction, on the 1{,}012 future-eligible units,
leave-one-claim-out. Fitted predictors are class-balanced logistic regressions on the named
variables; $\Delta$ is the predictor minus \compact on identical claim resamples.}
\label{tab:rq3}
\end{table}

\section{Heterogeneity Analyses}
\label{app:hetero}
Stratified by pre-cutoff $I^2$ (units with fewer than three effect estimates, where $I^2$ is not
estimable, form their own stratum), \compact ranks later change at 0.943 \ci{0.899}{0.981} where
$I^2$ is not estimable (167 future-eligible units, 25 events), 0.864 \ci{0.684}{1.000} at $I^2\le0.50$
(23 units, one event), 0.964 \ci{0.911}{0.996} at $0.50<I^2\le0.90$ (222 units, 12 events) and
0.937 \ci{0.901}{0.971} at $I^2>0.90$ (600 units, 33 events); the later-change rate is highest,
0.150, where $I^2$ cannot be estimated. Three
statements should be kept apart. Incremental benefit: the six heterogeneity variables change
\compact by $-0.004$ \ci{-0.013}{0.002} AUROC, so they add nothing measurable conditional on the
reference inputs. Model-specific performance: fitted on their own they rank later change at 0.486
\ci{0.339}{0.642}. Association: $I^2$ alone gives 0.385 when oriented so that higher heterogeneity
means more expected change, which is 0.615 under the reversed orientation and is therefore
evidence of an inverse association in this benchmark, not of zero information; the refitted
coefficient of $I^2$ is consistently negative in every domain (\cref{app:crossdomain}).
This analysis does not establish that heterogeneity carries no information in general.

\section{Rule Comparators, Single Signals and Learned-Model Tuning}
\label{app:comparators}
\cref{tab:appops} gives the operational definition of every rule comparator and single signal
and \cref{tab:appgrids} the tuning grids of the learned models. Every rule reads the common
pre-cutoff representation. The rules are simple threshold policies \emph{inspired by} the named
evidence-assessment principles (cumulative voting, conclusiveness monitoring, living-review
triggers, evidence grading, pooled estimation); they are not implementations or evaluations of
GRADE, Trial Sequential Analysis, living reviews or meta-analysis, and beating them substantiates
no comparison with those methodologies. Thresholding a continuous signal into a few states also
discards ranking resolution, so the continuous single signals are the more informative
comparators. For ranking, a rule's state is scored \stable
$=0$, \forming (or \unassess) $=1$, \unstable $=2$. Learned models are tuned by nested five-fold
GroupKFold over the training claims of every outer fold; the setting with the highest AUROC on the
pooled inner out-of-fold scores wins, with ties to the most regularised setting, and is refitted on
the training claims before the held-out claim is scored. Hard decisions for every method use the
F1-maximising threshold on the training claims of the fold.

\begin{table*}[!t]\centering\scriptsize
\setlength{\tabcolsep}{4pt}
\begin{tabular}{@{}p{0.19\textwidth}p{0.22\textwidth}p{0.55\textwidth}@{}}
\toprule
Comparator & Inputs & Operational rule \\
\midrule
Cumulative-agreement vote & agreement & \stable if agreement $\ge0.70$; \unstable if agreement $<0.55$; else \forming \\
Volume threshold & n\_pre\_t & \stable if at least 150 pre-cutoff records; else \forming \\
Recency only & recency\_density & \stable if at least 35\% of pre-cutoff records fall in the last five years; else \forming \\
Conclusiveness proxy (TSA-style) & naive\_tsa & \stable if accumulated precision reaches half the information-size proxy (naive\_tsa $\ge0.5$); else \forming \\
Conflict-aware rule & i\_squared, dissent\_recent, agreement & \unstable if $I^2>0.9$ and recent dissent $>0.4$; \stable if agreement $\ge0.7$; else \forming \\
Replacement-time retrieval & agreement\_recent & \stable if recent-window agreement $\ge0.7$; else \forming \\
Synthesis-presence workflow & synth\_frac, agreement & \stable if syntheses are present ($\ge3\%$ of records) and agreement $\ge0.65$; else \forming \\
Evidence-grade workflow & n\_pre\_t, agreement, rct\_frac & \stable if at least 100 records, agreement $\ge0.7$ and trial share $\ge0.2$; else \forming \\
Living-review workflow & vol\_growth, agreement & \forming while evidence volume is growing (vol\_growth $>1.2$); else \stable if agreement $\ge0.7$, else \forming \\
Scoping-review workflow & span\_years, agreement & \stable if the literature spans at least 20 years and agreement $\ge0.65$; else \forming \\
Pooled-estimate workflow & n\_rr, i\_squared & \unstable if $I^2>0.95$; \stable if at least 10 pooled estimates and $I^2<0.75$; else \forming \\
Decline-triggered workflow & decline\_signal, agreement & \unstable on a decline signal (decline\_signal $\ge0.45$); else \stable if agreement $\ge0.7$, else \forming \\
\midrule
\multicolumn{3}{@{}l}{\emph{Single signals, oriented so that a higher score means change is more expected}} \\
$1-|p_t|$ (direction margin) & pooled\_direction & score $=1-|p_t|$ \\
Low agreement & agreement & score $=1-$ agreement \\
Recent dissent & dissent\_recent & score $=$ dissent\_recent \\
Agreement contraction & agreement\_contraction & score $=-$ agreement\_contraction \\
Absent consensus & has\_consensus & score $=1-$ has\_consensus \\
Heterogeneity $I^2$ & i\_squared & score $=I^2$ \\
Imprecision & ci\_crosses\_null & score $=$ share of confidence intervals crossing the null \\
Low evidence volume & n\_pre\_t & score $=-$ n\_pre\_t \\
Recent accrual share & recency\_density & score $=$ share of records in the last five years \\
Low conclusiveness & naive\_tsa & score $=1-$ naive\_tsa \\
Decline signal & decline\_signal & score $=$ decline\_signal \\
\bottomrule
\end{tabular}
\caption{Operational definitions of the rule comparators and single signals. Variable names refer
to the inventory of \cref{tab:feature-inventory}.}
\label{tab:appops}
\end{table*}

\begin{table}[!t]\centering\scriptsize
\setlength{\tabcolsep}{3pt}
\resizebox{\columnwidth}{!}{%
\begin{tabular}{@{}p{0.36\columnwidth}p{0.60\columnwidth}@{}}
\toprule
Learned model (inputs) & Setting and nested tuning grid; mode over folds \\
\midrule
Generalised additive model (7) & one spline term per \compact input, 10 splines; $\lambda\in\{100,10,1,0.1\}$ \\
Random forest (7) & 400 trees; min\_samples\_leaf $\in\{5,3,1\}$ \\
Gradient boosting, XGBoost (7) & learning rate 0.05; max\_depth $\in\{2,3\}\times$ n\_estimators $\in\{100,300\}$ \\
\midrule
Random forest (all) & 400 trees; min\_samples\_leaf $\in\{5,3,1\}$ \\
Elastic net (all) & $C\in\{0.01,0.1,1,10\}\times$ l1\_ratio $\in\{0.2,0.5,0.8\}$ \\
Gradient boosting, XGBoost (all) & learning rate 0.05; max\_depth $\in\{2,3\}\times$ n\_estimators $\in\{100,300\}$ \\
SVM, RBF kernel (all) & $C\in\{0.1,1,10\}\times\gamma\in\{\text{scale},0.01\}$, balanced weights \\
$k$-nearest neighbours (all) & $k\in\{5,15,45\}$, distance weights \\
Neural network (all) & one hidden layer of 32 units, $\alpha\in\{10^{-3},10^{-2},10^{-1}\}$, early stopping \\
Decision tree (all) & max\_depth $\in\{2,3,4\}$ \\
Logistic regression (all) & unpenalised, standardised inputs \\
Balanced logistic (all) & the \compact estimator on all pre-cutoff variables \\
Gaussian naive Bayes (all) & no tuning \\
Best single variable & the single oriented variable with the highest AUROC on the training claims \\
\bottomrule
\end{tabular}}
\caption{Learned baselines, standardised inputs, tuned per outer fold by nested five-fold GroupKFold
over the training claims. ``7'' means the seven \compact inputs, ``all'' the full pre-cutoff
representation. Implementations: scikit-learn \citep{Pedregosa2011}, XGBoost
\citep{ChenGuestrin2016}, pyGAM \citep{pyGAM}.}
\label{tab:appgrids}
\end{table}

\section{Prompted LLM Forecasters on Label Digests}
\label{app:llm}
\paragraph{Summary.} On the 1{,}001-unit tie-free subset, the strongest valid digest prompts reach
AUROC 0.841 (Qwen2.5-72B, zero-shot) and 0.843 (Llama-3.1-70B, one-shot), against 0.938 for
\compact refitted on the same units (paired differences $+0.096$ \ci{0.031}{0.172} and $+0.095$
\ci{0.054}{0.147}). This is not a controlled architecture comparison: the baseline is supervised,
prompting is not; named prompts may draw on pretraining knowledge; and smaller models are
option-order sensitive (see below). The result is limited to the tested digest-prompting
protocol and is not a claim about language models in general. The same models on the calendar
cohort, with and without pre-origin abstracts, are in \cref{app:text}.

Each model sees a digest of the claim's pre-cutoff direction labels: the number of records published
before the cutoff by label, the span of years, the largest directional share over all years and over
the last six years before the cutoff, and year-by-year counts (the earliest eight and the latest
twenty years when the history is longer than 28 years). No post-cutoff record enters the prompt.
\cref{fig:llmprompt} gives the system instructions. The forecast is asked with the claim named
(``does \{exposure\} affect \{outcome\}?'') and blinded (``does an exposure (name withheld) affect a
health outcome (name withheld)?''). The answer is read from the next-token log-probabilities of the
option letters, and the option probabilities are a softmax over the letters; the share of next-token
probability falling on the letters is recorded for every run (0.95 to 1.00 for seven models, 0.68
to 0.90 for DeepSeek-R1-Distill-8B, and below 0.011 for DeepSeek-R1-Distill-70B, whose rows are
therefore invalid and marked $^{\dagger}$). Every task runs twice with the option letters rotated
and every reported value averages the two orders' probabilities. In the one-shot setting one solved
example from a unit of a different claim, drawn at random with a fixed seed, is shown as an earlier
chat turn with its true answer, so the held-out claim stays unseen. The reasoning distills open a
thinking block after the assistant tag, which is closed empty so that the next token is the answer;
Qwen3-8B runs with thinking disabled. The 32B and 70B-class models are pre-quantised 4-bit
checkpoints. Nothing is fitted. The prompt asks whether the predominant direction will reverse, which is a
protective$\leftrightarrow$harmful reversal; it does not distinguish a no-predominant-direction
state, so the stored forecasts are scored on the 1{,}001-unit tie-free subset ($p_t\neq0$ and
$p^{+}_t\neq0$), where \cref{eq:target} reduces to exactly that event, against \compact
refitted leave-one-claim-out on the same units (0.938 \ci{0.907}{0.964}, AUPRC 0.417). No prompt
was re-run. \cref{tab:llmfull} gives the AUROC and AUPRC of every model and setting on that
subset; the paired $\Delta$AUROC of the refit \compact against every valid forecast row excludes
zero (smallest gap $+0.095$ \ci{0.054}{0.147}, Llama-3.1-70B one-shot named), while the paired
$\Delta$AUPRC against the strongest forecast, Qwen2.5-72B zero-shot named, is $+0.142$
\ci{-0.018}{0.287} and covers zero. Scale helps, mostly on the named forecast: Llama-3.1's
zero-shot named AUROC rises from 0.371 at 8B to 0.814 at 70B and Qwen2.5's from 0.496 at 7B to
0.841 at 72B.

\begin{figure}[!t]
\small
\begin{tabular}{@{}p{0.95\linewidth}@{}}
\toprule
\textsc{forecast, system}: You forecast how a body of scientific evidence will develop after a
historical moment. You see only studies published before the cutoff year. When the studies
published after the cutoff are pooled, will the predominant direction of their protective and
harmful findings be the opposite of the predominant direction before the cutoff? Answer with
exactly one letter: A if the predominant direction will reverse; B if the predominant direction will
stay the same. Answer with the letter only. \\[3pt]
\textsc{user}: \{claim, named or blinded\} \textbackslash n Cutoff year: \{t\} \textbackslash
n\textbackslash n \{digest\} \textbackslash n\textbackslash n Will the predominant direction reverse
after \{t\}? \\
\bottomrule
\end{tabular}
\caption{The forecasting prompt, shown with the first option order.}
\label{fig:llmprompt}
\end{figure}

\begin{table*}[!t]\centering\scriptsize
\setlength{\tabcolsep}{4pt}
\resizebox{\textwidth}{!}{%
\begin{tabular}{@{}lllllllll@{}}
\toprule
 & \multicolumn{4}{c}{Zero-shot} & \multicolumn{4}{c}{One-shot} \\
\cmidrule(lr){2-5}\cmidrule(lr){6-9}
 & \multicolumn{2}{c}{Claim named} & \multicolumn{2}{c}{Claim blinded} & \multicolumn{2}{c}{Claim named} & \multicolumn{2}{c}{Claim blinded} \\
\cmidrule(lr){2-3}\cmidrule(lr){4-5}\cmidrule(lr){6-7}\cmidrule(lr){8-9}
Model & AUROC & AUPRC & AUROC & AUPRC & AUROC & AUPRC & AUROC & AUPRC \\
\midrule
Llama-3.1-8B-Instruct & .371 [.284, .452] & .045 & .393 [.325, .459] & .048 & .636 [.564, .706] & .087 & .598 [.525, .671] & .073 \\
Qwen2.5-7B-Instruct & .496 [.423, .583] & .059 & .622 [.518, .719] & .083 & .663 [.596, .726] & .121 & .576 [.513, .638] & .077 \\
Mistral-7B-Instruct-v0.3 & .604 [.462, .724] & .080 & .660 [.559, .748] & .101 & .528 [.451, .605] & .070 & .500 [.426, .574] & .062 \\
Qwen3-8B (thinking off) & .675 [.546, .793] & .121 & .702 [.590, .794] & .121 & .633 [.563, .700] & .079 & .588 [.513, .654] & .070 \\
DeepSeek-R1-Distill-Llama-8B & .388 [.301, .475] & .049 & .309 [.234, .392] & .041 & .461 [.361, .555] & .057 & .401 [.312, .491] & .051 \\
Qwen2.5-32B-Instruct (4-bit) & .790 [.714, .851] & .153 & \textbf{.805} [.735, .860] & .173 & .804 [.758, .847] & .180 & .694 [.621, .760] & .108 \\
Llama-3.1-70B-Instruct (4-bit) & .814 [.756, .858] & .170 & .793 [.731, .845] & .164 & \textbf{.843} [.790, .886] & .184 & \textbf{.829} [.774, .874] & .173 \\
Qwen2.5-72B-Instruct (4-bit) & \textbf{.841} [.764, .911] & .275 & .732 [.605, .827] & .162 & .770 [.705, .831] & .134 & .619 [.561, .677] & .073 \\
DeepSeek-R1-Distill-Llama-70B (4-bit)$^{\dagger}$ & n/a & n/a & n/a & n/a & n/a & n/a & n/a & n/a \\
\midrule
\compact, refit on the tie-free units & \textbf{.938} [.907, .964] & \textbf{.417} & \textbf{.938} [.907, .964] & \textbf{.417} & \textbf{.938} [.907, .964] & \textbf{.417} & \textbf{.938} [.907, .964] & \textbf{.417} \\
\bottomrule
\end{tabular}}
\caption{Direct reversal forecasts on the 1{,}001-unit tie-free subset (60 events, prevalence
0.060), where the target is a protective$\leftrightarrow$harmful reversal: AUROC [claim-clustered
95\% CI] and AUPRC, every value averaged over two option orders. \compact is refitted
leave-one-claim-out on the same units and never sees the claim name, so its value is the same in
every column. $^{\dagger}$Invalid read-out (0.02 to 1.0\% of the next-token probability on an
option letter); its values are unavailable rather than meaningful scores. Conclusions are bounded
to this digest-prompting setup.}
\label{tab:llmfull}
\end{table*}

\paragraph{Option-order sensitivity.} The 7B to 8B models are dominated by letter-position
preference: their $P(\text{reverse})$ under the two option orders correlates at $r=-0.27$ to $0.05$
zero-shot (Qwen3-8B scores 0.21 under one order and 0.78 under the other). The 32B and 70B-class
models are steadier: Qwen2.5-72B's forecasts correlate at $r=0.89$ (named) and $0.85$ (blinded)
across orders, and Qwen2.5-32B's at $r=0.68$. Llama-3.1-70B's zero-shot forecasts anticorrelate
across orders ($r=-0.45$, $-0.61$). Averaging the two orders removes the position bias but not the
instability.

\section{Text-Reading Forecasters on the Calendar Cohort}
\label{app:text}
Both families use only records published before the origin and are evaluated in the calendar cells
of \cref{tab:protocol}; trained models use claim-disjoint units whose outcomes are observable at the
origin.

\paragraph{Abstract embeddings.} Each record's title and abstract (the title alone when no abstract
is available) is embedded once with a biomedical sentence encoder (S-PubMedBert-MS-MARCO,
768 dimensions, L2-normalised). A unit's text features are the mean embedding of all its pre-origin
records, the mean of its 20 most recent pre-origin records and the difference between the two,
reduced to 64 dimensions by principal components fitted on the training units only. A class-balanced
$L_2$ logistic regression is fitted per origin on text alone, on text with the two composition inputs,
and on the composition inputs alone (\cref{tab:embed}). Text alone ranks later change at 0.606-0.646, and adding it to the composition
inputs lowers AUROC by an established margin at every horizon.

\begin{table}[!t]\centering\footnotesize
\setlength{\tabcolsep}{3pt}
\resizebox{\columnwidth}{!}{%
\begin{tabular}{@{}lccccl@{}}
\toprule
Horizon & Margin & Text & Text + comp. & Comp. & Text added [95\% CI] \\
\midrule
3 y & .907 & .606 & .746 & .913 & $-$.166 [$-$.253, $-$.078] \\
5 y & .897 & .646 & .781 & .920 & $-$.139 [$-$.265, $-$.040] \\
10 y & .867 & .608 & .758 & .876 & $-$.118 [$-$.194, $-$.045] \\
\bottomrule
\end{tabular}}
\caption{Embedding forecasters on the calendar cohort (AUROC). ``Comp.'' is the signed two-variable
score; ``Text added'' is text + composition minus composition, paired over claims.}
\label{tab:embed}
\end{table}

\paragraph{Language models.} Each prompt states the exposure, the outcome, the origin year and the
pre-origin counts of \textsc{harmful}, \textsc{protective} and \textsc{null} records with the current
majority direction (the digest). The abstract condition adds the titles and abstracts of the eight most
recent pre-origin records, each truncated to its first 1{,}000 characters, so that every prompt fits the
context window without selection by content. The masked condition replaces the exposure and outcome
names with placeholders in the prompt and in the abstracts. The model answers a two-option question
(the predominant direction reverses or stays) under both option orders, and the score is the
probability of the reversal option averaged over the two orders. At 7-8B (Llama-3.1-8B and Qwen2.5-7B) every cell lies
between 0.478 and 0.579 and is established below the margin rule in every cell; the per-cell values are released
with the benchmark. At 70B and 72B (\cref{tab:llmtext}) no configuration establishes superiority over the margin rule; the only positive point difference, Qwen at ten years (0.875 against 0.867), is $+0.008$ with an interval that
spans zero, and the remaining differences are negative point estimates whose intervals are released
with the per-unit scores. Abstracts add $+0.013$ to $+0.023$ when the claim is named and $+0.085$ to $+0.104$ when it
is masked. Naming the claim is worth $+0.119$ to $+0.129$ on the digest alone but $+0.037$, $+0.019$
and $-0.004$ once abstracts are supplied (established only at three years), and its value is not
concentrated at early origins, as recall of later literature would predict. Because post-origin pretraining
can give these models information unavailable at the historical origin, these runs are not clean
prospective tests; nevertheless, none establishes a gain over the margin rule, and this bounds the conclusion to the tested digest and eight-abstract prompts, not to text-reading methods in general.

\begin{table}[!t]\centering\footnotesize
\setlength{\tabcolsep}{4pt}
\begin{tabular}{@{}llccc@{}}
\toprule
Model & Input & 3 y & 5 y & 10 y \\
\midrule
Llama 70B & digest, named & .811 & .837 & .819 \\
Llama 70B & abstracts, named & .825 & .850 & .842 \\
Llama 70B & abstracts, masked & .788 & .831 & .846 \\
Qwen 72B & digest, named & .831 & .864 & .875 \\
\midrule
Margin rule & -- & .907 & .897 & .867 \\
\bottomrule
\end{tabular}
\caption{Prompted language models on the calendar cohort (AUROC).}
\label{tab:llmtext}
\end{table}

\section{Reproducibility}
\label{app:repro}
The artefacts contain the claim list with its fixed query for every claim, the retrieved PubMed
identifiers with publication years and the retrieval audit, the extractor-assigned direction labels
with their confidences (Llama-3.1-8B-Instruct, pinned revision, bfloat16, closed four-label
read-out), the claim-cutoff benchmark with its pre-cutoff columns and post-cutoff fields under the
tie convention of \cref{eq:target}, the 145-claim subset benchmark, the two external benchmarks as claim-cutoff
units, the stored option probabilities of every language-model run with the model revision,
settings and the complete prompts of \cref{fig:prompt,fig:llmprompt}, every result
table as CSV and JSON with a metadata block recording the sample, split, bootstrap (claim-clustered
percentile, 2{,}000 resamples, seed 3) and decision rules, and the code: one library holding the
feature builder, the leave-one-claim-out fitting, the metrics and the bootstrap;
one script per experiment; and one launcher that regenerates every table. Randomness enters only through the seeded bootstrap resamplers, the seeded
tree and network learners, and the seeded random draws of the cutoff designs and one-shot
examples. Every numbered analysis checks its own consistency with the system at run time: the
benchmark rebuilds byte-identically from the records and labels, the PMID-disjoint rebuild
reproduces every unit when nothing is removed, and every additional experiment asserts that its
\compact scores equal those of the system. The audit analyses (numeric anchor against both
annotations, by era and by anchor scale with the signed-anchor confusion matrices, agreement audit
with the branch-free refits, the common-support intervention, its branch slices and the branch
membership mask, the exact pair decomposition with per-pair contributions, cumulative and
three-state endpoints, orientation audit, two-variable transfer, calendar cohort, stationary nulls
with the partition grid, invariance audit and Null~D draws, wild cluster, clustered-ROC and
refitting bootstraps, fold-level coefficients and slopes, class-by-bin and documented-reversal
tabulations, and second annotation) are one script each over a
frozen snapshot of the released benchmark, records and code, with a manifest recording input
hashes, seeds, library versions and output files; each asserts at run time that its rebuild of the
shipped columns is exact. Abstract text is not redistributed.
\paragraph{Independent rebuild and reproducibility levels.} The release contains per-unit predictions and outcomes for every reported metric and interval, association-level labels from both pipelines (primary first-token scores and second-pipeline label probabilities) with identifiers and years, intersection labels, the calendar-cohort census, the external unit files, the adjudication frame and the protocol code. An independent rebuild from the released labels reproduces every unit and event count (1{,}012 units with 71, 52 and 55 events; calendar cohort 756/58, 625/45 and 492/35), matches the frozen feature builder on all 1{,}012 units and reproduces the reported AUROCs (0.948, 0.857, 0.898; 0.913, 0.920, 0.876). Three reproducibility claims are distinct and the release supports them to different degrees: reproduction of stored metrics (fully, from the per-unit files); reconstruction from released labels (counts, pooled directions, agreement under every definition, every endpoint, the stationary probabilities and the null controls); and re-execution of text-dependent methods, which requires the original abstracts, is possible for the source and calendar cells by retrieving the released PMIDs, and is not possible for the external corpora, which ship as unit-level label counts without identifiers. Two evaluation tracks follow from this last distinction: label-feature methods are eligible for every check of \cref{sec:protocol}, text-reading methods for checks 1-4 only, since the external corpora ship without document inputs; text-reading methods can reach Tier~A and a source-only Tier~C marked as such, never Tier~B, and an unavailable check is reported as unavailable, not as passed. The artefacts are supplied through the anonymous repository named in the main text
(\cref{app:snapshot} lists the rebuild levels it supports); the exact retrieval timestamp was not recorded but is bounded as described, and the candidate-list provenance is not reconstructible. The feature inventory that follows is a data dictionary for
the released columns.

\onecolumn
\begingroup
\footnotesize
\setlength{\tabcolsep}{4pt}
\renewcommand{\arraystretch}{1.12}
\begin{longtable}{@{}p{0.215\linewidth}p{0.105\linewidth}p{0.63\linewidth}@{}}
\caption{Inventory of the 55 pre-cutoff variables, with the block used in the ablation and the
operational definition of each column. $R$ denotes the pre-cutoff records (with a parsed year and an
extractor-assigned direction), $t$ the cutoff year, and $\mathrm{eff}$ the set of effect estimates
recovered from 95\% confidence intervals in the title and abstract. $\dagger$ and $\ddagger$ mark inputs of a descriptive
state rule that is released with the data but not used in this paper; $^{\ast}$ marks
one of the seven \compact inputs. All sentinels are 0.0 unless stated.}
\label{tab:feature-inventory}\\
\toprule
\textbf{Variable} & \textbf{Block} & \textbf{Definition} \\
\midrule
\endfirsthead
\multicolumn{3}{@{}l}{\small\itshape \cref{tab:feature-inventory} continued from previous page}\\
\toprule
\textbf{Variable} & \textbf{Block} & \textbf{Definition} \\
\midrule
\endhead
\midrule
\multicolumn{3}{r@{}}{\small\itshape continued on next page}\\
\endfoot
\bottomrule
\endlastfoot
\texttt{pooled\_direction}\,$^{\ast}$ & direction & Operational pooled direction $p_t$: the raw mean $(h-l)/(h+l)$ over signed pre-cutoff annotations (\textsc{protective}$=-1$, \textsc{harmful}$=+1$; \textsc{null} excluded, \textsc{unclear} omitted) when at least eight signed records are present; otherwise 0 (the gate takes precedence). The ungated raw mean is $r_t$ for a non-empty signed sample. \\
\texttt{agreement}\,$\dagger$\,$^{\ast}$ & direction & Consensus share over all resolved directions: if the null share is $\ge0.50$ it is that null share, else $\max(p,k-p)/k$ over the $k$ signed votes ($p$ = harmful); this is $a$. \\
\texttt{agreement\_recent}\,$\dagger$ & direction & The same share restricted to records with year $\ge t-6$; falls back to \texttt{agreement} when fewer than 4 signed recent votes exist; this is $a_r$. \\
\texttt{dissent\_recent}\,$\dagger$ & direction & Fraction of signed votes with year $\ge t-6$ disagreeing with the reference sign (the consensus verdict, or the overall majority sign when the verdict is 0); 0.0 under the same $<4$ fallback; this is $\delta_r$. \\
\texttt{agreement\_contraction}\,$\dagger$ & direction & Majority share of the later half of the signed-vote sequence minus that of the earlier half, split at $\lfloor k/2\rfloor$ in record order; this is $\gamma$. \\
\texttt{direction\_trajectory} & direction & OLS slope of the signed direction on publication year; 0.0 if fewer than 4 points or zero year range. \\
\texttt{null\_frac} & direction & Share of resolved directions that are \textsc{null}. \\
\texttt{n\_directional}\,$\ddagger$\,$^{\ast}$ & direction & Count of pre-cutoff records with a resolved (non-\textsc{unclear}) direction. \\
\texttt{n\_nonnull}\,$\ddagger$ & direction & Count of pre-cutoff records with a signed (non-\textsc{null}) direction; the direction block is zeroed when this is $<8$. \\
\texttt{consensus\_verdict} & direction & Sign of the consensus $\in\{-1,0,+1\}$; 0 when the null share is $\ge0.50$ or agreement $<0.60$. \\
\texttt{has\_consensus} & direction & $\mathbf{1}[\texttt{agreement}\ge0.60 \wedge \texttt{consensus\_verdict}\neq0]$. \\
\texttt{tau2\_DL} & heterogeneity & DerSimonian-Laird $\hat\tau^2=\max\!\big(0,(Q-(k-1))/C\big)$, $C=\sum w-\sum w^2/\sum w$, $w=1/v$, on the log-ratio or difference scale; 0.0 with fewer than 3 estimates. \\
\texttt{tau2\_REML} & heterogeneity & Iterative REML $\hat\tau^2$ (40 iterations, tolerance $10^{-8}$), initialised at \texttt{tau2\_DL}. \\
\texttt{tau2\_Hedges} & heterogeneity & $\max\!\big(0,\ \mathrm{Var}(y;\,\mathrm{ddof}{=}1)-\overline{v}\big)$, the method-of-moments estimator. \\
\texttt{mean\_tau2} & heterogeneity & Arithmetic mean of the three $\hat\tau^2$ estimates. \\
\texttt{i\_squared} & heterogeneity & $\max\!\big(0,(Q-(k-1))/Q\big)$ from the same $Q$; 0.0 when $Q\le0$ or fewer than 3 estimates. \\
\texttt{spread} & heterogeneity & 90th minus 10th percentile of the effect estimates $y$. \\
\texttt{decline\_signal} & decline & $(\mathrm{early}-\mathrm{late})/\mathrm{early}$, where early and late are the mean $|y|$ over the first and last $k=\max(2,\lfloor n/3\rfloor)$ estimates ordered by year; needs at least 6 estimates. \\
\texttt{early\_magnitude} & decline & Mean $|y|$ over the earliest $k$ effect estimates by year. \\
\texttt{late\_magnitude} & decline & Mean $|y|$ over the latest $k$ effect estimates by year. \\
\texttt{shrink\_ratio} & decline & $\mathrm{late}/\mathrm{early}$. \\
\texttt{synth\_frac} & design & Fraction of pre-cutoff texts matching the evidence-synthesis pattern (\emph{meta-analys*}, \emph{systematic review}, \emph{pooled analys*}). \\
\texttt{rct\_frac} & design & Fraction matching the trial pattern (\emph{randomi[sz]ed}, \emph{double-blind}, \emph{placebo-controlled}, \emph{clinical trial}). \\
\texttt{cohort\_frac} & design & Fraction matching the observational pattern (\emph{cohort}, \emph{prospective}, \emph{longitudinal}, \emph{follow-up}). \\
\texttt{synth\_trend} & design & OLS slope of the synthesis indicator on publication year. \\
\texttt{rct\_trend} & design & OLS slope of the trial indicator on publication year. \\
\texttt{synth\_recent} & design & Mean synthesis indicator over records with year $\ge t-5$; 0.0 when that window is empty. \\
\texttt{design\_upgrade\_frac} & design & Mean synthesis indicator on year $\ge t-5$ minus the mean on year $<t-5$, each term 0 when its window is empty. \\
\texttt{confounding\_trend} & design & OLS slope on year of the adjustment indicator (\emph{adjust*}, \emph{multivariable}, \emph{confound*}, \emph{propensity}, \emph{covariate}). \\
\texttt{n\_pre\_t}\,$^{\ast}$ & volume & Number of pre-cutoff records $|R|$; at least 40 by construction of the cutoff grid. \\
\texttt{n\_records} & volume & $|R|$ again; numerically identical to \texttt{n\_pre\_t} in every row. \\
\texttt{vol\_last5} & volume & Count of pre-cutoff records with year $\ge t-5$. \\
\texttt{vol\_growth} & volume & Add-one smoothed ratio $\big(|\{y\ge t{-}5\}|+1\big)/\big(|\{t{-}10\le y<t{-}5\}|+1\big)$. \\
\texttt{volume\_ratio} & volume & $(|R|-h)/\max(h,1)$ with $h=\lfloor|R|/2\rfloor$: later-half to earlier-half record count. \\
\texttt{recency\_density} & volume & Fraction of pre-cutoff records with year $\ge t-5$. \\
\texttt{accrual\_rate\_5yr} & volume & Records per year over the final five pre-cutoff years, $|\{y\ge t-5\}|/5$. \\
\texttt{span\_years} & derived & $\max(\text{year})-\min(\text{year})$ over the pre-cutoff records. \\
\texttt{n\_rr}\,$^{\ast}$ & derived & Number of usable effect estimates $|\mathrm{eff}|$ (confidence-interval pairs with $se>10^{-6}$ and $|y|<10$). \\
\texttt{n\_rr\_studies} & derived & $|\mathrm{eff}|$ again; identical to \texttt{n\_rr}. \\
\texttt{decline\_sufficient} & derived & Gate flag for the decline block; constant 0.0 in the modelled benchmark. \\
\texttt{naive\_tsa} & derived & $\min\!\big(1,\ \sum_{\mathrm{eff}} se^{-2}/400\big)$: a naive trial-sequential information-fraction proxy; 0.0 when no estimate exists. \\
\texttt{n\_large\_recent} & derived & Count of records with year $\ge t-5$ whose text reports a large enrolment (a five-or-more-digit or comma-grouped count followed by \emph{participants}, \emph{subjects}, \emph{patients}, \emph{individuals}, \emph{women} or \emph{men}). \\
\texttt{n\_ci} & derived & Number of 95\% confidence-interval matches across all pre-cutoff texts, counted before the validity filter, so $\texttt{n\_ci}\ge\texttt{n\_rr}$. \\
\texttt{ci\_rate} & derived & $\texttt{n\_ci}/\max(|R|,1)$: reported intervals per pre-cutoff record. \\
\texttt{ci\_width\_mean} & derived & Mean of $2\times1.96\,se$ over the effect estimates. \\
\texttt{ci\_width\_sd} & derived & Population standard deviation of the same interval widths. \\
\texttt{ci\_crosses\_null} & derived & Fraction of effect estimates with $|y|<1.96\,se$, i.e.\ intervals covering the null. \\
\texttt{val\_mean}\,$^{\ast}$ & derived & Mean of the effect estimates $y$ on their working scale. \\
\texttt{val\_sd}\,$^{\ast}$ & derived & Population standard deviation of $y$. \\
\texttt{val\_spread} & derived & A copy of \texttt{spread}; identical in every row. \\
\texttt{val\_cv} & derived & $\texttt{val\_sd}/\max(|\texttt{val\_mean}|,10^{-6})$. \\
\texttt{n\_numeric} & derived & $|\mathrm{eff}|$ once more; identical to \texttt{n\_rr}. \\
\texttt{numeric\_rate} & derived & Fraction of pre-cutoff records yielding at least one usable effect estimate. \\
\texttt{numeric\_rate\_recent} & derived & The same fraction restricted to records with year $\ge t-5$; 0.0 when that window is empty. \\
\texttt{numeric\_trend} & derived & OLS slope on year of the per-record indicator ``this abstract yields an effect estimate''. \\
\end{longtable}
\endgroup
\twocolumn
% \input{figures/appendix_final}

\end{document}